% concatenated from mcrae_yang_conjecture_7_16.tex
\documentclass[12pt,leqno]{amsart}
\usepackage{pifont}
\usepackage{stmaryrd}
%\usepackage{amsfonts}
%\usepackage{amscd}
\usepackage{dsfont}
\usepackage{color}
\usepackage{mathrsfs}
\usepackage{amsmath}
\usepackage{amssymb}
\usepackage[hidelinks]{hyperref}
\usepackage{bookmark}
\usepackage[bottom]{footmisc}
\usepackage{verbatim}
\usepackage{extarrows}
\usepackage{mathtools}
\usepackage{orcidlink}

\oddsidemargin -.1in \evensidemargin -.1in \textwidth 6.5in
\textheight 8.2in
\linespread{1.3}

\numberwithin{equation}{section}
\newtheorem{thm}{Theorem}[section]
\newtheorem{lem}[thm]{Lemma}
\newtheorem{ex}[thm]{Example}
\newtheorem{prop}[thm]{Proposition}
\newtheorem{cor}[thm]{Corollary}
\newtheorem{defi}[thm]{Definition}
\newtheorem{rmk}[thm]{Remark}
\newtheorem{notation}[thm]{Notation}

\newcommand{\End}{\operatorname{End}}
\newcommand{\Hom}{\operatorname{Hom}}
\newcommand{\Ext}{\operatorname{Ext}}
\newcommand{\Soc}{\operatorname{Soc}}
\newcommand{\rad}{\operatorname{rad}}
\newcommand{\ch}{\operatorname{ch}}
\newcommand{\one}{\mathbf{1}}
\newcommand{\id}{\operatorname{id}}
\newcommand{\wt}{\operatorname{wt}}
\newcommand{\xqz}[1]{\lfloor #1 \rfloor}
\newcommand{\Res}{\operatorname{Res}}
\newcommand{\Aut}{\operatorname{Aut}}
\newcommand{\Vir}{\mathfrak{Vir}}
\newcommand{\TL}{\mathcal{TL}}
\newcommand{\VirMod}{\mathcal L}
\newcommand{\boxt}{\boxtimes}
\newcommand{\KL}{KL}
\newcommand{\Oc}{\mathcal O}
\newcommand{\Kac}{\mathcal K}
\renewcommand{\mod}{\operatorname{mod}}

\def\Z{\mathbb{Z}}
\def\N{\mathbb{N}}
\def\A{\mathbb{A}}
\def\C{\mathbb{C}}
\def\R{\mathbb{R}}
\def\Q{\mathbb{Q}}
\def\F{\mathbb{F}}

\allowdisplaybreaks

\begin{document}


\title[Exactness and DS realization]{Exactness and Drinfeld--Sokolov Realization of the McRae--Yang Tensor Functors}

\author{Shun Xu~\orcidlink{0009-0006-8080-8107}}

\address{School of Mathematical Sciences, Anhui University, Anhui, Hefei, 230601, China}

\email{shunxu@ahu.edu.cn}

\subjclass[2020]{Primary 17B69, 17B68; Secondary 18M15, 81R10}

\keywords{Kazhdan--Lusztig category, quantum Drinfeld--Sokolov reduction, Virasoro algebra, logarithmic tensor category, projective object, Kac module}

\begin{abstract}
Let $p,q\geq2$ be coprime, set $k=-2+p/q$, and let
$c_{p,q}=1-6(p-q)^2/(pq)$.  McRae and Yang constructed right exact
braided tensor functors from the non-semisimple Kazhdan--Lusztig
categories of $V^{-2+p/q}(\mathfrak{sl}_2)$ and
$V^{-2+q/p}(\mathfrak{sl}_2)$ to the Virasoro category
$\Oc_{c_{p,q}}$, and conjectured that these functors are exact and agree
with the corresponding quantum Drinfeld--Sokolov reductions.  We prove
this conjecture.  For the $(p,q)$ branch, we determine the images of all
indecomposable projective objects without assuming exactness, using
projective tensor-product recursions and exact generalized
conformal-residue projections.  Projective faithfulness is then combined
with a one-row Virasoro extension analysis to identify the non-wall images
and to prove full faithfulness on projectives.  The Virasoro input is
matched objectwise with Nakano's logarithmic extension theorem away from
the vacuum edge; the exceptional vacuum edge is handled directly by the
staggered-module theory of Kyt\"ol\"a--Ridout, including the higher prime
singular-vector branch.  Independently, principal Drinfeld--Sokolov
reduction is shown to be exact and faithful on the whole finite-length
affine category and to have the same Weyl, simple, and projective images.
Compatibility with the non-standard affine twist identifies the canonical
nilpotent endomorphisms on projectives and removes the remaining scalar
ambiguity in the comparison of adjacent projective morphisms.  Projective
density for right exact functors then yields a natural isomorphism
\[
 F_{p,q}\cong
 H^0_{DS,+}\big|_{\KL^k(\mathfrak{sl}_2)}.
\]
The same argument after interchanging $p$ and $q$ identifies the second
McRae--Yang functor with the transposed Drinfeld--Sokolov reduction.
\end{abstract}

\maketitle

\section{Introduction}

Quantum Drinfeld--Sokolov reduction relates representations of affine Lie
algebras to representations of $W$-algebras by a BRST construction.  At
admissible non-integral levels, both sides are non-semisimple, and the
resulting functor interacts with extensions, projective objects, and tensor
products in ways that are not visible from the semisimple theory.  For
$\widehat{\mathfrak{sl}}_2$, McRae and Yang obtained a tensor-categorical
counterpart of this reduction and conjectured that the two constructions
coincide.  The aim of this paper is to prove their conjecture.

Let $p,q\geq2$ be relatively prime and put
\begin{equation}\label{eq:parameters}
 k=-2+\frac pq,
 \qquad
 c_{p,q}=1-\frac{6(p-q)^2}{pq}.
\end{equation}
We write $\KL^k(\mathfrak{sl}_2)$ for the finite-length,
grading-restricted generalized module category of the universal affine
vertex algebra $V^k(\mathfrak{sl}_2)$.  McRae and Yang determined its
projective objects and identified the projective subcategory with the
root-of-unity tilting category for quantum $\mathfrak{sl}_2$
\cite{McRaeYang}.  Their universal property produces a right exact braided
tensor functor
\begin{equation}\label{eq:F-intro}
 F_{p,q}:\KL^k(\mathfrak{sl}_2)\longrightarrow\Oc_{c_{p,q}},
 \qquad F_{p,q}(V_2)\cong\Kac_{2,1},
\end{equation}
where $\Oc_{c_{p,q}}$ is the finite-length Virasoro tensor category of
\cite{McRaeSopin}.  They conjectured that $F_{p,q}$ is exact and that,
after forgetting the tensor structure, it is naturally isomorphic to the
quantized Drinfeld--Sokolov reduction; the parameter-interchanged functor
should satisfy the analogous statement
\cite[Conjecture~7.16]{McRaeYang}.

Our main theorem proves both assertions.

\begin{thm}\label{thm:main}
Let $p,q\geq2$ be relatively prime and let $k$ and $c_{p,q}$ be as in
\eqref{eq:parameters}.  There is a natural isomorphism of $\C$-linear
functors
\[
 F_{p,q}\cong H^0_{DS,+}\big|_{\KL^k(\mathfrak{sl}_2)}.
\]
Consequently $F_{p,q}$ is exact and faithful, and its restriction to the
projective subcategory is fully faithful.  The transposed functor of
McRae--Yang satisfies the analogous statement; see
Corollary~\ref{cor:transposed-reduction}.  In particular,
Conjecture~7.16 of McRae--Yang holds.
\end{thm}

The first obstacle is that the known structure of affine projectives cannot
simply be pushed through $F_{p,q}$.  For $n\geq1$ and
$1\leq r\leq p-1$, the projective cover $P_{np+r}$ has a nonsplit Weyl
filtration
\[
 0\longrightarrow V_{np-r}\longrightarrow P_{np+r}
 \longrightarrow V_{np+r}\longrightarrow0,
\]
but right exactness of $F_{p,q}$ gives no information about the first
arrow.  We therefore reconstruct the image of every indecomposable
projective intrinsically in the Virasoro category.  The reconstruction uses
the projective tensor-product recursions for $V_2\boxtimes P_j$, the exact
functor $K_2\boxtimes-$ with $K_2=\Kac_{2,1}$, and exact projection to
generalized conformal-residue sectors modulo $\Z$.  In particular, no
monomorphism in the reconstructed Kac filtrations is obtained by applying
$F_{p,q}$ to a nonsplit source sequence.

A second ingredient is faithfulness on the projective subcategory.  The
source projectives form reflection chains whose adjacent morphisms and
one-dimensional radical loops are explicitly known from the
root-of-unity tilting model.  For $p\geq3$, a Temperley--Lieb tensor-ideal
argument detects the first critical Jones--Wenzl projector and prevents a
nonzero projective morphism from being annihilated; the case $p=2$ is
handled directly from the reconstructed logarithmic objects and the
non-standard affine twist.  This yields projective faithfulness before any
target Hom-space computation is used.

The Virasoro extension theory needed to identify the reconstructed objects
is isolated in a one-row form.  We write $K_j=\Kac_{j,1}$ and let $S_j$
denote its simple quotient.  For
\[
 a=np-r,\qquad b=np+r,\qquad c=(n+2)p-r,
\]
the one-row Kac modules $K_a$ and $K_b$ are matched with consecutive terms
in Nakano's Felder dictionary.  When $a\geq2$, every logarithmic extension
\[
 0\longrightarrow K_a\longrightarrow E\longrightarrow K_b
 \longrightarrow0
\]
that occurs in our category is shown, object by object, to satisfy all
ambient hypotheses of Nakano's logarithmic extension theorem.  Nakano's
theorem is stated for an arbitrary logarithmic representative of the
relevant Ext class and gives the required distinguished-quotient
nonsplitting.  The sole exception to the support argument is the vacuum
edge $a=1$.  There the right Verma module is of braid type; the smaller
prime branch is treated by Kyt\"ol\"a--Ridout's Proposition~7.5, while
for the higher prime branch we use their determinant refinement immediately
following that proposition.  This supplies the one-row recognition theorem
for all ordered coprime pairs $(p,q)$ and leads to the identification
\[
 F_{p,q}(P_{np+r})\cong P(\tau_{n,r})
\]
for every non-wall projective.  The resulting short exact sequences give
the target Hom-space upper bounds, and projective faithfulness supplies the
matching lower bounds.  Hence
\[
 F_{p,q}\big|_{\mathcal P^k}
 \quad\text{is fully faithful}.
\]

Drinfeld--Sokolov reduction is established independently of this
projective reconstruction.  We compare the standard rank-one BRST
conventions, prove higher cohomology vanishing by finite-length
d\'evissage, and obtain an exact faithful functor
\[
 H:=H^0_{DS,+}:\KL^k(\mathfrak{sl}_2)\longrightarrow\Oc_{c_{p,q}}.
\]
It satisfies
\[
 H(V_j)\cong K_j,\qquad H(L_j)\cong S_j,
\]
and the same one-row recognition theorem identifies $H(P_j)$ with the
projective image already found for $F_{p,q}$.  The argument proving
exactness of $H$ is logically independent of $F_{p,q}$ and of the
projective comparison.

Objectwise agreement on projectives is not yet enough for the main
result.  The remaining issue is to make these isomorphisms natural along
the reflection chains.  The non-standard affine twist provides the needed
normalization.  On a non-wall affine projective, the nilpotent part of the
conformal Hamiltonian is the canonical generator of the radical of the
endomorphism algebra.  The BRST exponential identity carries this
canonical element to the corresponding Virasoro nilpotent endomorphism.
Consequently any objectwise comparison intertwines the radical loop.  A
single scalar normalization on one arrow of each adjacent pair then forces
the opposite arrow as well.  This produces a natural isomorphism on all
projectives.  A general projective-density statement for additive right
exact functors extends it uniquely to the whole abelian category, proving
Theorem~\ref{thm:main}.

The theorem also determines the tensor-categorical status of the
McRae--Yang functor.  Exactness and faithfulness follow from the comparison
with $H$, while full faithfulness on projectives has already been obtained
internally.  Transporting the braided tensor structure of $F_{p,q}$ across
the natural isomorphism gives the restricted Drinfeld--Sokolov functor a
braided tensor structure.  This last observation is a transport statement;
no independent monoidal construction on the BRST complex is assumed.
After interchanging $p$ and $q$, the same proof applies.  A rigid
normalization argument identifies the parameter-swapped generating Kac
object, together with its evaluation and coevaluation maps, with the second
normalization in \cite[Theorem~7.15]{McRaeYang}; this yields the second
functor appearing in Conjecture~7.16.

The paper is organized as follows.  Section~\ref{sec:prelim} develops the
source projective data, conformal-residue calculus, the one-row Felder
dictionary, and the Virasoro extension input.  Section~\ref{sec:DS}
proves the independent Drinfeld--Sokolov package: exactness, faithfulness,
Weyl and simple images, and exponential compatibility.  In
Section~\ref{sec:images} we reconstruct all projective images under
$F_{p,q}$ and establish projective faithfulness without assuming
exactness.  Section~\ref{sec:fullness} identifies the logarithmic
non-wall images and computes the target Hom spaces, proving full
faithfulness on projectives.  Section~\ref{sec:comparison} compares
$F_{p,q}$ with Drinfeld--Sokolov reduction on projectives, proves
naturality by locking the radical scalars, and extends the comparison to
the whole category.

\section{Structural input and the one-row Felder dictionary}\label{sec:prelim}

\subsection{Affine projectives, Hom spaces, and the minus twist}

Write $V_r$ for the Weyl module induced from the $r$-dimensional simple
$\mathfrak{sl}_2$-module, $L_r$ for its simple quotient, and
$P_r\twoheadrightarrow L_r$ for the projective cover.  We collect here the
part of the projective structure of $\KL^k(\mathfrak{sl}_2)$ that will be
used later.  By \cite[Theorem~2.2]{McRaeYang},
\[
 V_{np}=L_{np}\qquad(n\ge1).
\]
Theorem~4.8 of \cite{McRaeYang} identifies the projective covers and gives
the projective tensor-product recursions used below, while the nonsplit
Weyl filtration and Loewy structure of the non-wall projectives are those
of \cite[Theorem~4.4]{McRaeYang}.  The projective Hom spaces used below
are the explicit Hom spaces displayed in the proof of
\cite[Theorem~6.6]{McRaeYang}.  Thus the source-side projective data are
read directly from the affine projective category; the tilting equivalence
will enter only in the Temperley--Lieb proof of projective faithfulness.

We call a positive label $j$ a \emph{wall label} if $p\mid j$ and a
\emph{non-wall label} otherwise.  The category has enough projectives.
For $1\leq r<p$ one has $P_r=V_r$, while for every wall label $j$,
\[
 P_j=V_j=L_j;
\]
in particular $P_{np}=V_{np}=L_{np}$ is simple-projective for $n\geq1$.
All of these projective covers are indecomposable; the wall simplicity
is the preceding application of \cite[Theorem~2.2]{McRaeYang}, while the
projective-cover assertions are supplied by \cite[Theorem~4.8]{McRaeYang}.
If $n\geq1$ and $1\leq r\leq p-1$, then
\begin{equation}\label{eq:source-projective-filtration}
 0\longrightarrow V_{np-r}\longrightarrow P_{np+r}
 \longrightarrow V_{np+r}\longrightarrow0
\end{equation}
is non-split.  The monomorphism and epimorphism in
\eqref{eq:source-projective-filtration} will only ever be used up to
nonzero scalar; no normalization of either map is imposed.  Moreover,
\begin{equation}\label{eq:source-composition-factors}
 [P_{np+r}]
 =2[L_{np+r}]+[L_{np-r}]+[L_{(n+2)p-r}].
\end{equation}
The tensor products $V_2\boxtimes P_j$ satisfy the projective recursions of
\cite[Theorem~4.8]{McRaeYang}; for $p=2$ we use the two exceptional
recursions as collected in \cite[Theorem~1.4(1)]{McRaeYang}.  These recursions will be the source-side input in
Section~\ref{sec:images}.

For $1\leq r\leq p-1$ define the reflection chain
\begin{equation}\label{eq:reflection-chain}
 s_{2m}(r)=2mp+r,\qquad
 s_{2m+1}(r)=2(m+1)p-r
 \qquad(m\geq0).
\end{equation}
When the initial label is fixed we suppress it and write $s_i=s_i(r)$.

\begin{lem}[Reflection-chain partition]\label{lem:reflection-partition}
The non-wall positive labels are partitioned by the reflection chains:
\[
 \{j\geq1:p\nmid j\}
 =\bigsqcup_{r=1}^{p-1}\{s_i(r):i\geq0\}.
\]
More explicitly, if $j=np+r$ with $n\geq0$ and $1\leq r\leq p-1$,
then $j=s_n(r)$ when $n$ is even and $j=s_n(p-r)$ when $n$ is odd.
\end{lem}

\begin{proof}
If $n=2m$, then $j=2mp+r=s_{2m}(r)$.  If $n=2m+1$, then
$j=(2m+1)p+r=2(m+1)p-(p-r)=s_{2m+1}(p-r)$.  Conversely each label in
\eqref{eq:reflection-chain} is positive and not divisible by $p$.
The Euclidean decomposition $j=np+r$ with $n\ge0$ and $1\le r<p$ is
unique, and its parity determines uniquely which of $r$ or $p-r$ is the
initial label of the reflection chain.  It therefore determines the chain
and the chain index uniquely; in particular, two distinct reflection chains
cannot meet.
\end{proof}

For later comparison of Hom spaces we record the projective morphisms along
these chains.

\begin{prop}[Source projective Hom package]\label{prop:source-Hom-package}
For a fixed reflection chain \eqref{eq:reflection-chain},
\[
 \dim\Hom(P_{s_i},P_{s_j})=
 \begin{cases}
 2,&i=j\geq1,\\
 1,&i=j=0,\\
 1,&|i-j|=1,\\
 0,&|i-j|\geq2.
 \end{cases}
\]
If $P_{s_i(r)}$ and $P_{s_j(r')}$ belong to distinct reflection chains,
then both Hom spaces between them vanish.  The wall projectives $P_{np}$ are isolated
simple-projectives:
\[
 \End(P_{np})=\C,\qquad
 \Hom(P_{np},P_j)=\Hom(P_j,P_{np})=0\quad(j\ne np).
\]
In particular, $\Hom(P_{np},P_{mp})=0$ for $n\ne m$.

Moreover, for each adjacent pair one may choose nonzero arrows
\[
 x_i:P_{s_i}\longrightarrow P_{s_{i+1}},
 \qquad
 y_i:P_{s_{i+1}}\longrightarrow P_{s_i}
\]
such that
\begin{equation}\label{eq:source-zigzag-composition}
 \nu_{i+1}:=x_i\circ y_i
\end{equation}
is the nonzero radical endomorphism of $P_{s_{i+1}}$.
\end{prop}

\begin{proof}
By the projective Hom-space display in the proof of
\cite[Theorem~6.6]{McRaeYang}, the only nonzero Hom spaces between
indecomposable projectives are the scalar endomorphisms of the wall
projectives, the two-dimensional endomorphism spaces of the positive-index
non-wall projectives, and the one-dimensional Hom spaces between adjacent
vertices in a reflection chain.  All Hom spaces between distinct reflection
chains and all Hom spaces between a wall projective and a different
indecomposable projective vanish.  Translating the labels in that display
into \eqref{eq:reflection-chain} gives exactly the dimensions stated above.

The nonvanishing of the backtracking composite is a separate input from
these dimension statements.  More concretely, for every adjacent pair
there are unique integers $n\ge1$ and $1\le\rho\le p-1$ such that
\[
 s_{i+1}=np+\rho,\qquad s_i=np-\rho.
\]
(The remainder $\rho$ alternates between the initial chain label and its
reflection, so it should not be confused with the fixed parameter used to
name the chain.)  Item~(3) in the Hom-space display in the proof of
\cite[Theorem~6.6]{McRaeYang} gives, in its notation,
\[
 \Hom(P_{np+\rho},P_{np-\rho})=\C F^-_{np+\rho},\qquad
 \Hom(P_{np-\rho},P_{np+\rho})=\C F^+_{np-\rho},
\]
and, crucially, the explicit endomorphism-space basis
\begin{equation}\label{eq:MY-backtracking-basis}
 \End(P_{np+\rho})
 =\C\id_{P_{np+\rho}}
  \oplus
  \C\bigl(F^+_{np-\rho}\circ F^-_{np+\rho}\bigr).
\end{equation}
The same proof later identifies the corresponding tilting backtracking
composition with a nonzero multiple of this second basis element.  We take
\[
 x_i:=F^+_{np-\rho},\qquad y_i:=F^-_{np+\rho}.
\]
Equation~\eqref{eq:MY-backtracking-basis} therefore gives
$x_i\circ y_i\ne0$ directly; this nonvanishing is not inferred from the
Hom-space dimensions.

Here and below, $\rad\End(P)$ denotes the Jacobson radical of the
finite-dimensional algebra $\End(P)$.  At a positive-index vertex
$P_{s_{i+1}}$ the projective Hom computation gives
$\dim_\C\End(P_{s_{i+1}})=2$.  Since $P_{s_{i+1}}$ is indecomposable of
finite length, Fitting's lemma implies that its endomorphism algebra $A$ is
local.  Hence $A/\rad A$ is a finite-dimensional division algebra over the
algebraically closed field $\C$, so $A/\rad A\cong\C$.  It follows that
$\dim_\C\rad A=1$.

The nonzero backtracking composite is not invertible.  Indeed, if
$x_i\circ y_i$ were invertible, then
\[
 g:=(x_i\circ y_i)^{-1}\circ x_i:P_{s_i}\longrightarrow P_{s_{i+1}}
\]
would satisfy $g\circ y_i=\id_{P_{s_{i+1}}}$.  Hence $y_i$ would be a
split monomorphism and $P_{s_{i+1}}$ a direct summand of the
indecomposable object $P_{s_i}$, forcing
$P_{s_i}\cong P_{s_{i+1}}$, contrary to their distinct simple heads.
Hence the composite is a nonzero element of the one-dimensional Jacobson
radical and therefore spans it.  Since both adjacent Hom spaces
are one-dimensional, any other nonzero choice differs by nonzero scalars and
has the same nonvanishing property.  We make no assertion here about the
opposite composite $y_i\circ x_i$; only the positive-index radical loop
$x_i\circ y_i$ will be used below.
\end{proof}

The category is $\Z/2\Z$-graded: if $M$ lies in degree
$\bar\epsilon$, we take $\epsilon\in\{0,1\}$ as its representative; then
the $h_0$-eigenvalues on $M$ lie in $\epsilon+2\Z$.  Besides the standard twist, McRae and Yang use the non-standard twist
\begin{equation}\label{eq:minus-twist}
 \theta^-_M=(-1)^\epsilon e^{2\pi iL^{\mathrm{aff}}_0},
\end{equation}
with the normalization of \cite[Theorem~5.4]{McRaeYang}.  Their functor
$F_{p,q}$ preserves this twist by \cite[Theorem~7.15]{McRaeYang}.  Our representative
$\epsilon\in\{0,1\}$ and the sign in \eqref{eq:minus-twist} agree with
the parity convention used in that theorem.

For later use, we distinguish the ordered-pair notation.  Write
$F^{(p,q)}:=F_{p,q}$ for the functor normalized by the $(2,1)$ Kac object
for the ordered pair $(p,q)$.  Exchanging the parameters gives the
ordered-pair construction $F^{(q,p)}$; under the normalization of
\cite[Theorem~7.15]{McRaeYang}, its distinguished $(2,1)$ Kac object is
$\mathcal K^{(q,p)}_{2,1}\cong\mathcal K^{(p,q)}_{1,2}$, and this is the
second functor denoted $F_{q,p}$ there.  We shall use this identification
in Corollary~\ref{cor:transposed-reduction}.

For $n\geq1$ and $1\leq r\leq p-1$, the projective $P_{np+r}$ is
logarithmic by \cite[Theorem~4.13]{McRaeYang}.  Fix once and for all a nonzero generator
$N_{np+r}\in\rad\End(P_{np+r})$; no normalization of this generator is
assumed.  Proposition~\ref{prop:source-Hom-package} gives
\[
 \End(P_{np+r})=\C\id\oplus\C N_{np+r},\qquad
 \rad\End(P_{np+r})=\C N_{np+r}.
\]
The radical of this finite-dimensional local algebra is a two-sided
nilpotent ideal.  Hence
$N_{np+r}^2\in\rad\End(P_{np+r})=\C N_{np+r}$, say
$N_{np+r}^2=aN_{np+r}$.  Nilpotence forces $a=0$.
Thus
\[
 N_{np+r}^2=0.
\]
For an affine module $M$ and $\xi\in\C/\Z$, write $M^{[\xi]}$ for the sum
of generalized $L^{\mathrm{aff}}_0$-weight spaces whose weights lie in
$\xi$.  We now isolate the relation between the projective radical and the
nilpotent part of the conformal Hamiltonian.

\begin{lem}[Radical and affine conformal nilpotent]\label{lem:source-L0-radical}
Let $P=P_{np+r}$ with $n\geq1$ and $1\leq r\leq p-1$, and put
$N_P:=N_{np+r}$.  Let $\bar\epsilon$ be the $\mathbb Z/2\mathbb Z$-degree of
$P$ and let $\epsilon\in\{0,1\}$ be its chosen representative.  Because
$P$ is grading restricted,
\[
 P=\bigoplus_{\lambda\in\C}P^{\mathrm{gen}}_\lambda
\]
is an algebraic direct sum, and every vector has only finitely many
nonzero generalized-eigenvalue components.  For each generalized eigenvalue $\lambda$, put
\[
 D_\lambda:=
 (L^{\mathrm{aff}}_0-\lambda\,\id)|_{P^{\mathrm{gen}}_\lambda},
 \qquad
 D_P:=\bigoplus_\lambda D_\lambda.
\]
The operators $D_\lambda$ assemble to a
$V^k(\mathfrak{sl}_2)$-module endomorphism $D_P$, which is locally
nilpotent in the algebraic sense: for every $v\in P$ there exists $N(v)$
such that $D_P^{N(v)}v=0$.  Since $P$ is logarithmic, $D_P\ne0$.
Consequently
\[
 D_P=\delta_PN_P\qquad(\delta_P\in\C^\times).
\]
If $\xi_P\in\C/\Z$ is the unique conformal-weight residue of $P$, then
\begin{equation}\label{eq:source-twist-nilpotent}
 \theta^-_P=\lambda_P(\id+2\pi iD_P)
 =\lambda_P(\id+\gamma_PN_P),
 \qquad
 \lambda_P=(-1)^\epsilon e^{2\pi i\xi_P},\qquad
 \gamma_P=2\pi i\,\delta_P\neq0 .
\end{equation}
\end{lem}

\begin{proof}
Let $v\in V^k(\mathfrak{sl}_2)$ be conformally homogeneous and put
\[
 \Delta=\wt v-m-1.
\]
The universal affine vertex algebra is $\Z_{\ge0}$-graded, so
$\wt v\in\Z_{\ge0}$ and therefore $\Delta\in\Z$.  The commutator formula
$[L^{\mathrm{aff}}_0,v_m]=\Delta\,v_m$ implies
\[
 \bigl(L^{\mathrm{aff}}_0-(\lambda+\Delta)\bigr)^Nv_m
 =v_m\bigl(L^{\mathrm{aff}}_0-\lambda\bigr)^N
 \qquad(N\ge1)
\]
on $P^{\mathrm{gen}}_\lambda$.  In particular,
\[
 D_{\lambda+\Delta}v_m=v_mD_\lambda.
\]
Thus the operators $D_\lambda$ assemble to an affine-module endomorphism
$D_P$.

Write $D_P=a\id+bN_P$.  If $a\neq0$, then
\[
 D_P=a\bigl(\id+(b/a)N_P\bigr)
\]
is invertible because $N_P^2=0$.  This contradicts local nilpotence:
if $0\ne v\in P$ and $D_P^m v=0$, invertibility of $D_P^m$ would force
$v=0$.  Hence $a=0$.  Since $P$ is logarithmic by
\cite[Theorem~4.13]{McRaeYang}, $D_P\neq0$, so $b=\delta_P\in\C^\times$.

Finally, every vector of $P$ has only finitely many generalized
$L^{\mathrm{aff}}_0$-eigencomponents, so the regrouping
\[
 P=\bigoplus_{\xi\in\C/\Z}P^{[\xi]}
\]
is an algebraic direct sum.  It is a direct sum of affine submodules
because every affine mode changes
conformal weight by an integer.  If two residue summands were nonzero,
one of them and the direct sum of all the remaining summands would give a
nontrivial direct-sum decomposition of $P$.  Indecomposability therefore
forces exactly one summand to be nonzero; denote its class by $\xi_P$.
Here $e^{2\pi i\xi_P}$ means $e^{2\pi i\lambda}$ for any representative
$\lambda$ of $\xi_P$.  Set
\[
 S_P:=L^{\mathrm{aff}}_0-D_P.
\]
On $P^{\mathrm{gen}}_\lambda$ one has $S_P=\lambda\id$.  Since every
actual generalized eigenvalue of $P$ belongs to the unique residue class
$\xi_P$, it follows algebraically, on each generalized eigenspace and hence
on their direct sum, that
\[
 e^{2\pi iS_P}=e^{2\pi i\xi_P}\id_P.
\]
Moreover $D_P^2=\delta_P^2N_P^2=0$, so
\[
 e^{2\pi iD_P}=\id+2\pi iD_P
 =\id+2\pi i\delta_PN_P.
\]
Since $S_P$ and $D_P$ commute, combining the two exponentials with
\eqref{eq:minus-twist} proves \eqref{eq:source-twist-nilpotent}.
\end{proof}

\subsection{One-row Virasoro Kac modules and conformal-residue sectors}

Let $\Kac_{r,s}$ denote the Virasoro Kac module of central charge
$c_{p,q}$, and abbreviate
\[
 K_j:=\Kac_{j,1},
 \qquad
 S_j:=L(c_{p,q},h_{j,1}),
 \qquad
 h_j:=h_{j,1}.
\]
Thus
\begin{equation}\label{eq:weights}
 h_j=\frac{(qj-p)^2-(p-q)^2}{4pq}.
\end{equation}
For the arithmetic formulas below only, we extend $h_j$ to every
$j\in\Z$ by \eqref{eq:weights}.  No Virasoro module $K_j$ or $S_j$ is
attached to a non-positive label.  The symbol $K_0=0$ will be used only as
a bookkeeping convention in recursion formulas; it is not a Virasoro Kac
module.  The labels $S_j$
are likewise labels rather than an assertion that all $S_j$ are pairwise
nonisomorphic; the only collision relevant below is isolated explicitly
in Corollary~\ref{cor:positive-label-collisions}.

We use the category $\Oc_{c_{p,q}}$ of McRae--Sopin: its objects are the
finite-length grading-restricted generalized Virasoro modules of central
charge $c_{p,q}$ whose simple composition factors are irreducible quotients
of reducible Verma modules; see \cite[Theorem~3.4]{McRaeSopin}.  It is a
$\C$-linear abelian category.  If $S,T$ are simple highest-weight Virasoro
modules, then $\dim_\C\Hom(S,T)\le1$: a morphism is determined by the image
of a nonzero highest-weight vector of lowest conformal weight and can be
nonzero only when the two simples are isomorphic.  We prove Hom-finiteness
by induction on
\[
 \ell(M)+\ell(N).
\]
If both $M$ and $N$ are simple, the preceding observation gives the base
case.  Suppose $\ell(M)+\ell(N)>2$.  If $\ell(N)>1$, choose
\[
 0\longrightarrow N'\longrightarrow N\longrightarrow T\longrightarrow0
\]
with $T$ simple.  Left exactness of $\Hom(M,-)$ gives
\[
 0\longrightarrow\Hom(M,N')\longrightarrow\Hom(M,N)
 \longrightarrow\Hom(M,T).
\]
Both outer spaces are finite dimensional by induction, since
$\ell(M)+\ell(N')<\ell(M)+\ell(N)$ and
$\ell(M)+1<\ell(M)+\ell(N)$.  Hence $\Hom(M,N)$ is finite dimensional.
If $\ell(N)=1$, then $\ell(M)>1$; choose
\[
 0\longrightarrow M'\longrightarrow M\longrightarrow S\longrightarrow0
\]
with $S$ simple.  The contravariant left-exact functor $\Hom(-,N)$ gives
\[
 0\longrightarrow\Hom(S,N)\longrightarrow\Hom(M,N)
 \longrightarrow\Hom(M',N),
\]
and again both outer terms are finite dimensional by induction.  Thus all
Hom spaces in $\Oc_{c_{p,q}}$ are finite dimensional.

Every finite-length object is a finite direct sum of indecomposables, by
induction on its length.  Fitting's lemma implies that the endomorphism
algebra of each indecomposable finite-length object is local.  Together
with Hom-finiteness, this proves that $\Oc_{c_{p,q}}$ is Krull--Schmidt.
We call an object \emph{logarithmic} when $L_0$ is not semisimple.

Every one-row Kac module $K_j$ is ordinary, hence $L_0$-semisimple.
The one-row structure results of \cite{McRaeSopin} give that $K_{np}$ is
simple and, for $n\geq0$ and $1\leq r\leq p-1$,
\begin{equation}\label{eq:Kac-length-two}
 0\longrightarrow S_{(n+2)p-r}\longrightarrow K_{np+r}
 \longrightarrow S_{np+r}\longrightarrow0
\end{equation}
is non-split; see the $s=1$ specialization of the boundary and bulk
composition series in \cite[Section~2.3]{McRaeSopin}, together with
\cite[Remark~2.2]{McRaeSopin}.  Since $K_{np+r}$ is generated by its
lowest-weight vector, its simple head is $S_{np+r}$.  The one-row
composition series has length two with remaining factor
$S_{(n+2)p-r}$, which fixes the orientation in
\eqref{eq:Kac-length-two}.  We shall repeatedly use the reflected
specialization obtained by writing
$np-r=(n-1)p+(p-r)$: for $n\ge1$,
\begin{equation}\label{eq:Kac-reflected-length-two}
 0\longrightarrow S_{np+r}\longrightarrow K_{np-r}
 \longrightarrow S_{np-r}\longrightarrow0.
\end{equation}
By \cite[Theorem~5.10]{McRaeSopin}, the individual object $K_2$ is rigid
and self-dual.  Thus $K_2^\vee\simeq K_2$, and its rigidity gives natural
adjunctions
\[
 \Hom(K_2\boxtimes X,Y)\cong\Hom(X,K_2^\vee\boxtimes Y),
 \qquad
 \Hom(X,K_2\boxtimes Y)\cong\Hom(K_2^\vee\boxtimes X,Y).
\]
Hence $K_2\boxtimes-$ is simultaneously a left and a right adjoint.  It is
therefore both left exact and right exact, hence exact.  Their fusion theorem gives
\begin{equation}\label{eq:K2-fusion}
 0\longrightarrow K_{j-1}\longrightarrow K_2\boxtimes K_j
 \longrightarrow K_{j+1}\longrightarrow0,
\end{equation}
which splits if and only if $p\nmid j$
\cite[Theorem~6.5]{McRaeSopin}.

We next collect the elementary conformal-weight arithmetic needed
throughout the rest of the paper.

\begin{lem}[Weight arithmetic]\label{lem:weight-arithmetic}
For arbitrary integers $u,v$,
\begin{equation}\label{eq:weight-difference}
 h_u-h_v=\frac{(u-v)(q(u+v)-2p)}{4p}.
\end{equation}
In particular,
\begin{equation}\label{eq:weight-equality}
 h_j=h_\ell
 \quad\Longleftrightarrow\quad
 j=\ell\ \text{ or }\ q(j+\ell)=2p.
\end{equation}
For $n\geq1$ and $1\leq r\leq p-1$,
\begin{align}
 h_{np+r-1}-h_{np-r+1}&=(r-1)(nq-1),\label{eq:parallel-one}\\
 h_{np+r+1}-h_{np-r-1}&=(r+1)(nq-1),\label{eq:parallel-two}\\
 h_{np+r-1}-h_{np-r-1}&=r(nq-1)-\frac{rq}{p},\label{eq:cross-one}\\
 h_{np+r+1}-h_{np-r+1}&=r(nq-1)+\frac{rq}{p}.\label{eq:cross-two}
\end{align}
We shall also use repeatedly
\begin{align}
 h_{np+r}-h_{np-r}&=r(nq-1),\label{eq:strip-central-difference}\\
 h_{(n+2)p-r}-h_{np+r}
   &=(p-r)((n+1)q-1).\label{eq:strip-upper-difference}
\end{align}
The first two differences in the preceding four-line display are integral
and the last two are nonintegral.
For the wall calculation put
\[
 A=(n-1)p,\quad B=(n+1)p,\quad C=(n-1)p+2,\quad D=(n+1)p-2.
\]
Then
\begin{align}
 h_B-h_A&=p(nq-1),\label{eq:wall-parallel-one}\\
 h_D-h_C&=(p-2)(nq-1),\label{eq:wall-parallel-two}\\
 h_D-h_A&=(p-1)(nq-1)-\frac{(p-1)q}{p},\label{eq:wall-cross-one}\\
 h_B-h_C&=(p-1)(nq-1)+\frac{(p-1)q}{p}.\label{eq:wall-cross-two}
\end{align}
The crossed wall differences are nonintegral for $p\geq3$.  If $p=2$,
their fractional parts are $\pm q/2$, which are nonzero modulo $\Z$
because $q$ is odd.
\end{lem}

\begin{proof}
Formula \eqref{eq:weight-difference} follows by subtracting
\eqref{eq:weights}.  Equality of weights gives
$(qj-p)^2=(q\ell-p)^2$, hence \eqref{eq:weight-equality}.  The remaining
formulas follow from \eqref{eq:weight-difference}; coprimality of $p$ and
$q$ gives the asserted nonintegrality.
\end{proof}

\begin{cor}[Positive-label collisions]\label{cor:positive-label-collisions}
For $j,\ell>0$,
\[
 S_j\cong S_\ell
 \quad\Longleftrightarrow\quad
 \bigl(j=\ell\bigr)
 \ \text{or}\ 
 \bigl(q=2\ \text{and}\ j+\ell=p\bigr).
\]
In the second case $p$ is odd.
\end{cor}

\begin{proof}
For fixed central charge $c_{p,q}$, an irreducible highest-weight Virasoro
module is uniquely determined by its lowest conformal weight.  Thus \eqref{eq:weight-equality} reduces
the question to $j=\ell$ or $q(j+\ell)=2p$.  In the second case
coprimality and $q\ge2$ force $q=2$ and $j+\ell=p$.  Conversely these
relations give $h_j=h_\ell$, hence $S_j\cong S_\ell$.  This is an
equality-of-lowest-weights statement, not merely a congruence modulo
$\Z$.  Finally $\gcd(p,2)=1$ implies that $p$ is odd.
\end{proof}

We shall use two different notions below.  Equality $h_j=h_\ell$ governs
isomorphism of simple highest-weight modules, whereas congruence
$h_j\equiv h_\ell\pmod{\Z}$ governs the conformal-residue decomposition.
The latter is strictly weaker.

\begin{cor}[Residue separation]\label{cor:residue-separation}
Let $n\ge1$ and $1\le r\le p-1$, and put $a=np-r$, $b=np+r$.  Then
\[
 [h_{a+1}]=[h_{b-1}],\qquad
 [h_{a-1}]=[h_{b+1}]
 \quad\text{in }\C/\Z,
\]
and these two residue classes are distinct.  For the wall labels
\[
 A=(n-1)p,\quad B=(n+1)p,\quad
 C=(n-1)p+2,\quad D=(n+1)p-2,
\]
one has
\[
 [h_A]=[h_B],\qquad [h_C]=[h_D],
\]
and the two displayed residue classes are distinct (with the terms of
label $0$ omitted when $n=1$).  These assertions remain valid for $p=2$
with the interpretation in Lemma~\ref{lem:weight-arithmetic}.
\end{cor}

\begin{proof}
The congruences are exactly
\eqref{eq:parallel-one}--\eqref{eq:parallel-two} and
\eqref{eq:wall-parallel-one}--\eqref{eq:wall-parallel-two}.  Distinctness
follows from the nonintegrality of the crossed differences in
\eqref{eq:cross-one}--\eqref{eq:cross-two} and
\eqref{eq:wall-cross-one}--\eqref{eq:wall-cross-two}.
\end{proof}

\begin{lem}[Residue support of one-row Kac modules]\label{lem:Kac-residue-support}
For every $j\ge1$,
\[
 K_j=K_j^{[h_j]}.
\]
\end{lem}

\begin{proof}
Since the second Kac label is $s=1<q$, the one-row module
$K_j=\Kac_{j,1}$ lies in the Verma-quotient regime described in
\cite[Remark~2.2]{McRaeSopin}; in particular it is generated by its
lowest-weight vector of weight $h_j$.
Virasoro modes change conformal weight by integers, so every weight of
$K_j$ belongs to $h_j+\Z_{\ge0}$.
\end{proof}

\begin{lem}[Conformal-residue decomposition]\label{lem:block-projection}
For $M\in\Oc_{c_{p,q}}$ and
$\xi\in\C/\Z$, let $M^{[\xi]}$ be the sum of the generalized
$L_0$-weight spaces whose weights belong to $\xi$.  Then
\[
 M=\bigoplus_{\xi\in\C/\Z}M^{[\xi]}
\]
with only finitely many nonzero summands.  Each $M^{[\xi]}$ is a Virasoro
submodule, and the projection $M\mapsto M^{[\xi]}$ is exact.  Moreover,
\[
 \Hom(M^{[\xi]},N^{[\eta]})=0\qquad(\xi\neq\eta),
\]
and if $M,N$ are supported in distinct residue classes, then the Yoneda extension group vanishes:
\begin{equation}\label{eq:cross-residue-Ext}
 \Ext^1_{\Oc_{c_{p,q}}}(M,N)=0.
\end{equation}
We refer to the summands $M^{[\xi]}$ as \emph{conformal-residue sectors}.
No assertion about the indecomposable categorical blocks of
$\Oc_{c_{p,q}}$ is intended.
\end{lem}

\begin{proof}
Since $L_m$ changes generalized conformal weight by the integer $-m$, each
$M^{[\xi]}$ is a Virasoro submodule.  A simple grading-restricted module is supported in one residue class.
For finite support, take a short exact sequence
$0\to A\to M\to S\to0$ with $S$ simple.  A generalized weight vector of
$M$ whose residue does not occur in $S$ maps to zero and therefore lies in
$A$.  Induction on Jordan--H\"older length thus shows directly that the
residue support of $M$ is contained in the union of the finite residue
supports of $A$ and $S$.

We now prove exactness of residue projection independently.  Let
$B\twoheadrightarrow C$ be a surjective Virasoro homomorphism and take
$c\in C^{[\xi]}$.  Choose a lift $b\in B$.  By local finiteness the cyclic
$\C[L_0]$-space $W_b:=\C[L_0]b$ is finite dimensional.  Primary
decomposition for $L_0|_{W_b}$, followed by grouping the finitely many
generalized eigenvalues modulo $\Z$, gives a finite decomposition
$b=\sum_\eta b_\eta$ with $b_\eta\in B^{[\eta]}$.  Since the
homomorphism commutes with $L_0$,
the components with $\eta\neq\xi$ cannot contribute to $c$; hence
$b_\xi$ maps to $c$.  Thus $B^{[\xi]}\twoheadrightarrow C^{[\xi]}$.  For a short exact
sequence $0\to A\to B\to C\to0$ one also has
\[
 \ker\bigl(B^{[\xi]}\to C^{[\xi]}\bigr)
 =A\cap B^{[\xi]}=A^{[\xi]},
\]
so residue projection is exact.  The same weight preservation gives the
cross-residue Hom vanishing.

Finally, apply exact residue projection to an extension
\[
 0\longrightarrow N^{[\eta]}\longrightarrow E\longrightarrow
 M^{[\xi]}\longrightarrow0,\qquad \xi\neq\eta .
\]
It yields
\[
 E^{[\xi]}\xrightarrow{\sim}M^{[\xi]},\qquad
 N^{[\eta]}\xrightarrow{\sim}E^{[\eta]},
\]
and all other residue summands of $E$ vanish.  Hence
$E=E^{[\xi]}\oplus E^{[\eta]}$.  The inverse of
$E^{[\xi]}\xrightarrow{\sim}M^{[\xi]}$, followed by the inclusion
$E^{[\xi]}\hookrightarrow E$, is a Virasoro-module section of
$E\twoheadrightarrow M^{[\xi]}$.  Thus the extension splits.  This proves
\eqref{eq:cross-residue-Ext}.
\end{proof}

\begin{lem}[Exponential detects logarithmicity]\label{lem:exponential-logarithmic}
On every finite-dimensional generalized $L_0$-eigenspace, the exponential
is understood algebraically.  If $M$ is grading restricted, then
$e^{2\pi iL_0}$ is diagonalizable as a linear operator if and only if
$L_0$ is semisimple.
\end{lem}

\begin{proof}
Set $E=e^{2\pi iL_0}$ and take
$x\in M^{\mathrm{gen}}_\lambda$.  Since $L_0-\lambda$ is nilpotent on
the finite-dimensional generalized eigenspace containing $x$, there is an
$N$ such that
\[
 (E-e^{2\pi i\lambda})^N x=0.
\]
Assume first that $E$ is diagonalizable on $M$.  Writing $x$ as a finite
sum of $E$-eigenvectors, the displayed relation forces every eigenvalue
occurring in $x$ to equal $e^{2\pi i\lambda}$.  Hence
$Ex=e^{2\pi i\lambda}x$.  Writing
$L_0=\lambda\id+D$ on $M^{\mathrm{gen}}_\lambda$ therefore gives
$e^{2\pi iD}=\id$.  But
\[
 e^{2\pi iD}-\id=DQ(D),
 \qquad Q(0)=2\pi i\ne0,
\]
where $Q(D)$ is a polynomial because $D$ is nilpotent.  Thus $Q(D)$ is
invertible and $D=0$.  Hence $L_0$ is semisimple on every generalized
eigenspace and therefore on $M$.

Conversely, if $L_0$ is semisimple, then $M$ is the algebraic direct sum
of its $L_0$-eigenspaces and $E$ acts by a scalar on each of them.  Thus
$E$ is diagonalizable.
\end{proof}

We shall also repeatedly use the following elementary finite-length
cancellation principle.

\begin{lem}[Finite-length cancellation]\label{lem:length-cancellation}
If $M=A\oplus B$ has finite length and $M$ and $A$ have the same
Jordan--H\"older multiplicities, then $B=0$.
\end{lem}

\begin{proof}
Additivity of Jordan--H\"older multiplicities under direct sums forces all
multiplicities of $B$ to be zero; equivalently $\ell(B)=0$.
\end{proof}

\subsection{Nakano's extension subcategory and Ext conventions}

Set
\[
 p_+=\min\{p,q\},\qquad p_-=\max\{p,q\}.
\]
Because $p,q\ge2$ are coprime, $p\ne q$, and hence
\[
 2\le p_+<p_-,\qquad \gcd(p_+,p_-)=1.
\]
Thus the standing numerical hypotheses of \cite{Nakano} are satisfied.
Since the Virasoro central charge is symmetric in its two parameters,
\begin{equation}\label{eq:central-charge-symmetry}
 c_{p,q}=c_{q,p}=c_{p_+,p_-}.
\end{equation}
Thus all Virasoro modules considered below have the same central charge as
those in Section~5 of \cite{Nakano}.  All references to Nakano in this paper
use \nolinkurl{arXiv:2305.12448v2} (9 August 2026).

There is a minor but important point of scope.  We do not identify
$\Oc_{c_{p,q}}$, or any subcategory of it, with the ambient category of
\cite[Definition~3.14]{Nakano}.  Instead, we isolate below a full Serre
subcategory of $\Oc_{c_{p,q}}$ containing all finite-length Virasoro
modules needed in the one-row argument, and we formulate the required
extension statements directly in that category.  This avoids any
dependence on the literal support equality in Definition~3.14 while
retaining the module-theoretic content of the calculations in Section~5
of \cite{Nakano}.

\begin{lem}[Extension closure at fixed central charge]
\label{lem:gr-extension}
Let
\[
 0\longrightarrow A\longrightarrow B\longrightarrow C\longrightarrow0
\]
be a short exact sequence in the category of Virasoro modules of fixed
central charge $c$, so that the central element acts as $c\,\id$ on all
three terms.  If $A$ and $C$ are grading-restricted generalized Virasoro
modules, then so is $B$.  Moreover, for every $\lambda\in\C$,
\[
 0\longrightarrow A^{\mathrm{gen}}_\lambda
 \longrightarrow B^{\mathrm{gen}}_\lambda
 \longrightarrow C^{\mathrm{gen}}_\lambda
 \longrightarrow0
\]
is exact.
\end{lem}

\begin{proof}
Let $b\in B$, and write $\bar b$ for its image in $C$.  Local finiteness
of $L_0$ on $C$ gives a nonzero polynomial $p$ with
$p(L_0)\bar b=0$, so $p(L_0)b\in A$.  Local finiteness on $A$ then gives a
nonzero polynomial $q$ with $q(L_0)p(L_0)b=0$.  Thus $L_0$ acts locally
finitely on $B$, and hence $B$ is the algebraic direct sum of its generalized
eigenspaces.

Fix $c\in C^{\mathrm{gen}}_\lambda$ and choose a lift $b\in B$.  The
finite-dimensional $\C[L_0]$-span of $b$ has only finitely many generalized
eigenvalues.  On this finite-dimensional space the spectral projector onto
the $\lambda$-primary summand is a polynomial in $L_0$.  Applying that
projector to $b$ gives a lift in $B^{\mathrm{gen}}_\lambda$ of $c$.  This
proves surjectivity on the $\lambda$-generalized eigenspaces; injectivity at
the left and exactness in the middle are inherited from the original short
exact sequence.  Therefore
\[
 0\longrightarrow A^{\mathrm{gen}}_\lambda
 \longrightarrow B^{\mathrm{gen}}_\lambda
 \longrightarrow C^{\mathrm{gen}}_\lambda\longrightarrow0
\]
is exact and
\[
 \dim B^{\mathrm{gen}}_\lambda
 =\dim A^{\mathrm{gen}}_\lambda+\dim C^{\mathrm{gen}}_\lambda<\infty.
\]
The generalized-eigenspace exact sequence also gives
\[
 \operatorname{Spec}(B)=\operatorname{Spec}(A)\cup\operatorname{Spec}(C).
\]
Indeed, nonzero generalized eigenspaces of $A$ inject into those of $B$,
and surjectivity onto the generalized eigenspaces of $C$ shows that every
weight in $\operatorname{Spec}(C)$ occurs in $B$.  Conversely, if
$B^{\mathrm{gen}}_\lambda\ne0$, then either its image in
$C^{\mathrm{gen}}_\lambda$ is nonzero or it meets
$A^{\mathrm{gen}}_\lambda$ nontrivially.

It remains to verify lower truncation, and here no global finiteness
assumption on the set of lowest weights is needed.  Fix $\lambda\in\C$.
Since $A$ and $C$ are grading restricted, there are integers $N_A,N_C$
such that
\[
 A^{\mathrm{gen}}_{\lambda-n}=0\quad(n\ge N_A),
 \qquad
 C^{\mathrm{gen}}_{\lambda-n}=0\quad(n\ge N_C).
\]
Exactness on generalized eigenspaces therefore gives
\[
 B^{\mathrm{gen}}_{\lambda-n}=0
 \qquad\bigl(n\ge\max\{N_A,N_C\}\bigr).
\]
If $b\in B^{\mathrm{gen}}_\lambda$, then $[L_0,L_n]=-nL_n$ implies
$L_nb\in B^{\mathrm{gen}}_{\lambda-n}$, so $L_nb=0$ for all sufficiently
large $n$.  Together with local finiteness and finite-dimensional
generalized eigenspaces proved above, this is precisely the grading
restriction for $B$.
\end{proof}


\begin{lem}[Finite-length generalized-weight control]
\label{lem:Nakano-ambient}
Let $M\in\Oc_{c_{p,q}}$.  Then $L_0$ acts locally finitely on $M$, every
generalized $L_0$-eigenspace is finite dimensional, and
\[
 \operatorname{Spec}(M)\subset H_0+\Z_{\ge0}
\]
for a finite set $H_0\subset\C$.  The restricted contragredient $M'$ has
the same properties and finite length, and the canonical map
$M\to(M')'$ is an isomorphism.
\end{lem}

\begin{proof}
The first assertions follow from a composition series and repeated
application of Lemma~\ref{lem:gr-extension}.  Define
\[
 M':=\bigoplus_{\lambda\in\C}
       \bigl(M^{\mathrm{gen}}_\lambda\bigr)^*
\]
using the Virasoro anti-involution $L_n\mapsto L_{-n}$ and $C\mapsto C$.
Since $L_0$ is fixed by this anti-involution,
\[
 (M')^{\mathrm{gen}}_\lambda
 \cong \bigl(M^{\mathrm{gen}}_\lambda\bigr)^*.
\]
Thus $M'$ has the same support, finite generalized-weight multiplicities,
and lower truncation.  Restricted duality is exact on modules with
finite-dimensional generalized weight spaces.  Moreover every simple lowest-weight Virasoro module occurring here is
self-contragredient.  Indeed, the restricted dual of $L(c_{p,q},h)$ is
again a simple lowest-weight module of central charge $c_{p,q}$ and lowest
weight $h$, and uniqueness of the simple quotient of the corresponding
Verma module gives
\[
 L(c_{p,q},h)'\cong L(c_{p,q},h).
\]
Thus a finite composition series of $M$ dualizes to a finite composition
series of $M'$ with the same simple isomorphism classes.  In particular,
$M'\in\Oc_{c_{p,q}}$, and $M\in\mathscr C$ will imply
$M'\in\mathscr C$ once $\mathscr C$ is defined below.  Finally, finite
dimensionality of every generalized weight space gives the canonical
biduality isomorphism $M\xrightarrow{\sim}(M')'$.
\end{proof}

Define the full subcategory
\[
\begin{aligned}
 \mathscr C^{\mathrm{fl}}_{p_+,p_-}
 :=\bigl\{M\in\Oc_{c_{p,q}}:\;&\text{every simple Jordan--H\"older
 constituent of $M$}\\
 &\text{has lowest weight in }H_{p_+,p_-}\bigr\},
\end{aligned}
\]
where $H_{p_+,p_-}$ is the set of Felder weights used in
\cite[Definition~5.1]{Nakano}.  We write
$\mathscr C:=\mathscr C^{\mathrm{fl}}_{p_+,p_-}$.

\begin{lem}\label{lem:Nakano-Serre}
The full subcategory $\mathscr C$ is a Serre subcategory of
$\Oc_{c_{p,q}}$.  Every object of $\mathscr C$ has a finite socle series.
\end{lem}

\begin{proof}
Subobjects and quotients of a finite-length object have only
Jordan--H\"older factors occurring in the original object.  In a short
exact sequence the multiset of Jordan--H\"older factors of the middle term
is the union of those of the end terms.  This proves the Serre property.  In particular, $\mathscr C$ is an
abelian extension-closed full subcategory of $\Oc_{c_{p,q}}$.
If $M/\Soc^j(M)\ne0$, then the finite-length quotient has a simple
submodule, so the socle filtration strictly increases until it reaches
$M$, after at most $\ell(M)$ steps.
\end{proof}

\begin{lem}[Yoneda extensions in the Serre subcategory]
\label{lem:Nakano-exact-compatibility}
For $X,Y\in\mathscr C$, the inclusion
$\mathscr C\hookrightarrow\Oc_{c_{p,q}}$ induces an isomorphism
\[
 \Ext^1_{\mathscr C}(X,Y)
 \xrightarrow{\ \sim\ }
 \Ext^1_{\Oc_{c_{p,q}}}(X,Y)
\]
of Yoneda extension spaces.
\end{lem}

\begin{proof}
If
\[
 0\longrightarrow Y\longrightarrow E\longrightarrow X\longrightarrow0
\]
is exact in $\Oc_{c_{p,q}}$, the Serre property gives $E\in\mathscr C$.
Thus every ambient Yoneda extension already lies in $\mathscr C$.
Conversely, $\mathscr C$ is full, so equivalences of short exact sequences,
Baer sums, scalar pushouts, and splittings are computed by the same
Virasoro-module morphisms in the two categories.  The two Yoneda spaces
are therefore identical.
\end{proof}

Our orientation convention is
\begin{equation}\label{eq:Ext-convention}
 \Ext^1_{\mathscr C}(X,Y)
 \quad\text{parametrizes}\quad
 0\longrightarrow Y\longrightarrow E\longrightarrow X\longrightarrow0.
\end{equation}
This is also the quotient-first convention used immediately after
\cite[Definition~5.1]{Nakano}.

\begin{lem}[Six-term Hom--Yoneda exactness]
\label{lem:prelim-Hom-Yoneda}
Let
\[
 0\longrightarrow A\longrightarrow B\longrightarrow C\longrightarrow0
\]
be a short exact sequence in a $\C$-linear abelian category.  For every
object $X$ there are sequences
\[
\begin{aligned}
0\longrightarrow&\Hom(C,X)\longrightarrow\Hom(B,X)
 \longrightarrow\Hom(A,X)\\
 \longrightarrow&\Ext^1(C,X)\longrightarrow\Ext^1(B,X)
 \longrightarrow\Ext^1(A,X)
\end{aligned}
\]
and
\[
\begin{aligned}
0\longrightarrow&\Hom(X,A)\longrightarrow\Hom(X,B)
 \longrightarrow\Hom(X,C)\\
 \longrightarrow&\Ext^1(X,A)\longrightarrow\Ext^1(X,B)
 \longrightarrow\Ext^1(X,C).
\end{aligned}
\]
The first sequence is exact through the displayed term
$\Ext^1(B,X)$, and the second is exact through the displayed term
$\Ext^1(X,B)$.  No assertion of exactness is made at either final
displayed $\Ext^1$-term.  Here $\Ext^1$ is Yoneda $\Ext^1$.
\end{lem}

\begin{proof}
The first connecting morphism is pushout along $A\to X$ and the second
is pullback along $X\to C$.  Exactness at the displayed $\Ext^1$-terms
through $\Ext^1(B,X)$ and $\Ext^1(X,B)$ is the usual Yoneda criterion:
an extension becomes split after pullback (respectively pushout) precisely
when it is obtained by pushout (respectively pullback) from the preceding
term.  Extending either sequence one step further would require the next
Yoneda obstruction (equivalently an $\Ext^2$ term when a derived-functor
long exact sequence is available), so no terminal surjectivity is being
claimed here.  The argument uses only short exact sequences and their
universal properties; no hypothesis of enough projectives or injectives
is involved.
\end{proof}

\begin{lem}[One-dimensional Yoneda Ext spaces]\label{lem:one-dimensional-ext}
Let $\mathscr E$ be a full extension-closed subcategory of a $\C$-linear
abelian category, and suppose
$\dim_\C\Ext^1_{\mathscr E}(X,Y)=1$.  Then all nonzero extension classes
of $X$ by $Y$ have isomorphic middle terms.
\end{lem}

\begin{proof}
If $\xi,\xi'\neq0$, then $\xi'=a\xi$ for some $a\in\C^\times$.
Scalar multiplication by $a$ is represented by pushout along
$a\,\id_Y$ (equivalently, by the corresponding pullback along an
automorphism of $X$).  Pushout or pullback along an endpoint
isomorphism does not change the isomorphism class of the middle term.
\end{proof}

We now prove the local Virasoro extension statements required below.
Although the notation follows Section~5 of \cite{Nakano}, the assertions
are formulated and verified in the Serre category $\mathscr C$.  This is
necessary because the ambient category in
\cite[Definition~3.14]{Nakano} imposes an additional global support
condition which is not part of the definition of $\mathscr C$.

Throughout the remainder of this local extension discussion, fix
\[
 \tau=(\alpha_1,\alpha_2,\alpha_3)\in\mathcal T_{p_+,p_-}
 \quad\text{with}\quad
 h_{\alpha_1}<h_{\alpha_2}<h_{\alpha_3}.
\]
This is the only case used later; Proposition~\ref{prop:Felder-labels}
will show that every one-row triple $\tau_{n,r}$ has this strict ordering.
Abbreviate
\[
 S_i:=L(h_{\alpha_i})\qquad(i=1,2,3).
\]

\begin{lem}[Comparison with the classical Virasoro category $\mathcal O$]
\label{lem:local-category-O}
Let $S,T\in\mathscr C$ be distinct simple lowest-weight Virasoro modules.
Write
$\mathcal O^{\mathrm{Vir}}_{c_{p,q}}$ for the classical Virasoro
category $\mathcal O$ used in \cite{BoeNakanoWiesner}.  Then every
extension of $S$ by $T$ in $\mathscr C$ belongs to
$\mathcal O^{\mathrm{Vir}}_{c_{p,q}}$, and every length-two extension
of these two simples in $\mathcal O^{\mathrm{Vir}}_{c_{p,q}}$ belongs
to $\mathscr C$.  Consequently
\[
 \Ext^1_{\mathscr C}(S,T)
 \cong
 \Ext^1_{\mathcal O^{\mathrm{Vir}}_{c_{p,q}}}(S,T).
\]
\end{lem}

\begin{proof}
Let
\[
 0\longrightarrow T\longrightarrow E\longrightarrow S\longrightarrow0
\]
be exact in $\mathscr C$.  The canonical nilpotent part $N_E$ of $L_0$ on $E$ is a
Virasoro-module endomorphism: on the generalized $\lambda$-eigenspace it
is $L_0-\lambda$, and $[L_0,L_n]=-nL_n$ intertwines these nilpotent
parts.  Since $L_0$ is semisimple on both endpoints,
\[
 N_E|_T=0,\qquad \operatorname{Im}N_E\subset T.
\]
Thus $N_E$ factors through the quotient $S=E/T$: there is a morphism
$\overline N:S\to T$ such that
\[
 N_E=(T\hookrightarrow E)\circ\overline N\circ(E\twoheadrightarrow S).
\]
Because $S$ and $T$ are distinct simples, $\Hom(S,T)=0$, so $N_E=0$.
Hence $L_0$ is semisimple on $E$.

Let $t$ be a lowest-weight generator of $T$, let $s$ be a lowest-weight
generator of $S$, and choose a lift $\widetilde s\in E$ of $s$.  Then
\[
 E=U(\Vir)t+U(\Vir)\widetilde s,
\]
so $E$ is finitely generated.  Moreover the positive Virasoro
subalgebra acts locally finitely.  Indeed, for a homogeneous vector
$v$ of weight $h+N$, every monomial in positive modes of total degree
larger than $N+C$ annihilates $v$, where $C$ is chosen so that all
weights of $E$ lie above $h-C$.  Only finitely many PBW monomials have
total positive degree at most $N+C$, so $U(\Vir_{>0})v$ is finite
dimensional.  These are the defining finiteness conditions for the classical Virasoro
category $\mathcal O$ used in \cite[Section~2.2]{BoeNakanoWiesner}; hence
$E\in\mathcal O^{\mathrm{Vir}}_{c_{p,q}}$.

Conversely, let $E$ be a length-two extension of $S$ by $T$ in the
classical category $\mathcal O$.  By definition it is a weight module,
so $L_0$ acts semisimply.  Its weight support is contained in the union
of the two simple supports and is therefore bounded below, while the
short exact sequence on each weight space gives finite-dimensional
weight spaces.  Thus $E$ is grading-restricted.  Since its only
Jordan--H\"older constituents are $S$ and $T$, it belongs to the Serre
subcategory $\mathscr C$.  The two classes of short exact sequences,
and the equivalence relation between them, are therefore identical for
these endpoints.
\end{proof}

\begin{lem}[Finite reduction for simple $\Ext^1$ in Virasoro category $\mathcal O$]
\label{lem:BNW-finite-reduction}
Let $L(\lambda)$ and $L(\mu)$ be fixed simple objects in one block of
$\mathcal O^{\mathrm{Vir}}_{c_{p,q}}$.  After enlarging a prescribed
finite interval of the block containing $\lambda$ and $\mu$, their
$\Ext^1$ may be computed in one of the finite highest-weight categories
used in Section~4 of \cite{BoeNakanoWiesner}.  More precisely, there is
a finite block, truncation, or quotient $\mathcal C_0$ containing the
chosen interval such that
\[
 \Ext^1_{\mathcal O^{\mathrm{Vir}}_{c_{p,q}}}
   \bigl(L(\lambda),L(\mu)\bigr)
 \cong
 \Ext^1_{\mathcal C_0}
   \bigl(L(\lambda),L(\mu)\bigr).
\]
The same finite category can be chosen simultaneously for finitely many
pairs of simples.
\end{lem}

\begin{proof}
There is nothing to prove for a finite block.  Suppose first that the
block is infinite with a minimal element.  Choose a truncation
$\mathcal C(\xi)$ in the sense of Section~4.1 of
\cite{BoeNakanoWiesner} which contains the prescribed finite interval.
The truncation is a full subcategory defined by its allowed composition
factors.  Hence every length-two extension of two simples in
$\mathcal C(\xi)$ again lies in $\mathcal C(\xi)$, and equivalences and
splittings of such short exact sequences are the same in the full block
and in the truncation.  Their Yoneda $\Ext^1$ groups are therefore
identical.

If the block is infinite with a maximal element, choose a finite coideal
$\Omega$ containing the prescribed interval and all vertices at
length-distance at most one from the finitely many endpoints under
consideration.  Proposition~4.2(C) of \cite{BoeNakanoWiesner}, with
$N\ge1$, identifies the full-block $\Ext^1$ groups with those in the
finite quotient category $\mathcal C(\Omega)$.  Enlarging $\Omega$ if
necessary treats finitely many pairs at once.  In each case the finite
weight poset contains the prescribed interval with its original order,
and the doubled-Hasse description of Section~4.4 applies there.
\end{proof}

\begin{lem}[Hasse positions of a strict Felder triple]
\label{lem:Felder-Hasse-chain}
Let $\lambda_i$ be the vertex of the Boe--Nakano--Wiesner weight poset
corresponding to $S_i=L(h_{\alpha_i})$.  Up to reversing the convention
for the partial order, one has two successive cover relations
\[
 \lambda_1\lessdot\lambda_2\lessdot\lambda_3.
\]
In particular, $\lambda_1$ and $\lambda_3$ are not joined by a Hasse
edge.
\end{lem}

\begin{proof}
For a strict Felder triple, \cite[Definition~5.9]{Nakano} provides the
ordinary nonsplit length-two highest-weight modules
$L(h_{\alpha_1},h_{\alpha_2})$ and
$L(h_{\alpha_2},h_{\alpha_3})$.  These are concrete Virasoro modules in
$L_{c_{p_+,p_-}}\text{-mod}$, hence are grading-restricted generalized
modules.  Their Jordan--H\"older factors are respectively $(S_1,S_2)$
and $(S_2,S_3)$.  Since the three lowest weights belong to
$H_{p_+,p_-}$, both middle terms lie in the Serre subcategory
$\mathscr C$.  Their defining short exact sequences remain nonsplit in
$\mathscr C$ because $\mathscr C$ is full.
Lemma~\ref{lem:local-category-O} therefore places these two nonsplit
extensions in the classical Virasoro category $\mathcal O$.  They
determine nonzero classes between $(S_1,S_2)$ and between $(S_2,S_3)$
in that category.

Apply Lemma~\ref{lem:BNW-finite-reduction} simultaneously to the three
vertices, choosing the finite model to contain the order interval joining
them.  The two nonzero extension classes remain nonzero in this finite
category.  By the doubled-Hasse description of
\cite[Section~4.4]{BoeNakanoWiesner} (equivalently
\cite[Theorem~10]{BoeNakanoWiesner}), the pairs
$(\lambda_1,\lambda_2)$ and $(\lambda_2,\lambda_3)$ are Hasse edges
there, and hence in the retained interval of the original block.

It remains only to determine the relative direction of the two edges; no
claim about embeddings of Fock modules is used.  In the Virasoro
highest-weight order underlying \cite{BoeNakanoWiesner}, an edge between
simple highest weights $h<h'$ can only have the direction corresponding
to a Verma embedding
\[
 M(c_{p,q},h')\hookrightarrow M(c_{p,q},h).
\]
Indeed, a nonzero homomorphism of Verma modules sends the highest-weight
vector to a singular vector, whose relative grade is $h'-h\in\Z_{\ge0}$;
the opposite direction would require a singular vector of negative
relative grade.  Since
\[
 h_{\alpha_1}<h_{\alpha_2}<h_{\alpha_3},
\]
the two Hasse edges consequently have the same orientation.  Thus, up to
reversing the convention for the partial order,
\[
 \lambda_1\lessdot\lambda_2\lessdot\lambda_3.
\]
In particular $\lambda_2$ is a comparable vertex strictly between the
two endpoints, so $\lambda_1$ and $\lambda_3$ cannot themselves be
joined by a cover relation.
\end{proof}

\begin{lem}[Simple extensions in the Felder triple]
\label{lem:local-simple-Ext}
For distinct simples in the above strict triple,
\[
 \Ext^1_{\mathscr C}(S_i,S_{i+1})
 \cong
 \Ext^1_{\mathscr C}(S_{i+1},S_i)
 \cong\C
 \qquad(i=1,2),
\]
and
\[
 \Ext^1_{\mathscr C}(S_1,S_3)
 =
 \Ext^1_{\mathscr C}(S_3,S_1)=0.
\]
No assertion about self-extensions is made here; the self-extension
vanishing needed for the one-row application is proved separately in
Lemma~\ref{lem:one-row-self-Ext} using primitive positive Kac
representatives.
\end{lem}

\begin{proof}
Lemma~\ref{lem:local-category-O} identifies the Yoneda groups for
distinct endpoints in $\mathscr C$ with those in the classical Virasoro
category $\mathcal O$.  Choose, by
Lemma~\ref{lem:BNW-finite-reduction}, one finite BNW model containing
$\lambda_1,\lambda_2,\lambda_3$ and their order interval.  By
Lemma~\ref{lem:Felder-Hasse-chain}, the first two pairs are Hasse-adjacent
in this finite poset, whereas $\lambda_1$ and $\lambda_3$ are not.
Section~4.4 of \cite{BoeNakanoWiesner} identifies the simple
$\Ext^1$-quiver of this finite category with the doubled Hasse diagram.
The adjacent pairs therefore have one-dimensional $\Ext^1$ in either
orientation, while the endpoints have zero $\Ext^1$ in either
orientation.  Lemma~\ref{lem:BNW-finite-reduction} transfers these values
back to the full classical block, and then
Lemma~\ref{lem:local-category-O} transfers them to $\mathscr C$.
\end{proof}

Let $A(\tau)$ and $B(\tau)$ denote the middle terms of nonzero classes
in
\[
 \Ext^1_{\mathscr C}(S_1,S_2)
 \quad\text{and}\quad
 \Ext^1_{\mathscr C}(S_2,S_3),
\]
respectively.  By Lemma~\ref{lem:one-dimensional-ext} these are
well-defined up to isomorphism, and
\begin{align}
 0&\longrightarrow S_2\longrightarrow A(\tau)
   \longrightarrow S_1\longrightarrow0,
   \label{eq:Nakano-tower-A}\\
 0&\longrightarrow S_3\longrightarrow B(\tau)
   \longrightarrow S_2\longrightarrow0.
   \label{eq:Nakano-tower-B}
\end{align}
These isomorphism classes agree with the modules denoted
$L(h_{\alpha_1},h_{\alpha_2})$ and
$L(h_{\alpha_2},h_{\alpha_3})$ in
\cite[Definition~5.9]{Nakano}.  Indeed, Nakano's two middle terms have
length two with the same Felder simple factors, so they satisfy the
weight-space finiteness conditions above and belong to $\mathscr C$.
They represent nonzero classes in the corresponding one-dimensional
Yoneda spaces of Lemma~\ref{lem:local-simple-Ext};
Lemma~\ref{lem:one-dimensional-ext} therefore identifies their middle
terms with $A(\tau)$ and $B(\tau)$, respectively.

Applying Lemma~\ref{lem:prelim-Hom-Yoneda} to
\eqref{eq:Nakano-tower-B} with target $S_1$, and using
Lemma~\ref{lem:local-simple-Ext}, gives
\begin{equation}\label{eq:Nakano-first-Ext}
 \Ext^1_{\mathscr C}\bigl(B(\tau),S_1\bigr)\cong\C.
\end{equation}
Indeed, the three preceding Hom spaces vanish,
$\Ext^1_{\mathscr C}(S_2,S_1)\cong\C$, and
$\Ext^1_{\mathscr C}(S_3,S_1)=0$.  We denote by $K(\tau)$ the middle
term of a nonzero class, so that
\begin{equation}\label{eq:Nakano-tower-K}
 0\longrightarrow S_1\longrightarrow K(\tau)
 \longrightarrow B(\tau)\longrightarrow0.
\end{equation}
This is also the Virasoro module isomorphism class denoted $K(\tau)$ in
\cite[Definition~5.10(2)]{Nakano}.  Nakano's module has length three,
hence lies in $\mathscr C$, and represents a nonzero class in the
one-dimensional space \eqref{eq:Nakano-first-Ext};
Lemma~\ref{lem:one-dimensional-ext} gives the identification.  No
identification of the two ambient Ext categories is involved.

We shall need two local consequences of the staggered-module
calculation.  We first record the precise highest-weight structure needed
to place the relevant extensions in that framework.

\begin{lem}[Highest-weight realization of the Felder extensions]
\label{lem:local-highest-weight-realization}
The modules $A(\tau)$ and $B(\tau)$ are ordinary highest-weight
Virasoro modules.  More generally, if
\[
 0\to T\to M\to S\to0
\]
is a nonsplit extension of simple lowest-weight modules with
$h_S<h_T$ and semisimple $L_0$, then $M$ is generated by any lift of a
lowest-weight vector of $S$.
\end{lem}

\begin{proof}
Let $v$ be a lowest-weight generator of $S$, choose a weight-vector lift
$\widetilde v\in M$, and set $U=U(\Vir)\widetilde v$.  The map
$U\to S$ is surjective.  Since $T$ is simple, either $U\cap T=0$ or
$U\cap T=T$.  The first alternative gives $U\cong S$ and hence a
section of $M\twoheadrightarrow S$, contrary to nonsplitting.  Thus
$T\subset U$, so $U=M$.  Positive Virasoro modes annihilate
$\widetilde v$ because they would have weights strictly below $h_S$.
Therefore $M$ is highest-weight.  Applying this to
\eqref{eq:Nakano-tower-A} and \eqref{eq:Nakano-tower-B}, whose
ordinariness follows from Lemma~\ref{lem:local-category-O}, proves the
first assertion.
\end{proof}

\begin{lem}[Staggered realization]
\label{lem:local-staggered-realization}
Let
\[
 0\longrightarrow A(\tau)\longrightarrow E\longrightarrow
 B(\tau)\longrightarrow0
\]
be logarithmic in $\mathscr C$.  Then this sequence is a rank-two
staggered Virasoro module in the sense of
\cite[Section~3]{KytolaRidout}, with left module $A(\tau)$ and right
module $B(\tau)$.  If $v_{\alpha_1}$ is a lowest-weight generator of
$A(\tau)$, there is a Shapovalov element
\[
 \eta_{12}:=\eta(h_{\alpha_1},h_{\alpha_2})
\]
for which the distinguished simple submodule is
\[
 S_2=U(\Vir)\eta_{12}v_{\alpha_1}\subset A(\tau),
 \qquad
 \deg\eta_{12}=h_{\alpha_2}-h_{\alpha_1}.
\]
This is the singular-vector datum used in the logarithmic coupling of
\cite[Theorem~5.11]{Nakano}.
\end{lem}

\begin{proof}
By Lemma~\ref{lem:local-highest-weight-realization}, both end terms are
ordinary highest-weight Virasoro modules.  Let $N_E$ be the canonical
nilpotent part of $L_0$.  Semisimplicity on the endpoints gives
\[
 N_E(A(\tau))=0,\qquad
 N_E(E)\subset A(\tau),\qquad N_E^2=0,
\]
and logarithmicity gives $N_E\ne0$.

We also verify indecomposability, which is part of the definition in
\cite[Section~3]{KytolaRidout}.  Since $A(\tau)$ is generated by its
one-dimensional lowest-weight space of weight $h_{\alpha_1}$, whereas
all weights of $B(\tau)$ are at least $h_{\alpha_2}>h_{\alpha_1}$,
\[
 \Hom\bigl(A(\tau),B(\tau)\bigr)=0.
\]
Both endpoint modules are cyclic highest-weight modules with
one-dimensional lowest-weight spaces, so their endomorphism algebras are
$\C$.  Let $e\in\End(E)$ be an idempotent, and write
$\iota:A(\tau)\hookrightarrow E$ and
$\pi:E\twoheadrightarrow B(\tau)$ for the structure maps.  Since
$\pi e\iota=0$, the endomorphism $e$ preserves $A(\tau)$ and induces
idempotents
\[
 e_A\in\End(A(\tau)),\qquad e_B\in\End(B(\tau)),
\]
so $e_A,e_B\in\{0,\id\}$.  If both are zero, then
$e(E)\subset A(\tau)$ and $e|_{A(\tau)}=0$, hence $e^2=0$ and therefore
$e=0$.  If both are the identity, the same argument applied to $\id-e$
gives $e=\id$.  In the mixed case $e_A=0$, $e_B=\id$, one has
$e\iota=0$ and $\pi e=\pi$; consequently
$\pi|_{e(E)}:e(E)\to B(\tau)$ is an isomorphism and supplies a section of
$\pi$.  The other mixed case is reduced to this one by replacing $e$
with $\id-e$.  Thus a nontrivial idempotent would split the extension.
But a split extension of the ordinary modules $A(\tau)$ and $B(\tau)$
has semisimple $L_0$, contrary to the logarithmicity of $E$.  Hence $E$
is indecomposable.  Together with $N_E^2=0$ and $N_E\ne0$, this places
$E$ in the rank-two staggered class of \cite{KytolaRidout}.

The nonsplit highest-weight extension
\eqref{eq:Nakano-tower-A} has unique simple submodule $S_2$, so the
singular vector generating it from the lowest-weight vector of
$A(\tau)$ is well defined up to nonzero scalar.  We normalize the
corresponding Shapovalov element as $\eta_{12}$.  This notation agrees
with the element $\eta(h_{\alpha_1},h_{\alpha_2})$ fixed immediately
before \cite[Theorem~5.11]{Nakano}.  Notice that no assertion is made
here that $\eta_{12}$ is the first singular vector of the ambient Verma
module; this distinction is relevant at the vacuum edge.
\end{proof}

\begin{lem}[Splitting of the distinguished quotient forces zero coupling]
\label{lem:split-forces-zero-coupling}
Let
\[
 0\longrightarrow A(\tau)\longrightarrow E\longrightarrow B(\tau)
 \longrightarrow0
\]
be logarithmic in $\mathscr C$, and let $\beta_E$ be the logarithmic
coupling defined from the Shapovalov element $\eta_{12}$ of
Lemma~\ref{lem:local-staggered-realization}.  If the induced sequence
\begin{equation}\label{eq:distinguished-quotient-general}
 0\longrightarrow S_1\longrightarrow E/S_2\longrightarrow B(\tau)
 \longrightarrow0
\end{equation}
splits, then $\beta_E=0$.
\end{lem}

\begin{proof}
Put $h_i=h_{\alpha_i}$.  Choose the lowest-weight vector of the
$B(\tau)$-summand supplied by a splitting of
\eqref{eq:distinguished-quotient-general}, and lift it to a vector
$v_2\in E[h_2]$.  In the notation preceding
\cite[Theorem~5.11]{Nakano},
\[
 \sigma(\eta_{12})v_2=c_Ev_1,
 \qquad
 (L_0-h_2)v_2=d_E\eta_{12}v_1,
 \qquad
 \beta_E=c_E/d_E.
\]
The image of $v_2$ in $E/S_2$ is a lowest-weight vector in the split
$B(\tau)$-summand, so every positive Virasoro mode annihilates it.  Hence
$\sigma(\eta_{12})v_2\in S_2$.  This vector has weight $h_1$, whereas
$S_2$ has lowest weight $h_2>h_1$; therefore
$\sigma(\eta_{12})v_2=0$ and $c_E=0$.

Moreover $d_E\ne0$.  Indeed, the nilpotent part $N_E$ is a Virasoro
homomorphism from the cyclic right module $B(\tau)$ into $A(\tau)$; if
it vanished on a lowest-weight lift, it would vanish on all of $E$,
contrary to logarithmicity.  Thus $\beta_E=0$.
\end{proof}

\begin{lem}[The auxiliary length-three extension]
\label{lem:local-auxiliary-extension}
Assume $\Ext^1_{\mathscr C}(S_2,S_2)=0$.  Let
\[
 0\longrightarrow S_2\longrightarrow E_2\longrightarrow
 B(\tau)\longrightarrow0
\]
represent a nonzero class in $\mathscr C$.  Then $E_2$ is logarithmic
and contains $B(\tau)'$ as a submodule; equivalently,
\begin{equation}\label{eq:auxiliary-E2}
 0\longrightarrow B(\tau)'\longrightarrow E_2
 \longrightarrow S_2\longrightarrow0
\end{equation}
is exact.
\end{lem}

\begin{proof}
Pull back the displayed extension along the distinguished inclusion
$S_3\hookrightarrow B(\tau)$.  By the six-term sequence for
\eqref{eq:Nakano-tower-B}, the map
\[
 \Ext^1_{\mathscr C}(B(\tau),S_2)
 \longrightarrow
 \Ext^1_{\mathscr C}(S_3,S_2)
\]
is injective because the preceding group
$\Ext^1_{\mathscr C}(S_2,S_2)$ vanishes by hypothesis.  Hence the
pullback is nonzero.
Since $\Ext^1_{\mathscr C}(S_3,S_2)\cong\C$, its middle term is
$B(\tau)'$, the contragredient of \eqref{eq:Nakano-tower-B}.  Thus
$B(\tau)'\hookrightarrow E_2$, and quotienting gives
\[
 E_2/B(\tau)'\cong B(\tau)/S_3\cong S_2,
\]
which proves \eqref{eq:auxiliary-E2}.

It remains to show that $L_0$ is not semisimple on $E_2$.  Suppose to the
contrary that it is semisimple.  Let $v$ be a lowest-weight generator of
$B(\tau)$, of weight $h_{\alpha_2}$, and choose a weight-vector lift
$\widetilde v\in E_2$.  Both endpoints have no weights below
$h_{\alpha_2}$, so every positive Virasoro mode annihilates
$\widetilde v$.  Hence $U:=U(\Vir)\widetilde v$ is a highest-weight
submodule mapping surjectively onto $B(\tau)$.  Since the kernel $S_2$ is
simple, either $U\cap S_2=0$ or $U\cap S_2=S_2$.  The first alternative
splits the original extension.  In the second alternative the
lowest-weight line of $S_2$ would lie in $U$.  But the
$h_{\alpha_2}$-weight space of the highest-weight module $U$ is
$\C\widetilde v$, whereas the lowest-weight vector of the kernel is
linearly independent of $\widetilde v$ because the latter has nonzero
image in $B(\tau)$.  This is impossible.  Thus every nonzero class in
$\Ext^1_{\mathscr C}(B(\tau),S_2)$ has logarithmic middle term.
\end{proof}

\begin{lem}[Nakano logarithmic existence input]
\label{lem:Nakano-existence-input}
Let $\tau=(\alpha_1,\alpha_2,\alpha_3)\in\mathcal T_{p_+,p_-}$ satisfy
$h_{\alpha_1}<h_{\alpha_2}<h_{\alpha_3}$.  Nakano's construction
provides a logarithmic extension
\begin{equation}\label{eq:Nakano-existence-E}
 0\longrightarrow A(\tau)\longrightarrow E_\tau\longrightarrow
 B(\tau)\longrightarrow0
\end{equation}
whose middle term belongs to $\mathscr C$.
\end{lem}

\begin{proof}
Since the strict inequality excludes $\mathcal T^0_{p_+,p_-}$, the
Virasoro subquotient construction described in the paragraph immediately
preceding \cite[Theorem~5.11]{Nakano} supplies the logarithmic module
\eqref{eq:Nakano-existence-E}.  Its middle term has finite length, its
generalized weight spaces are finite dimensional and lower truncated in
Nakano's ambient category, and all its Jordan--H\"older factors are the
Felder simples fixed above.  Consequently it is an object of
$\Oc_{c_{p,q}}$ and, by the Serre property, of $\mathscr C$.  No use of
\cite[Proposition~5.12]{Nakano} or of the later definition of $P(\tau)$
is made at this stage.
\end{proof}

\begin{thm}[Yoneda consequence of the distinguished-quotient property]
\label{thm:Nakano-input}
Let $\tau=(\alpha_1,\alpha_2,\alpha_3)\in\mathcal T_{p_+,p_-}$
satisfy $h_{\alpha_1}<h_{\alpha_2}<h_{\alpha_3}$.  Assume
\begin{equation}\label{eq:self-Ext-S2-hypothesis}
 \Ext^1_{\mathscr C}(S_2,S_2)=0,
\end{equation}
and assume the following distinguished-quotient property:
\begin{equation}\label{eq:NS-tau}
 (\mathrm{NS}_\tau)\qquad
 \begin{gathered}
 \text{every logarithmic extension }0\to A(\tau)\to E\to B(\tau)\to0
 \text{ in }\mathscr C\\
 \text{induces a nonsplit extension }
 0\to S_1\to E/S_2\to B(\tau)\to0.
 \end{gathered}
\end{equation}
Then
\begin{equation}\label{eq:extracted-P-Ext}
 \Ext^1_{\mathscr C}\bigl(K(\tau),S_2\bigr)\cong\C.
\end{equation}
Consequently all nonzero classes in \eqref{eq:extracted-P-Ext} have
isomorphic middle terms; we denote this isomorphism class by $P(\tau)$.
The proof also identifies it with Nakano's module of the same notation.
\end{thm}

\begin{proof}
We prove the Ext dimension using only the hypothesis
\eqref{eq:NS-tau} and Yoneda exactness; no identification of $\mathscr C$
with the ambient category of \cite[Definition~3.14]{Nakano} is made.
Fix the logarithmic extension supplied by
Lemma~\ref{lem:Nakano-existence-input}:
\begin{equation}\label{eq:local-E1}
 0\longrightarrow A(\tau)\longrightarrow E_1\longrightarrow
 B(\tau)\longrightarrow0.
\end{equation}
Quotienting \eqref{eq:local-E1} by the distinguished
$S_2\subset A(\tau)$ gives an exact sequence
\[
 0\longrightarrow S_1\longrightarrow E_1/S_2
 \longrightarrow B(\tau)\longrightarrow0.
\]
By the hypothesis \eqref{eq:NS-tau}, the displayed extension is
nonsplit and therefore represents a nonzero class in
\[
 \Ext^1_{\mathscr C}(B(\tau),S_1)\cong\C.
\]
Lemma~\ref{lem:one-dimensional-ext} now gives
\[
 E_1/S_2\cong K(\tau).
\]
Consequently
\begin{equation}\label{eq:E1-over-K}
 0\longrightarrow S_2\longrightarrow E_1\longrightarrow
 K(\tau)\longrightarrow0.
\end{equation}

We first prove
\begin{equation}\label{eq:E1-Ext-vanishing}
 \Ext^1_{\mathscr C}(E_1,S_2)=0.
\end{equation}
Apply $\Hom(-,S_2)$ to \eqref{eq:Nakano-tower-A}.  The connecting map
\[
 \End(S_2)\longrightarrow\Ext^1_{\mathscr C}(S_1,S_2)
\]
sends $\id_{S_2}$ to the nonzero class of $A(\tau)$, and is therefore
an isomorphism between one-dimensional spaces.  Together with the hypothesis
\eqref{eq:self-Ext-S2-hypothesis}, this gives
\begin{equation}\label{eq:A-S2-Ext-zero}
 \Ext^1_{\mathscr C}(A(\tau),S_2)=0.
\end{equation}
If \eqref{eq:E1-Ext-vanishing} failed, the six-term sequence for
\eqref{eq:local-E1}, together with \eqref{eq:A-S2-Ext-zero}, would give
\[
 \Ext^1_{\mathscr C}(B(\tau),S_2)\ne0.
\]
Choose a nonzero extension $E_2$ in this space.  By
Lemma~\ref{lem:local-auxiliary-extension}, $E_2$ is logarithmic and
fits into \eqref{eq:auxiliary-E2}.

Because $B(\tau)'$ has composition factors $S_2,S_3$, one has
$\Hom(B(\tau)',S_1)=0$.  Applying $\Hom(-,S_1)$ to
\eqref{eq:auxiliary-E2} therefore gives an injection
\[
 \Ext^1_{\mathscr C}(S_2,S_1)
 \hookrightarrow
 \Ext^1_{\mathscr C}(E_2,S_1).
\]
Take the image of a nonzero class and represent it by
\[
 0\longrightarrow S_1\longrightarrow X\longrightarrow E_2
 \longrightarrow0.
\]
Equivalently, $X$ is the pullback of the nonzero extension
$A(\tau)'$ of $S_2$ by $S_1$ along $E_2\twoheadrightarrow S_2$.
Thus there is also an exact sequence
\[
 0\longrightarrow B(\tau)'\longrightarrow X
 \longrightarrow A(\tau)'\longrightarrow0.
\]
Restricted contragredience is exact on $\mathscr C$ by
Lemma~\ref{lem:Nakano-ambient}; set $E_3:=X'$.  Dualizing the last two
sequences gives
\[
 0\longrightarrow E_2'\longrightarrow E_3\longrightarrow S_1
 \longrightarrow0,
 \qquad
 0\longrightarrow A(\tau)\longrightarrow E_3\longrightarrow
 B(\tau)\longrightarrow0.
\]
The first sequence shows that $E_3$ is logarithmic.  We make the
Yoneda naturality in the next step explicit.  After dualizing, the class
of
\[
 0\to E_2'\to E_3\to S_1\to0
\]
is the pushout of the nonzero class
$[A(\tau)]\in\Ext^1_{\mathscr C}(S_1,S_2)$ along the inclusion
$S_2\hookrightarrow E_2'$.  Its image under
\[
 \Ext^1_{\mathscr C}(S_1,E_2')
 \longrightarrow
 \Ext^1_{\mathscr C}(S_1,B(\tau))
\]
is therefore the pushout along the composite
\[
 S_2\hookrightarrow E_2'\twoheadrightarrow B(\tau),
\]
which is zero.  By construction, the lower sequence is the pushout of
\eqref{eq:Nakano-tower-A} along $S_2\hookrightarrow E_2'$.  Therefore
the induced map $A(\tau)\to E_3$ restricts on the distinguished
$S_2\subset A(\tau)$ to the kernel inclusion $S_2\hookrightarrow E_2'$.
Thus the copy of $S_2$ used in the following quotient is exactly the
distinguished copy occurring in \eqref{eq:NS-tau}.  The first exact sequence above shows that
\begin{equation}\label{eq:E3-first-split}
 0\longrightarrow B(\tau)\longrightarrow E_3/S_2
 \longrightarrow S_1\longrightarrow0
\end{equation}
splits.  We must still check that the distinguished exact sequence with
the two endpoint roles reversed is split; the abstract isomorphism of the
middle term with a direct sum does not by itself imply this.

Put $M=E_3/S_2$, let $i_B:B(\tau)\hookrightarrow M$ be the inclusion in
\eqref{eq:E3-first-split}, and choose a section
$s:S_1\to M$ of its quotient map.  Thus
\[
 M=i_B(B(\tau))\oplus s(S_1).
\]
Let $j:S_1\hookrightarrow M$ and $q:M\twoheadrightarrow B(\tau)$ be the
kernel inclusion and quotient map in the distinguished sequence
\[
 0\longrightarrow S_1\xrightarrow{\ j\ }M
 \xrightarrow{\ q\ }B(\tau)\longrightarrow0.
\]
The composition factors of $B(\tau)$ are $S_2$ and $S_3$, so
$\Hom(S_1,B(\tau))=0$.  Hence the $i_B(B(\tau))$-component of $j$ is
zero and
\[
 j=s\circ a
 \qquad\text{for some }a\in\End(S_1)=\C.
\]
Since $j\ne0$, the scalar $a$ is invertible; therefore
$j(S_1)=s(S_1)$.  From $q\circ j=0$ it follows that $q\circ s=0$.
Consequently $q$ is determined by the endomorphism
$q\circ i_B:B(\tau)\to B(\tau)$.  This endomorphism is surjective because
$q$ is surjective and $q(s(S_1))=0$.  The module $B(\tau)$ has finite
length, so every surjective endomorphism of it is an automorphism.  Thus
\[
 i_B\circ(q\circ i_B)^{-1}:B(\tau)\longrightarrow M
\]
is a section of $q$.  The distinguished sequence is therefore split,
contradicting \eqref{eq:NS-tau}.  Thus \eqref{eq:E1-Ext-vanishing} holds.

It remains to compute the dimension in
\eqref{eq:extracted-P-Ext}.  From
\eqref{eq:Nakano-tower-B},
\[
 \Hom(B(\tau),S_2)\cong\End(S_2)\cong\C.
\]
Also $\Hom(A(\tau),S_2)=0$: a nonzero restriction to the submodule
$S_2$ would split \eqref{eq:Nakano-tower-A}, while a map vanishing on
that submodule would factor through
$\Hom(S_1,S_2)=0$.  Hence \eqref{eq:local-E1} gives
\[
 \Hom(E_1,S_2)\cong\Hom(B(\tau),S_2)\cong\C.
\]
Similarly \eqref{eq:Nakano-tower-K} gives
\[
 \Hom(K(\tau),S_2)\cong\Hom(B(\tau),S_2)\cong\C.
\]
The map
$\Hom(K(\tau),S_2)\to\Hom(E_1,S_2)$ induced by
\eqref{eq:E1-over-K} is nonzero and hence an isomorphism.  The
restriction
\[
 \Hom(E_1,S_2)\longrightarrow\End(S_2)
\]
is zero, since its generator factors through
$E_1\twoheadrightarrow B(\tau)\twoheadrightarrow S_2$ and therefore
vanishes on the distinguished $S_2\subset A(\tau)$.  The six-term
sequence for \eqref{eq:E1-over-K}, together with
\eqref{eq:E1-Ext-vanishing}, now gives an isomorphism
\[
 \End(S_2)\xrightarrow{\ \sim\ }
 \Ext^1_{\mathscr C}(K(\tau),S_2).
\]
This proves \eqref{eq:extracted-P-Ext}.  Before proving
\cite[Proposition~5.12]{Nakano}, Nakano constructs the nonzero extension
labelled (5.2),
\[
 0\longrightarrow S_2\longrightarrow P_{\mathrm N}(\tau)
 \longrightarrow K(\tau)\longrightarrow0.
\]
Its middle term has finite length, finite-dimensional lower-truncated
generalized weight spaces, and only the Felder composition factors fixed
above; hence it belongs to $\mathscr C$.  Since the Ext space just computed is
one-dimensional, the one-dimensional extension lemma identifies
$P_{\mathrm N}(\tau)$ with the isomorphism class denoted $P(\tau)$ here.  This recovers Nakano's notation without using
Proposition~5.12 in the proof of the Ext dimension.
\end{proof}

By definition of $P(\tau)$, we have the final exact sequence
\begin{equation}\label{eq:Nakano-tower-P}
 0\longrightarrow L(h_{\alpha_2})\longrightarrow P(\tau)
 \longrightarrow K(\tau)\longrightarrow0.
\end{equation}

\begin{rmk}\label{rmk:no-socle}
No description of the socle or contragredient self-duality of a general
$P(\tau)$ is used.  In particular, we do not use the dotted socle arrow
whose status is left open in the discussion around Figure~5.1 of
\cite{Nakano}.
\end{rmk}

\subsection{The Felder dictionary for the one-row chain}

Nakano writes $\alpha_{u,v;m}$ for the Fock momentum whose lowest
conformal weight is $h_{u,v;m}$, and $F_{\alpha}$ for the corresponding
Fock module.  The set $\mathcal T_{p_+,p_-}$ consists of triples of such
labels whose Fock modules are consecutive in one Felder complex, with
ordered lowest weights as in Definition~5.6 of \cite{Nakano};
the subset $\mathcal T^0_{p_+,p_-}$ consists of the chain-type triples
in $\mathcal T_{p_+,p_-}$ for which the first two weights coincide; see
\cite[Definition~5.7]{Nakano}.  In particular, a triple with strictly
increasing first two lowest weights lies outside
$\mathcal T^0_{p_+,p_-}$.  We use no other part of the free-field
construction.  Recall
from \cite[(2.3)--(2.4)]{Nakano} that
\begin{equation}\label{eq:Nakano-weight-shift}
 h_{u,v;m}=h_{u-mp_+,v}=h_{u,v+mp_-}.
\end{equation}
For $n\geq1$ and $1\leq r\leq p-1$, put
\begin{equation}\label{eq:abc}
 a=np-r,\qquad b=np+r,\qquad c=(n+2)p-r.
\end{equation}

\begin{lem}[Primitive Shapovalov coefficient]
\label{lem:primitive-Shapovalov-coefficient}
Let $R,S$ be positive integers with $1\le R<p_+$.  For the normalized
Shapovalov element at the Kac weight $h_{R,S}$, the linear coefficient
$R_{R,S}$ in \cite[Proposition~5.2]{Nakano} is nonzero.
\end{lem}

\begin{proof}
Nakano's coefficient is
\[
 R_{R,S}=2\!\!\prod_{\substack{
 1-R\le k\le R,\ 1-S\le l\le S\\
 (k,l)\ne(0,0),(R,S)}}
 \left(
 k\Bigl(\frac{p_+}{p_-}\Bigr)^{-1/2}
 +l\Bigl(\frac{p_+}{p_-}\Bigr)^{1/2}
 \right).
\]
If one of the displayed factors vanished, then
$kp_-+lp_+=0$.  Since $\gcd(p_+,p_-)=1$, there would be an
$m\in\Z$ with
\[
 (k,l)=m(p_+,-p_-).
\]
But $1-R\le k\le R$ and $R<p_+$ exclude every $m\ne0$; the case
$m=0$ is the omitted pair $(0,0)$.  Thus every factor is nonzero.
\end{proof}

\begin{lem}[Harish--Chandra reduction of the Shapovalov operator]
\label{lem:Shapovalov-HC-reduction}
Let $X\in U(\Vir_-)$ be homogeneous and let $\sigma(X)$ be its image
under the standard anti-involution.  There is a polynomial
$P_X(z,\kappa)\in\C[z,\kappa]$ such that
\begin{equation}\label{eq:HC-Shapovalov-congruence}
 \sigma(X)X\equiv P_X(L_0,C)
 \pmod{U(\Vir)\Vir_{>0}}.
\end{equation}
Consequently, on any generalized highest-weight space $W$ of central
charge $c$, namely a subspace annihilated by every $L_m$ with $m>0$,
\[
 \left.\sigma(X)X\right|_W=P_X(L_0,c).
\]
If $X=\mathfrak S_{R,S}$ is the normalized Shapovalov element at the
Kac weight $h_{R,S}$, then $P_X(h,c)$ is its Shapovalov norm on a Verma
highest-weight vector.  In the normalization of
\cite[Proposition~5.2]{Nakano},
\[
 P_X(h_{R,S},c_{p,q})=0,\qquad
 \partial_zP_X(h_{R,S},c_{p,q})=R_{R,S}.
\]
\end{lem}

\begin{proof}
The element $\sigma(X)X$ has degree zero.  Reorder it by PBW with all
negative modes to the left, then powers of $L_0$ and $C$, and positive
modes to the right.  Modulo the right ideal
$U(\Vir)\Vir_{>0}$ every term containing a positive mode vanishes.  A
degree-zero PBW monomial with no positive or negative modes lies in
$U(\C L_0\oplus\C C)$, which proves
\eqref{eq:HC-Shapovalov-congruence}.  The first displayed operator
identity follows immediately on a generalized highest-weight space.

Applying the congruence to the highest-weight line of the Verma module
shows that $P_X(h,c)$ is exactly the scalar Shapovalov norm of $Xv_h$.
For $X=\mathfrak S_{R,S}$, Proposition~5.2 of \cite{Nakano} gives its
first-order expansion at $h=h_{R,S}$, and hence the two asserted values
of $P_X$ and its derivative.
\end{proof}

\begin{lem}[One-row self-extension vanishing]
\label{lem:one-row-self-Ext}
For $a,b,c$ in \eqref{eq:abc},
\[
 \Ext^1_{\mathscr C}(S_a,S_a)=
 \Ext^1_{\mathscr C}(S_b,S_b)=
 \Ext^1_{\mathscr C}(S_c,S_c)=0.
\]
\end{lem}

\begin{proof}
We first choose positive Kac representatives to which
Lemma~\ref{lem:primitive-Shapovalov-coefficient} applies.  If $p<q$, so
$(p_+,p_-)=(p,q)$, the weight-shift identity
\eqref{eq:Nakano-weight-shift} gives
\[
 \begin{aligned}
 h_a&=h_{\,r,\,q-1+(n-1)q},\\
 h_b&=h_{\,p-r,\,q-1+nq},\\
 h_c&=h_{\,r,\,q-1+(n+1)q}.
 \end{aligned}
\]
Their first coordinates lie in $\{1,\ldots,p_+-1\}$.  If $p>q$, so
$(p_+,p_-)=(q,p)$, the reflected positive representatives used below in
Proposition~\ref{prop:Felder-labels} give
\[
 \begin{aligned}
 h_a&=h_{\,1,\,p-r+(n-1)p},\\
 h_b&=h_{\,1,\,r+np},\\
 h_c&=h_{\,1,\,p-r+(n+1)p},
 \end{aligned}
\]
whose first coordinate is $1<p_+=q$.  In either ordering, each of the
three weights therefore has a positive representative $(R,S)$ with
$1\le R<p_+$, and Lemma~\ref{lem:primitive-Shapovalov-coefficient}
gives $R_{R,S}\ne0$.

Let $S=L(h_{R,S})$ be any of these three simples and suppose
\[
 0\longrightarrow S\xrightarrow{\iota}E\xrightarrow{\pi}S
 \longrightarrow0
\]
is a logarithmic self-extension.  Let $N_E$ be the canonical nilpotent
part of $L_0$.  Since both endpoint copies are ordinary,
\[
 N_E=\iota\circ(c\,\id_S)\circ\pi
 \qquad(c\in\C^\times).
\]
For a lowest-weight vector $v_h$, choose $u_0$ with $\pi(u_0)=v_h$ and
put $u_1=\iota(v_h)$.  Then
\[
 (L_0-h)u_0=cu_1.
\]
The two ordinary endpoints imply
$\operatorname{Spec}(E)\subset h+\Z_{\ge0}$.  Thus every positive
Virasoro mode annihilates the lowest generalized-weight space, because
$L_mE_h^{\rm gen}\subset E_{h-m}^{\rm gen}=0$ for $m>0$.  Put
$X=\mathfrak S_{R,S}$.  Lemma~\ref{lem:Shapovalov-HC-reduction} gives
on this generalized highest-weight space
\[
 \sigma(X)X=P_X(L_0,c_{p,q}).
\]
Here
\[
 P_X(h,c_{p,q})=0,\qquad
 \partial_zP_X(h,c_{p,q})=R_{R,S}\ne0.
\]
Since $L_0=h\,\id+N_E$ on $E_h^{\rm gen}$ and $N_E^2=0$, polynomial
functional calculus gives
\[
 P_X(L_0,c_{p,q})
 =R_{R,S}N_E.
\]
Therefore
\[
 \sigma(X)Xu_0=cR_{R,S}u_1\ne0.
\]
On the other hand,
$\pi(Xu_0)=Xv_h=0$, so $Xu_0$ lies in the simple kernel copy $S$.
Consequently $\sigma(X)Xu_0$ also belongs to $S$; it has conformal
weight $h$, and is therefore a scalar multiple of $u_1$.  Normalize the
nondegenerate contravariant form on $S$ by
$\langle u_1,u_1\rangle=1$.  Since $Xu_1=0$, contravariance gives
\[
 \big\langle u_1,\sigma(X)Xu_0\big\rangle
 =\big\langle Xu_1,Xu_0\big\rangle=0.
\]
It follows that $\sigma(X)Xu_0=0$, contradicting
$\sigma(X)Xu_0=cR_{R,S}u_1\ne0$.  Therefore no logarithmic
self-extension exists.

If $L_0$ is semisimple on a self-extension, the finite-generation and
$\Vir_{>0}$-local-finiteness argument of
Lemma~\ref{lem:local-category-O} places it in the classical Virasoro
category $\mathcal O$.  Lemma~\ref{lem:BNW-finite-reduction} reduces its
self-$\Ext^1$ to a finite BNW model.  The doubled-Hasse description of
Section~4.4 of \cite{BoeNakanoWiesner} has no loops, so the self-extension
entry is zero.  Thus ordinary self-extensions split as well, proving the
claim.
\end{proof}

\begin{prop}[Felder-label matching]\label{prop:Felder-labels}
There is an explicitly labelled triple
$\tau_{n,r}\in
\mathcal T_{p_+,p_-}\setminus\mathcal T^0_{p_+,p_-}$
whose ordered conformal weights are $(h_a,h_b,h_c)$.  More precisely, if
$p<q$, then
\begin{equation}\label{eq:tau-p-less-q}
 \tau_{n,r}=
 \bigl(\alpha_{r,q-1;n-1},
       \alpha_{p-r,q-1;n},
       \alpha_{r,q-1;n+1}\bigr),
\end{equation}
whereas if $p>q$, then
\begin{equation}\label{eq:tau-p-greater-q}
 \tau_{n,r}=
 \bigl(\alpha_{q-1,r;-n+1},
       \alpha_{q-1,p-r;-n},
       \alpha_{q-1,r;-n-1}\bigr).
\end{equation}
In either case
\begin{equation}\label{eq:tau-nr}
 (h_{\alpha_1},h_{\alpha_2},h_{\alpha_3})=(h_a,h_b,h_c)
\end{equation}
and
\begin{equation}\label{eq:A-B-tau}
 A(\tau_{n,r})\cong K_a,
 \qquad
 B(\tau_{n,r})\cong K_b.
\end{equation}
\end{prop}

\begin{proof}
Since $p,q\ge2$ are coprime, $p\neq q$.  To avoid confusing the one-row
label $r$ with Nakano's Felder coordinates, we denote the latter by
$R,S$ throughout this proof.  Suppose first that $p<q$, so
$(p_+,p_-)=(p,q)$.  Here $1\le r\le p-1=p_+-1$ and $q-1=p_--1$, so the displayed Felder
labels are admissible.  Apply \cite[Proposition~2.7(1)]{Nakano} with
\[
 (R,S,N)=(p-r,q-1,n).
\]
Then $R^\vee=p_+-R=r$, and the consecutive terms
$F_{R^\vee,S;N-1}\to F_{R,S;N}\to F_{R^\vee,S;N+1}$ are exactly the
three Fock modules in \eqref{eq:tau-p-less-q}.  The label conversion is
\[
\begin{array}{c|c|c}
 \text{Nakano label}&\text{weight}&\text{one-row label}\\ \hline
 \alpha_{r,q-1;n-1}&h_{r,q-1;n-1}&h_{np-r,1}=h_a\\
 \alpha_{p-r,q-1;n}&h_{p-r,q-1;n}&h_{np+r,1}=h_b\\
 \alpha_{r,q-1;n+1}&h_{r,q-1;n+1}&h_{(n+2)p-r,1}=h_c
\end{array}
\]
by \eqref{eq:Nakano-weight-shift}.  The three numerator identities are
\[
\begin{aligned}
 q\bigl(r-(n-1)p\bigr)-p(q-1)
   &=-\bigl(q(np-r)-p\bigr),\\
 q\bigl(p-r-np\bigr)-p(q-1)
   &=-\bigl(q(np+r)-p\bigr),\\
 q\bigl(r-(n+1)p\bigr)-p(q-1)
   &=-\bigl(q((n+2)p-r)-p\bigr).
\end{aligned}
\]
Thus the squares of the three Kac numerators are exactly those defining
$h_a,h_b,h_c$, respectively.

If $p>q$, then $(p_+,p_-)=(q,p)$.  Here
$q-1=p_+-1$ and $1\le r,p-r\le p-1=p_--1$.  Apply
\cite[Proposition~2.7(2)]{Nakano} with
\[
 (R,S,N)=(q-1,p-r,-n).
\]
Proposition~2.7(2) is stated with the Felder index ranging over
$\Z$, so the choice $N=-n$ is within its stated range.  Moreover,
Nakano's set $\mathcal A_{p_+,p_-}$ in Section~5.1 contains
$\alpha_{R,S;N}$ for arbitrary integral $R,S,N$, and
\cite[Definition~5.6]{Nakano} defines $\mathcal T_{p_+,p_-}$ by the
actual adjacency of the three Fock modules in a Felder complex; it does
not require replacing negative Felder indices by nonnegative ones.  Thus
the original labels in \eqref{eq:tau-p-greater-q} are legitimate
$\mathcal A_{p_+,p_-}$-labels.  Here $S^\vee=p_- -S=r$, and the consecutive
terms $F_{R,S^\vee;N+1}\to F_{R,S;N}\to F_{R,S^\vee;N-1}$ are exactly the
three Fock modules in \eqref{eq:tau-p-greater-q}.  Using $c_{q,p}=c_{p,q}$ and
$h_{j,1}^{(p,q)}=h_{1,j}^{(q,p)}$, one obtains
\[
\begin{array}{c|c|c}
 \text{Nakano label}&\text{weight for }(q,p)&\text{one-row label for }(p,q)\\ \hline
 \alpha_{q-1,r;-n+1}&h_{q-1,r;-n+1}&h_{np-r,1}=h_a\\
 \alpha_{q-1,p-r;-n}&h_{q-1,p-r;-n}&h_{np+r,1}=h_b\\
 \alpha_{q-1,r;-n-1}&h_{q-1,r;-n-1}&h_{(n+2)p-r,1}=h_c
\end{array}
\]
(the superscripts $(q,p)$ and $(p,q)$ are understood in the middle and
right columns, respectively).  To place these negative Felder indices in
the standard range used in Nakano's Proposition~5.5, we use the immediate
symmetry of the weight formula
\begin{equation}\label{eq:Nakano-negative-index}
 h_{u,v;-m}=h_{p_+-u,p_--v;m}.
\end{equation}
Indeed, for arbitrary $m\in\Z$,
\[
 p_-\bigl(u+mp_+\bigr)-p_+v
 =-\Bigl(
 p_-\bigl(p_+-u-mp_+\bigr)-p_+\bigl(p_--v\bigr)
 \Bigr),
\]
so the two Kac numerators differ by an overall sign and therefore have
the same square.  Thus, for $(p_+,p_-)=(q,p)$,
\[
\begin{aligned}
 h_{q-1,r;-n+1}&=h_{1,p-r;n-1},\\
 h_{q-1,p-r;-n}&=h_{1,r;n},\\
 h_{q-1,r;-n-1}&=h_{1,p-r;n+1}.
\end{aligned}
\]
The reflected coordinates $1$, $r$, and $p-r$ lie in the required
fundamental ranges $1\le1\le p_+-1$ and
$1\le r,p-r\le p_--1$.

We emphasize that \eqref{eq:Nakano-negative-index} is used only to
identify lowest conformal weights, and hence the corresponding simple
Virasoro modules, with representatives in the parameter ranges of
Proposition~5.5.  It is not an identification of the underlying Fock
modules.  Adjacency of the three Fock modules in the Felder complex was
established directly from Proposition~2.7(2) at the original, possibly
negative, Felder indices.  Since irreducible lowest-weight Virasoro modules
at fixed central charge are determined by their lowest weights, the
weight equalities above identify the corresponding simple Virasoro
modules.  Yoneda $\Ext^1$ is functorial under isomorphisms of either
endpoint, so the Ext dimensions obtained from positive-index
representatives transport to the original labels with the same
quotient--submodule orientation.

By \eqref{eq:weight-difference},
\[
 h_b-h_a=r(nq-1)>0,
 \qquad
 h_c-h_b=(p-r)((n+1)q-1)>0.
\]
The three labels therefore satisfy the two conditions used in
\cite[Definition~5.6]{Nakano}: their Fock modules occur consecutively in
the indicated Felder complex, and their lowest weights satisfy
\[
 h_a<h_b<h_c.
\]
Hence $\tau_{n,r}\in\mathcal T_{p_+,p_-}$.  The strict first inequality
excludes the equality condition defining $\mathcal T^0_{p_+,p_-}$, so
$\tau_{n,r}\notin\mathcal T^0_{p_+,p_-}$.  The three
lowest weights therefore lie in Nakano's Felder set $H_{p_+,p_-}$.
Since $S_a,S_b,S_c$ are finite-length grading-restricted simple objects in
$\Oc_{c_{p,q}}$, they belong to
$\mathscr C$:
\[
 S_a,S_b,S_c\in\mathscr C.
\]
The Kac sequences
\[
 0\longrightarrow S_b\longrightarrow K_a\longrightarrow S_a
 \longrightarrow0,
 \qquad
 0\longrightarrow S_c\longrightarrow K_b\longrightarrow S_b
 \longrightarrow0
\]
are non-split by \eqref{eq:Kac-length-two}.  Lemma~\ref{lem:Nakano-Serre}
therefore gives
\[
 K_a,K_b\in\mathscr C.
\]
No identification of $K_a$ with $A(\tau_{n,r})$, or of
$K_b$ with $B(\tau_{n,r})$, has been used in establishing this
membership.  By Lemma~\ref{lem:Nakano-exact-compatibility}, the two
nonsplit Kac sequences represent nonzero classes in the corresponding
Yoneda groups of $\mathscr C$.  By
Lemma~\ref{lem:local-simple-Ext}, the relevant Yoneda groups in
$\mathscr C$ have the same adjacency dimensions as in the classical
highest-weight calculation.  To identify the entries occurring here,
we use the Felder-label table in \cite[Proposition~5.5]{Nakano}.
We spell out the four parameter ranges.

If $p<q$, the first pair is
\[
 h_a=h_{r,q-1;n-1},\qquad h_b=h_{p-r,q-1;n}.
\]
For $n=1$ this is Proposition~5.5(1), and for $n\ge2$ it is
Proposition~5.5(2), with reflection in the first coordinate.  The second
pair
\[
 h_b=h_{p-r,q-1;n},\qquad h_c=h_{r,q-1;n+1}
\]
is Proposition~5.5(2) for every $n\ge1$.

Suppose now that $p>q$.  For the first pair with $n=1$ we do not convert
the source weight.  Instead
\[
 h_a=h_{q-1,r;0},\qquad
 h_b=h_{q-1,p-r;-1}=h_{1,r;1}
\]
by \eqref{eq:Nakano-negative-index}.  Thus Proposition~5.5(1) applies
directly with $R=q-1$, $S=r$, and $N=0$, for which $R^\vee=1$.  For
$n\ge2$ we use the positive-index forms
\[
 h_a=h_{1,p-r;n-1},\qquad h_b=h_{1,r;n},
\]
and Proposition~5.5(2), with reflection in the second coordinate.  The
second pair is, for every $n\ge1$,
\[
 h_b=h_{1,r;n},\qquad h_c=h_{1,p-r;n+1},
\]
again the corresponding second-coordinate case of Proposition~5.5(2).
Hence
\[
 \Ext^1_{\mathscr C}(S_a,S_b)\cong\C,
 \qquad
 \Ext^1_{\mathscr C}(S_b,S_c)\cong\C,
\]
with the orientation \eqref{eq:Ext-convention}.  The strict inequalities
$h_a<h_b<h_c$ also show that these three simple modules are pairwise
nonisomorphic.  Thus the order of the three terms agrees with Nakano's
ordered triple $(\alpha_1,\alpha_2,\alpha_3)$, not merely with the
direction of the corresponding Felder differential.  The two Kac
sequences are therefore nonzero classes in precisely these one-row Ext
groups, and
Lemma~\ref{lem:one-dimensional-ext} identifies their middle terms with the modules of
\cite[Definition~5.9]{Nakano}:
$K_a\cong A(\tau_{n,r})$ and $K_b\cong B(\tau_{n,r})$.
\end{proof}

\begin{lem}[One-row Serre-category membership]\label{lem:Nakano-membership}
Fix $n\geq1$ and $1\leq r\leq p-1$, and let $a,b,c$ be as in
\eqref{eq:abc}.  Then
\[
 S_a,S_b,S_c,K_a,K_b\in\mathscr C.
\]
If $0\to K_a\to E\to K_b\to0$ is a short exact sequence in
$\Oc_{c_{p,q}}$, then
$E$, the distinguished quotient $E/S_b$, and all three terms of the
induced sequence $0\to S_b\to E\to E/S_b\to0$ belong to
$\mathscr C$, where $S_b\subset E$ denotes the
image of the distinguished composite
\[
 S_b\hookrightarrow K_a\hookrightarrow E.
\]
\end{lem}

\begin{proof}
Proposition~\ref{prop:Felder-labels} realizes $h_a,h_b,h_c$ as consecutive
Felder weights.  Hence $S_a,S_b,S_c$ are objects of $\mathscr C$.  The Kac sequences and
Lemma~\ref{lem:Nakano-Serre} place $K_a,K_b$ in the same finite-length
subcategory.  The remaining assertion follows from its Serre closure.
\end{proof}


\begin{lem}[Nakano scope for the non-vacuum one-row case]
\label{lem:one-row-Nakano-scope}
Let $n\ge1$, $1\le r\le p-1$, and put
\[
 a=np-r,\qquad b=np+r.
\]
Assume $a\ge2$.  If
\begin{equation}\label{eq:one-row-scope-extension}
 0\longrightarrow K_a\longrightarrow E\longrightarrow K_b
 \longrightarrow0
\end{equation}
is exact in $\Oc_{c_{p,q}}$, then $E$ belongs to Nakano's category
$L_{c_{p_+,p_-}}\text{-mod}$ of \cite[Definition~5.1]{Nakano}.
\end{lem}

\begin{proof}
We verify the conditions of \cite[Definitions~3.14 and~5.1]{Nakano}
one by one.  The central element acts by $c_{p_+,p_-}=c_{p,q}$ on every
term, so Definition~3.14(1) is automatic.

For Definition~3.14(2), the one-row Kac character is
\[
 \ch_zK_a=z^{h_a}
 \frac{1-z^a}{\prod_{m\ge1}(1-z^m)}.
\]
Let $\mathfrak p(N)$ be the partition number, with
$\mathfrak p(N)=0$ for $N<0$.  The relative degree-$N$ subspace of
$K_a$ has dimension
\[
 \mathfrak p(N)-\mathfrak p(N-a).
\]
When $a\ge2$ this number is strictly positive for every $N\ge0$: for
$N<a$ this is immediate, while for $N\ge a$ adjoining a part of size
$a$ injects the partitions of $N-a$ into those of $N$, and the partition
$1^N$ is not in the image.  Hence
\begin{equation}\label{eq:Ka-full-support}
 \operatorname{Spec}K_a=h_a+\Z_{\ge0}
\end{equation}
with no missing weights.  Since
\[
 h_b-h_a=r(nq-1)\in\Z_{>0},
\]
both endpoint Kac modules have weights in $h_a+\Z_{\ge0}$.  Exactness on
generalized weight spaces in \eqref{eq:one-row-scope-extension}, together
with $K_a\hookrightarrow E$, gives the equality
\[
 \operatorname{Spec}E=h_a+\Z_{\ge0},
 \qquad
 0<\dim E[h_a+N]<\infty\quad(N\ge0).
\]
Thus Definition~3.14(2) holds with the single-element set
$H_0(E)=\{h_a\}$; in particular this is the required equality of support,
not merely a containment.

For Definition~3.14(3), Lemma~\ref{lem:Nakano-ambient} constructs the
restricted contragredient and shows that it has the same support and the
same finite generalized-weight multiplicities.  Dualizing
\eqref{eq:one-row-scope-extension} gives a finite-length, lower-truncated
Virasoro module, so the contragredient is again an object of the ambient
category of Definition~3.14.

It remains to check Definition~5.1.  The extension $E$ has finite length,
hence finite socle series, and every Jordan--H\"older factor is one of the
Felder simples occurring in $K_a$ or $K_b$.  Proposition~\ref{prop:Felder-labels}
places all of their lowest weights in $H_{p_+,p_-}$.  These are precisely
the two additional requirements in Definition~5.1.  Therefore this
particular middle term satisfies every defining condition of
$L_{c_{p_+,p_-}}\text{-mod}$.
\end{proof}

\begin{lem}[Kyt\"ol\"a--Ridout higher-prime obstruction]
\label{lem:higher-prime-determinant}
Assume $p>q$ and set
\[
 h^R=h_{2p-1,1}=h_{1,2q-1},\qquad
 \ell^-_1=2q-1,\qquad \ell^+_1=2p-1.
\]
Let $X^-_1v_{h^R}$ and $X^+_1v_{h^R}$ be the first prime singular
vectors of the braid-type Verma module $M(c_{p,q},h^R)$ at these
grades.  Suppose that a rank-one right-Verma staggered module has all
ordinary beta-invariants equal to zero and that the higher prime
$X^+_1v_{h^R}$ has the singular lift required for descent to the
quotient by $U(\Vir)X^+_1v_{h^R}$.  Then
\begin{equation}\label{eq:higher-prime-det-lower-bound}
 \operatorname{ord}_{h=h^R}\det B_{\ell^+_1}(h)
 \ge
 \mathfrak p(\ell^+_1-\ell^-_1)+2,
\end{equation}
where $B_N(h)$ is the Shapovalov Gram matrix at relative grade $N$ and
$\mathfrak p$ is the partition function.
\end{lem}

\begin{proof}
This is precisely the higher-first-prime refinement in the paragraph
immediately following \cite[Proposition~7.5]{KytolaRidout}.  In the
braid-type case with quotient-defining prime
$\overline X=X^+_1$, Kyt\"ol\"a--Ridout introduce the auxiliary
obstruction $\overline\beta$ for the required singular lift and show,
by the determinant argument following equations (7.29)--(7.31), that
when the ordinary beta-invariants vanish, the condition
$\overline\beta=0$ forces the Kac determinant at grade $\ell^+_1$ to
have a zero of order at least
\[
 \mathfrak p(\ell^+_1-\ell^-_1)+2.
\]
For a staggered module descending through
$U(\Vir)X^+_1v_{h^R}$, Proposition~7.2 of loc.~cit. supplies precisely
that singular lift, so the corresponding auxiliary obstruction
vanishes.  This gives \eqref{eq:higher-prime-det-lower-bound}.
\end{proof}

\begin{lem}[Vacuum-edge distinguished quotient]
\label{lem:vacuum-edge-staggered}
Let $n=1$ and $r=p-1$, so that
\[
 a=1,\qquad b=2p-1,
\]
and let
\[
 0\longrightarrow K_1\longrightarrow E\longrightarrow K_{2p-1}
 \longrightarrow0
\]
be logarithmic in $\Oc_{c_{p,q}}$.  Then the induced sequence
\begin{equation}\label{eq:vacuum-edge-quotient}
 0\longrightarrow S_1\longrightarrow E/S_{2p-1}
 \longrightarrow K_{2p-1}\longrightarrow0
\end{equation}
is nonsplit.
\end{lem}

\begin{proof}
This is the unique one-row case in which the support argument of
Lemma~\ref{lem:one-row-Nakano-scope} fails: the coefficient of relative
degree $1$ in $(1-z)/\prod_{m\ge1}(1-z^m)$ is zero.  By
Lemma~\ref{lem:local-staggered-realization}, $E$ is a rank-two staggered
module with distinguished Shapovalov datum
$\eta_{12}=\eta(h_1,h_{2p-1})$ and logarithmic coupling $\beta_E$.
We next identify the rank of the distinguished singular datum in the
sense of \cite{KytolaRidout}.  The vacuum weight satisfies
\[
 h_1=0=h_{1,1}=h_{p-1,q-1}.
\]
By the Feigin--Fuchs description of the vacuum Verma module, its two
first prime singular vectors occur at relative grades
\[
 1,\qquad (p-1)(q-1).
\]
The one-row vacuum Kac module is
\[
 K_1=M(c_{p,q},0)\big/U(\Vir)L_{-1}v,
\]
so the grade-one prime is killed.  On the other hand,
\[
 h_{2p-1}-h_1=(p-1)(q-1),
\]
and the image of the second first-prime singular vector is nonzero in
$K_1$ and generates its unique proper simple submodule
$S_{2p-1}$; this is also the $n=0,r=1$ case of
\eqref{eq:Kac-length-two}.  Hence $\eta_{12}v_1$ is, up to nonzero
scalar, the image of that first-prime vector.

We also verify primeness after passage to the quotient $K_1$.  The
nonsplit Kac sequence
\[
 0\longrightarrow S_{2p-1}\longrightarrow K_1
 \longrightarrow S_1\longrightarrow0
\]
shows that $K_1$ has length two and that
$U(\Vir)\eta_{12}v_1=S_{2p-1}$ is its unique nonzero proper
highest-weight submodule.  If $\eta_{12}v_1$ admitted a nontrivial
factorization through another proper singular vector of $K_1$, the
submodule generated by that intermediate singular vector would lie
strictly between $S_{2p-1}$ and $K_1$, which is impossible.  Thus
$\eta_{12}v_1$ is prime in $K_1$.  The associated vector
$\omega_0=(L_0-h_{2p-1})v_2$ consequently has rank one in the sense of
\cite[Section~2]{KytolaRidout}.  The
right-Verma cover therefore lies in case~(1) or~(1') of
\cite[Sections~6.4--6.5]{KytolaRidout}.  By
\cite[Theorems~6.14--6.15]{KytolaRidout}, its space of staggered data is
one-dimensional and is parametrized by a single ordinary beta-invariant.
Proposition~4.6 of loc.~cit. preserves the staggered data when the right
module is replaced by a highest-weight cover, and Section~7 observes that
the ordinary beta-invariant is unchanged under this replacement.

There is a harmless normalization issue which we make explicit.  The
Shapovalov element $\eta_{12}$ used in Nakano's coupling and the normalized
rank-one singular operator used in \cite{KytolaRidout} differ by a
nonzero scalar.  Accordingly, the two ordinary couplings may differ by a
nonzero scalar, but their vanishing loci coincide.  We write
$\beta_{\mathrm{KR}}$ for the ordinary invariant of the right-Verma cover;
then
\begin{equation}\label{eq:vacuum-beta-normalization}
 \beta_E=0\quad\Longleftrightarrow\quad
 \beta_{\mathrm{KR}}=0.
\end{equation}
Since the staggered datum has rank one, $\beta_{\mathrm{KR}}=0$ is exactly
the condition that all ordinary beta-invariants appearing in the
rank-one analysis of Section~7 vanish.

Set $t=q/p$ in the notation of \cite{KytolaRidout}.  The right lowest
weight satisfies
\begin{equation}\label{eq:vacuum-braid-weight}
 h_{2p-1,1}=h_{1,2q-1}.
\end{equation}
Moreover $p\nmid(2p-1)$ and $q\nmid1$.  Hence the corresponding right
Verma module is of braid type by the classification in
\cite[Section~2]{KytolaRidout}.  Its first two prime singular vectors occur
at relative grades $2p-1$ and $2q-1$.  In the braid notation of
\cite[Section~7]{KytolaRidout}, the lower-grade first prime is $X^-_1$
and the higher-grade first prime is $X^+_1$.  The quotient defining
$K_{2p-1}$ kills the $(2p-1,1)$ prime singular vector, hence its defining
operator is
\[
 X=\begin{cases}
 X^-_1,&p<q,\\
 X^+_1,&p>q.
 \end{cases}
\]
(The equality $p=q$ is impossible because $p,q\ge2$ are coprime.)
Corollary~4.7 of \cite{KytolaRidout} lifts $E$ to a staggered module
$\check E$ with the same left module and with right module equal to this
Verma module.  Since $E$ is obtained by descending through the quotient by
$UXv_{h^R}$, Proposition~7.2 supplies a singular lift of $Xv_{h^R}$ in
$\check E$; the necessary condition stated immediately before that
proposition also gives $X\omega_0=0$.

Suppose that $\beta_E=0$.  By
\eqref{eq:vacuum-beta-normalization}, the ordinary invariant of the
right-Verma cover also vanishes.  We distinguish the two prime branches.

If $p<q$, then $2p-1<2q-1$, so the defining operator is
$X=X^-_1$, the smallest positive-grade prime singular vector.  The
rank-one datum has only the single ordinary invariant
$\beta_{\mathrm{KR}}$, so all ordinary beta-invariants vanish.  Together
with $X\omega_0=0$, this verifies every hypothesis of Proposition~7.5 of
\cite{KytolaRidout}; that proposition applies to the present singular
lift.  It forces the right lowest weight to admit a positive Kac label
$(R,S)$ such that
\[
 p\mid R,\qquad q\mid S.
\]
Write $R=pA$ and $S=qB$.  Since the right lowest weight is
$h_{2p-1,1}$, the Kac formula gives
\[
 pq(A-B)=\pm\bigl(q(2p-1)-p\bigr)
        =\pm(2pq-p-q).
\]
Reducing modulo $p$ yields $0\equiv\mp q\pmod p$, contradicting
$\gcd(p,q)=1$.  Thus $\beta_E$ cannot vanish in this case.

Assume now that $p>q$.  Then
\[
 \ell^-_1=2q-1<\ell^+_1=2p-1,
\]
and the quotient defining $K_{2p-1}$ kills the higher first prime
$X=X^+_1$.  Proposition~7.5 is therefore not the applicable
minimal-prime statement.  Under the present assumption
$\beta_{\mathrm{KR}}=0$, all ordinary beta-invariants of the Verma-cover
staggered module vanish.  Since $E$ is obtained by descent through the
quotient by $UX^+_1v_{h^R}$, Proposition~7.2 of
\cite{KytolaRidout} gives the required singular lift; equivalently the
associated transverse obstruction $\bar\beta_{X^+_1}$ vanishes.
Lemma~\ref{lem:higher-prime-determinant}, the precise
Kyt\"ol\"a--Ridout higher-first-prime specialization needed here,
therefore gives
\begin{equation}\label{eq:vacuum-required-det-order}
 \operatorname{ord}_{h=h^R}\det B_{\ell^+_1}(h)
 \ge \mathfrak p(\ell^+_1-\ell^-_1)+2
 =\mathfrak p(2p-2q)+2.
\end{equation}

For this one vacuum-edge weight the actual order can be read off directly
from the Kac determinant formula.  Put
\[
 D=2pq-p-q>0.
\]
The equality $h_{R,S}=h^R=h_{2p-1,1}$ for positive integers $R,S$ is
equivalent to
\[
 |qR-pS|=D.
\]
The solutions of $qR-pS=D$ are
\[
 (R,S)=(2p-1+mp,\,1+mq),\qquad m\in\Z,
\]
while those of $qR-pS=-D$ are
\[
 (R,S)=(1+mp,\,2q-1+mq),\qquad m\in\Z.
\]
Under the determinant bound $RS\le2p-1$, positivity leaves only
\[
 (R,S)=(2p-1,1),\qquad (R,S)=(1,2q-1).
\]
Consequently the order of the Kac-determinant zero at level $2p-1$ is
exactly
\begin{equation}\label{eq:vacuum-actual-det-order}
 \mathfrak p\bigl((2p-1)-(2q-1)\bigr)+\mathfrak p(0)
 =\mathfrak p(2p-2q)+1.
\end{equation}
This contradicts \eqref{eq:vacuum-required-det-order}.  Hence
$\beta_E\ne0$ also on the higher-prime branch.

Thus $\beta_E\ne0$ for every ordered coprime pair $p,q\ge2$.  If
\eqref{eq:vacuum-edge-quotient} were split,
Lemma~\ref{lem:split-forces-zero-coupling} would instead give
$\beta_E=0$, a contradiction.  Therefore
\eqref{eq:vacuum-edge-quotient} is nonsplit.
\end{proof}

\begin{prop}[One-row distinguished quotient]
\label{prop:one-row-staggered-quotient}
Let $n\ge1$, $1\le r\le p-1$, and put
\[
 a=np-r,\qquad b=np+r.
\]
If
\[
 0\longrightarrow K_a\longrightarrow E\longrightarrow K_b
 \longrightarrow0
\]
is logarithmic in $\Oc_{c_{p,q}}$, then the distinguished quotient
fits into a nonsplit exact sequence
\begin{equation}\label{eq:one-row-distinguished-nonsplit}
 0\longrightarrow S_a\longrightarrow E/S_b\longrightarrow K_b
 \longrightarrow0,
\end{equation}
where $S_b\subset K_a\subset E$ is the distinguished simple submodule.
\end{prop}

\begin{proof}
By Lemma~\ref{lem:Nakano-membership}, all terms lie in $\mathscr C$, and
Proposition~\ref{prop:Felder-labels} identifies the endpoints with
$A(\tau_{n,r})$ and $B(\tau_{n,r})$.  If $a\ge2$, Lemma~\ref{lem:one-row-Nakano-scope} verifies, for this
particular and otherwise arbitrary middle term $E$, every condition of
\cite[Definitions~3.14 and~5.1]{Nakano}.  Proposition~\ref{prop:Felder-labels}
and Lemma~\ref{lem:local-staggered-realization} identify the endpoint
modules and the distinguished Shapovalov element with those used in the
Ext group (5.1) of \cite{Nakano}.  The quantifier in
\cite[Theorem~5.11]{Nakano} is explicitly universal: it takes an
\emph{arbitrary logarithmic module} representing a class in (5.1).  The
theorem therefore applies to the present $E$ and gives the stronger
conclusion that $E/S_b$ is indecomposable.  In particular,
\eqref{eq:one-row-distinguished-nonsplit} is nonsplit.  The only remaining
possibility is $a=1$, equivalently $(n,r)=(1,p-1)$, and this is
Lemma~\ref{lem:vacuum-edge-staggered}.
\end{proof}

\begin{cor}[Extension input for the one-row chain]\label{cor:Nakano-one-row}
Let $n\ge1$, $1\le r\le p-1$, and set
\[
 a=np-r,\qquad b=np+r,\qquad c=(n+2)p-r.
\]
Fix endpoint identifications
$K_a\xrightarrow{\sim}A(\tau_{n,r})$ and
$K_b\xrightarrow{\sim}B(\tau_{n,r})$.  If $E\in\mathscr C$ is
logarithmic and
\[
 0\longrightarrow K_a\longrightarrow E\longrightarrow K_b
 \longrightarrow0
\]
is exact in $\Oc_{c_{p,q}}$, then
\[
 0\longrightarrow S_a\longrightarrow E/S_b\longrightarrow K_b
 \longrightarrow0
\]
is nonsplit.  Moreover,
\[
 \Ext^1_{\mathscr C}(K_b,S_a)\cong\C,
 \qquad
 \Ext^1_{\mathscr C}(K(\tau_{n,r}),S_b)\cong\C.
\]
\end{cor}

\begin{proof}
The dictionary and the first Ext dimension are
Proposition~\ref{prop:Felder-labels} and
\eqref{eq:Nakano-first-Ext}.  By
Lemma~\ref{lem:one-row-self-Ext},
\[
 \Ext^1_{\mathscr C}(S_b,S_b)=0.
\]
Proposition~\ref{prop:one-row-staggered-quotient} proves
\eqref{eq:NS-tau} for $\tau=\tau_{n,r}$.  Thus all hypotheses of
Theorem~\ref{thm:Nakano-input} hold, and that theorem gives
\[
 \Ext^1_{\mathscr C}(K(\tau_{n,r}),S_b)\cong\C.
\]
\end{proof}

\begin{prop}[Preliminary package]\label{prop:preliminary-package}
The data from the source and target categories used later satisfy the
following properties.
\begin{enumerate}
\item The indecomposable affine projectives are partitioned into wall
objects and the reflection chains of Lemma~\ref{lem:reflection-partition};
their Hom spaces and radical generators are those of
Proposition~\ref{prop:source-Hom-package}, and the affine minus twist has a
nonzero radical part as in Lemma~\ref{lem:source-L0-radical}.
\item Tensoring by $K_2$ is exact.  A one-row Kac module is supported in a
single conformal-residue sector, and projection to a residue sector is exact.
\item For $a=np-r$, $b=np+r$, and $c=(n+2)p-r$, the concrete modules
$S_a,S_b,S_c,K_a,K_b$ lie in $\mathscr C$, with
$K_a\cong A(\tau_{n,r})$ and $K_b\cong B(\tau_{n,r})$.  Since
$\mathscr C$ is Serre, every short exact sequence used later with end
terms in $\mathscr C$ has its middle term there as well.
\end{enumerate}
\end{prop}

\begin{proof}
This is the content of Proposition~\ref{prop:source-Hom-package},
Lemmas~\ref{lem:source-L0-radical}, \ref{lem:Kac-residue-support}, and
\ref{lem:block-projection}, Proposition~\ref{prop:Felder-labels}, and
Lemma~\ref{lem:Nakano-membership}.  Exactness of $K_2\boxtimes-$ follows
above from its rigidity and self-duality.
\end{proof}

\begin{rmk}[Scope of the extension input]\label{rmk:dependency-boundary}
The one-row extension package above is formulated in the Serre
subcategory $\mathscr C\subset\Oc_{c_{p,q}}$.  We do not identify this
subcategory with the ambient category of \cite[Definition~3.14]{Nakano}.
Instead, Lemma~\ref{lem:one-row-Nakano-scope} verifies Nakano's ambient
hypotheses object by object when $a\ge2$, while the unique vacuum edge
$a=1$ is treated directly by the Kyt\"ol\"a--Ridout staggered-module
classification.  No unproved socle or self-duality description of a
general $P(\tau)$ is used.  In particular, no left exactness of $F_{p,q}$
enters Section~\ref{sec:images}.
\end{rmk}
\section{Drinfeld--Sokolov reduction}\label{sec:DS}

We now turn to the homological construction that will eventually be
compared with the McRae--Yang functor.  We use principal $+$ quantum
Drinfeld--Sokolov reduction and write
\[
 H:=H^0_{DS,+}\big|_{\KL^k(\mathfrak{sl}_2)}.
\]
Throughout this section, $W^k(\mathfrak{sl}_2)$ denotes the universal
principal $W$-algebra.  In type $A_1$, the principal DS conformal vector
$\omega_{DS}$ of \cite[Section~4.17]{ArakawaW} satisfies the Virasoro
relations with central charge $c(k)$.  Hence there is a homomorphism of
conformal vertex algebras
\[
 \varphi_k:V^{\mathrm{Vir}}_{c(k)}
 \longrightarrow W^k(\mathfrak{sl}_2)
\]
sending the universal Virasoro conformal vector to $\omega_{DS}$.
By \cite[Proposition~4.12.1]{ArakawaW}, in type $A_1$ the universal
principal $W$-algebra is freely strongly generated by a single nonzero
field of conformal weight two.  Its conformal vector $\omega_{DS}$ is
nonzero and the weight-two subspace is one-dimensional.  Hence
\[
 W^k(\mathfrak{sl}_2)_2=\C\omega_{DS},
\]
and the conformal homomorphism $\varphi_k$ is surjective.  The same PBW theorem gives
\[
 \operatorname{ch}W^k(\mathfrak{sl}_2)
 =\prod_{n\ge2}(1-z^n)^{-1}
 =\operatorname{ch}V^{\mathrm{Vir}}_{c(k)}.
\]
Both sides are nonnegatively graded with finite-dimensional graded
pieces; the surjective graded map $\varphi_k$ is consequently injective
degree by degree.  Thus
\[
 W^k(\mathfrak{sl}_2)\cong V^{\mathrm{Vir}}_{c(k)}
\]
as conformal vertex algebras.  Here both sides are universal vertex
algebras; no simple quotient is being taken.  In particular, the Hamiltonian induced
on BRST cohomology is the Virasoro operator denoted below by
$L^{DS}_0$.  Every BRST cohomology group below is first regarded as a
module for this universal Virasoro vertex algebra; membership in the
finite-length category $\Oc_{c_{p,q}}$ will be proved, not assumed.

We fix the normalization before invoking any BRST result.  Our affine
level and shifted level are
\[
 k_{\mathrm{aff}}:=k=-2+\frac pq,
 \qquad
 t:=k_{\mathrm{aff}}+h^\vee=k+2=\frac pq
 \quad(h^\vee=2).
\]
Let $\kappa_0$ be the invariant form used in
\cite[Section~3.1]{ArakawaFrenkel}, normalized so that the long root has
squared length $2$.  To avoid a clash
between the affine VOA level and the bilinear-form parameter of
Arakawa--Frenkel, write $\kappa_{\mathrm{AF}}$ for their bilinear-form parameter.
Their affine algebra is defined at bilinear-form level
$\kappa_{\mathrm{AF}}+\kappa_c$, and in the rank-one scalar convention of
\cite[Section~4.4]{ArakawaFrenkel},
\[
 \kappa_c=-2\kappa_0,\qquad
 \kappa_{\mathrm{AF}}=t\kappa_0,
\]
so the actual affine form is
\[
 \kappa_{\mathrm{AF}}+\kappa_c=(t-2)\kappa_0=k\kappa_0.
\]
This is the convention in which the affine algebra attached to
$\kappa_{\mathrm{AF}}$ has actual cocycle
$\kappa_{\mathrm{AF}}+\kappa_c$.  Thus their rank-one scalar parameter is
\[
 \gamma=\frac{\kappa_{\mathrm{AF}}}{\kappa_0}=t=\frac pq.
\]
Accordingly, the affine central extension used in
$W_{\kappa_{\mathrm{AF}}}(\mathfrak{sl}_2)$ is the same central
extension as that used in $W^k(\mathfrak{sl}_2)$.  Under the chain-level
identification below we therefore identify
\[
 W_{\kappa_{\mathrm{AF}}}(\mathfrak{sl}_2)
 =H^0_{DS,0}(V_{\kappa_{\mathrm{AF}}}(\mathfrak{sl}_2))
 \cong H^0_{DS,+}(V^k(\mathfrak{sl}_2))
 =W^k(\mathfrak{sl}_2).
\]
For the $r$-dimensional horizontal $\mathfrak{sl}_2$-module we take
\[
 \lambda=(r-1)\omega,\qquad
 \check\mu=0.
\]
No symbol $\kappa$ below is used simultaneously for the actual affine
level and for the shifted Arakawa--Frenkel bilinear form.  With the
normalization above, our affine Weyl module $V_r$ is the
Arakawa--Frenkel Weyl module
$\mathbb V_{\lambda,\kappa_{\mathrm{AF}}}$: both are induced from the
same $r$-dimensional irreducible horizontal $\mathfrak{sl}_2$-module,
with $\mathfrak{sl}_2[t]t$ acting trivially, and the affine cocycle is
defined by
\[
 \kappa_{\mathrm{AF}}+\kappa_c=k\kappa_0.
\]
At the present rational parameter $\kappa_{\mathrm{AF}}/\kappa_0=p/q$,
we do not identify this Weyl module with the corresponding irreducible
highest-weight module.
We next fix the comparison between the two BRST conventions used below.

\begin{lem}[Comparison of the rank-one BRST conventions]
\label{lem:BRST-convention-comparison}
Choose the positive-root vector so that the principal characters in
\cite{ArakawaFrenkel} and \cite{ArakawaW} both take the value $1$ on it.
For $\mathfrak{sl}_2$, the $\check\mu=0$ Arakawa--Frenkel complex and
Arakawa's principal $+$ complex are naturally isomorphic for every
smooth level-$k$ affine module, in particular for every object and
subquotient used in this paper.  The identification is functorial in $M$.  Under this identification their
cohomological gradings, differentials, and induced principal
$W$-algebra actions agree.  Consequently
\[
 H^\bullet_{DS,0}(M)\cong H^\bullet_{DS,+}(M)=H^\bullet_+(M)
\]
naturally in $M$.
\end{lem}

\begin{proof}
We make the mode conventions explicit.  Fix root vectors with
$(e_\alpha,e_{-\alpha})=1$.  In the notation of
\cite{ArakawaFrenkel},
\[
 \psi^{AF}_{\alpha,n}=e_\alpha t^n,
 \qquad
 (\psi^{AF}_{\alpha,n})^*=e_\alpha^*t^{n-1}dt.
\]
The residue pairing, together with the invariant-form identification
$e_\alpha^*\leftrightarrow e_{-\alpha}$, pairs
$e_\alpha t^m$ with $e_\alpha^*t^{n-1}dt$ by $\delta_{m+n,0}$.
Consequently there is no shift of the underlying Clifford mode
index.  Let
\[
 \iota_{\mathrm{gh}}:\mathcal F^{AF}_{\mathrm{gh}}
 \xrightarrow{\sim}\mathcal F_{\mathrm{gh}}
\]
be the Clifford-module isomorphism determined by
\[
 \begin{aligned}
 \psi^{AF}_{\alpha,n}&\longmapsto\psi_\alpha(n),&
 (\psi^{AF}_{\alpha,n})^*&\longmapsto\psi_{-\alpha}(n),\\
 \one_{\mathrm{gh}}&\longmapsto\one_{\mathrm{gh}}.&&
 \end{aligned}
\]
This assignment preserves the Clifford relations.  Indeed,
\[
 [\psi_\alpha(m),\psi_{-\alpha}(n)]_+
 =\delta_{m+n,0}
 =[(\psi^{AF}_{\alpha,n})^*,\psi^{AF}_{\alpha,m}]_+.
\]
Thus there is no hidden shift or scalar in the ghost-mode
identification.  The Laurent powers used for the positive-root and negative-root ghost
fields in Arakawa's notation do not alter the underlying Clifford mode
labels.  The vacuum conditions are exactly
\cite[(2.3)]{ArakawaFrenkel}:
\[
 \psi^{AF}_{\alpha,n}\one_{\mathrm{gh}}=0\quad(n\ge0),
 \qquad
 (\psi^{AF}_{\alpha,m})^*\one_{\mathrm{gh}}=0\quad(m>0),
\]
and these agree with Arakawa's positive-real-root annihilation convention
under the displayed identification.  The cohomological degrees also agree:
\[
 \deg\psi^{AF}_{\alpha,n}=-1=\deg\psi_\alpha(n),
 \qquad
 \deg(\psi^{AF}_{\alpha,n})^*=1=\deg\psi_{-\alpha}(n).
\]
Here the subscript $0$ in $H^\bullet_{DS,0}$ means
$\check\mu=0$ in the twisted family of \cite{ArakawaFrenkel}; it does
not mean that the principal Whittaker character is zero.  For
$\check\mu=0$, their character contribution is the mode
$(\psi^{AF}_{\alpha,1})^*$, which corresponds exactly to Arakawa's
$\psi_{-\alpha}(1)$.  Equivalently, with the chosen root vector the
principal character is
\[
 \Psi(e_\alpha t^n)=\delta_{n,-1}.
\]

For $\check\mu=0$, formulas~(2.4) and (2.7) of
\cite{ArakawaFrenkel} reduce in type $A_1$ to
\[
 d_{\check\mu=0}=d_{\mathrm{st}}+\widehat\Psi,
 \qquad
 \widehat\Psi=(\psi^{AF}_{\alpha,1})^*.
\]
On the other hand, formulas~(166)--(168) of \cite{ArakawaW} give
\[
 Q_+=Q^{\mathrm{st}}_++\psi_{-\alpha}(1).
\]
Because $\mathfrak n_+$ is one-dimensional and abelian, the cubic ghost
term vanishes, and the standard part is
\[
 d_{\mathrm{st}}
 =\sum_{n\in\Z}e_\alpha(-n)(\psi^{AF}_{\alpha,n})^*
 \longleftrightarrow
 \sum_{n\in\Z}J_\alpha(-n)\psi_{-\alpha}(n)
 =Q_+^{\mathrm{st}}.
\]
The character term satisfies
\[
 \widehat\Psi=(\psi^{AF}_{\alpha,1})^*
 \longleftrightarrow
 \psi_{-\alpha}(1).
\]
Set
\[
 I_M:=\id_M\otimes\iota_{\mathrm{gh}}.
\]
The preceding formulas give the chain identity
\[
 I_M\,d_{\check\mu=0}=Q_+\,I_M,
\]
with the same cohomological grading and vacuum vector.  Thus $I_M$ is a natural isomorphism of the two BRST complexes.  For
$M=V^k(\mathfrak{sl}_2)$ the same identification is an isomorphism of
differential graded vertex algebras; for general $M$ it is an
isomorphism of differential graded modules over this common vacuum
complex.

For $\check\mu=0$ the spectral-flow automorphism
$\sigma_{\check\mu}$ of \cite[Section~2.1]{ArakawaFrenkel} is the
identity.  Thus the induced $W$-action is the ordinary vacuum BRST action.
The preceding chain isomorphism is consequently an isomorphism of
differential graded modules over the vacuum BRST complex, and the induced
cohomology isomorphism is $W$-linear and natural in $M$.
\end{proof}

We use these naturally identified cohomology functors without further
notational distinction.  With the Sugawara conformal structure at affine
level $k$, the lowest conformal weight is
$\Delta^{\mathrm{aff}}_r=(r^2-1)/(4(k+2))$.  In the rank-one notation of
\cite[Section~4.4]{ArakawaFrenkel}, with $m=r-1$ and
$\gamma=k+2$, the reduced lowest weight is
$\Delta^\gamma_{m,0}=m(m+2)/(4\gamma)-m/2$.  Therefore
\begin{equation}\label{eq:DS-weight-dictionary}
\begin{aligned}
 \Delta^{\mathrm{aff}}_r&=\frac{r^2-1}{4(k+2)},\\
 \Delta^\gamma_{r-1,0}
 &=\Delta^{\mathrm{aff}}_r-\frac{r-1}{2}
  =\frac{q(r^2-1)-2p(r-1)}{4p}\\
 &=\frac{(qr-p)^2-(p-q)^2}{4pq}=h_r.
\end{aligned}
\end{equation}
With these conventions the reduced Virasoro central charge is
\begin{equation}\label{eq:DS-central-charge}
 c(k)=13-6(k+2)-\frac{6}{k+2}
 =1-\frac{6(p-q)^2}{pq}=c_{p,q}.
\end{equation}
Thus all reduced modules below are taken with the conformal structure
defining $\Oc_{c_{p,q}}$.

\begin{lem}[Hamiltonian on the BRST complex]\label{lem:BRST-Hamiltonian}
On
\[
 C_+(M)=M\otimes\mathcal F_{\mathrm{gh}}
\]
let $L^{C(M)}_0$ denote the zero mode of the conformal vector on the BRST
complex.  With the $+$ convention fixed above,
\begin{equation}\label{eq:BRST-Hamiltonian}
 L^{C(M)}_0=
 L^{\mathrm{aff}}_0-\frac12h_0+
 \bigl(L^{\mathrm f}_0-J^{\mathrm{gh}}_0\bigr).
\end{equation}
Here $L^{\mathrm f}_0$ is the charged-fermion energy operator and
\[
 J^{\mathrm{gh}}(z):=:\psi_\alpha(z)\psi_{-\alpha}(z):
\]
is the ghost contribution to the diagonal action of
$\bar\rho^\vee=h/2$, with zero mode $J^{\mathrm{gh}}_0$ and the sign
convention of \cite[Section~4.17]{ArakawaW}.  The operator
$L^{C(M)}_0$ commutes with the BRST differential and induces the Virasoro
operator $L^{DS}_0$ on $H^0_{DS,+}(M)$.
\end{lem}

\begin{proof}
We spell out the sign convention because it is used again in
Section~\ref{sec:comparison}.  In the notation of
\cite[Section~4.17]{ArakawaW}, the BRST conformal field is
\[
 L(z)=L^{\mathfrak g}(z)+L^{\mathrm f}(z)
      +\frac{d}{dz}\widehat{\bar\rho^\vee}(z).
\]
Here $L^{\mathfrak g}(0)$ is the Sugawara Hamiltonian on the affine
module, hence $L^{\mathfrak g}(0)=L^{\mathrm{aff}}_0$ in our notation.
For $\mathfrak{sl}_2$ one has $\rho^\vee=h/2$ and the unique positive root has height $1$.  Hence the ghost term in $\widehat{\bar\rho^\vee}(z)$ is exactly $J^{\mathrm{gh}}(z)=:\psi_\alpha(z)\psi_{-\alpha}(z):$, while the diagonal
action on the ghost factor gives
\begin{equation}\label{eq:diagonal-rho-action}
 \widehat{\bar\rho^\vee}(0)=\frac12h_0+J^{\mathrm{gh}}_0.
\end{equation}
We use the mode convention $a(z)=\sum_{n\in\Z}a_nz^{-n-1}$; hence
$(\partial a)_0=-a_0$.  The derivative term therefore contributes
$-(\frac12h_0+J^{\mathrm{gh}}_0)$ to the zero mode, and
\[
 L^{C(M)}_0
 =L^{\mathrm{aff}}_0+L^{\mathrm f}_0
   -\left(\frac12h_0+J^{\mathrm{gh}}_0\right),
\]
which is \eqref{eq:BRST-Hamiltonian}.  The identity $[Q_+,L(z)]=0$ in
\cite[Section~4.17]{ArakawaW} is an operator identity in the BRST vertex
algebra.  Therefore the same commutator identity holds on every BRST module
complex $C_+(M)$; in particular,
\[
 [d_M,L^{C(M)}_0]=0.
\]
Hence $L^{C(M)}_0$ descends to degree-zero BRST cohomology, where it is the
Hamiltonian $L^{DS}_0$ of the principal Virasoro conformal vector.
\end{proof}


\subsection{BRST d\'evissage and exactness}

The vanishing theorem will be applied only to simple highest-weight
composition factors.  We do not apply Arakawa's category-$\mathcal O$
exactness statement directly to logarithmic objects of
$\KL^k(\mathfrak{sl}_2)$.  Since our category also contains finite-length generalized modules, we
record the d\'evissage that passes from the simple factors to arbitrary
objects.  We use only that $\KL^k(\mathfrak{sl}_2)$ is finite length and
abelian, that its underlying-module functor to smooth level-$k$ affine modules
is exact, and that the BRST complex is functorially
defined on the latter category.  In particular it is defined on every
subquotient appearing in a composition series.

\begin{lem}[Smoothness]\label{lem:KL-smooth}
Every object of $\KL^k(\mathfrak{sl}_2)$ is smooth as an affine module, and
so is every subquotient.
\end{lem}

\begin{proof}
Fix a basis $x_1,x_2,x_3$ of $\mathfrak{sl}_2$.  For every
$v\in M$, the vertex-operator truncation property gives integers
$N_i(v)$ such that $x_i(n)v=0$ for all $n\ge N_i(v)$.  Taking
$N(v)=\max_iN_i(v)$ gives
\[
 (\mathfrak{sl}_2\otimes t^{N(v)}\C[t])v=0,
\]
which is precisely smoothness.  The assertion is inherited by submodules
and quotients, and the forgetful functor to smooth affine modules is
exact.
\end{proof}

\begin{lem}[BRST d\'evissage]\label{lem:BRST-devissage}
Let $\mathscr A$ be a finite-length abelian category whose short exact
sequences remain exact as sequences of smooth affine modules, and suppose that the
$+$ BRST complex is defined functorially on its objects.
Suppose that
\[
 H^i_{DS,+}(L)=0\qquad(i\neq0)
\]
for every simple object $L$ of $\mathscr A$.  Then the same vanishing holds
for every object of $\mathscr A$, and $H^0_{DS,+}$ is exact on
$\mathscr A$.
\end{lem}

\begin{proof}
No semisimplicity of $L_0$ and no membership of the objects of
$\mathscr A$ in the affine BGG category $\mathcal O$ is used here;
finite length is the d\'evissage input.  For a smooth affine module $M$
the standard module BRST complex is
\[
 C_+(M)=M\otimes\mathcal F_{\mathrm{gh}},
 \qquad
 C_+^j(M)=M\otimes\mathcal F_{\mathrm{gh}}^j.
\]
Its differential is functorial in $M$; smoothness makes the standard
mode sum defining the differential locally finite on every vector.  Each
cochain functor $C_+^j(-)$ is exact because tensoring vector spaces over
$\C$ is exact.  Hence every
short exact sequence in $\mathscr A$ induces a degreewise short exact
sequence of BRST complexes.

We induct on Jordan--H\"older length.  Choose
\[
 0\longrightarrow M'\longrightarrow M\longrightarrow L\longrightarrow0
\]
with $L$ simple.  If the assertion is known for $M'$, the long exact cohomology sequence
first gives the vanishing of $H^i_{DS,+}(M)$ for $i\ne0$.  Around degree
zero it contains the five-term segment
\[
 H^{-1}_{DS,+}(L)\longrightarrow H^0_{DS,+}(M')
 \longrightarrow H^0_{DS,+}(M)\longrightarrow H^0_{DS,+}(L)
 \longrightarrow H^1_{DS,+}(M').
\]
The two outer terms vanish by the simple case and the induction hypothesis,
so this reduces to
\[
 0\longrightarrow H^0_{DS,+}(M')\longrightarrow H^0_{DS,+}(M)
 \longrightarrow H^0_{DS,+}(L)\longrightarrow0.
\]
This proves the induction step.  More generally, for any
short exact sequence in $\mathscr A$, the associated long exact
cohomology sequence and the vanishing in all nonzero degrees reduce to a
short exact sequence in degree zero.  Hence $H^0_{DS,+}$ is exact on
$\mathscr A$.
\end{proof}

\begin{prop}\label{prop:H-simple}
For every $r\geq1$,
\[
 H(L_r)\cong S_r,
 \qquad
 H^i_{DS,+}(L_r)=0\quad(i\neq0).
\]
In particular, every simple object has nonzero simple reduction.
\end{prop}

\begin{proof}
Put
\[
 \widehat\lambda_r=(r-1)\omega+k\Lambda_0.
\]
Then $L_r=L(\widehat\lambda_r)$ is an object of Arakawa's affine
category $\mathcal O_k$, and $k\ne-2$.  Condition~(383) of
\cite[Section~9.1]{ArakawaW}, equivalently condition~(385) there,
requires
\[
 \langle\widehat\lambda_r+\widehat\rho,\beta^\vee\rangle
 \notin\Z
\]
for
\[
 \beta\in\{-\bar\alpha+n\delta:
 \bar\alpha\in\bar\Delta_+,\ 1\le n\le
 \operatorname{ht}\bar\alpha\}.
\]
For $\mathfrak{sl}_2$ this set consists only of $-\alpha+\delta$,
whose coroot is
$-\alpha^\vee+K_{\mathrm{aff}}$, with $K_{\mathrm{aff}}$ the affine
central element.  If $\widehat\rho$ is the affine Weyl vector, then
\[
 \left\langle
 \widehat\lambda_r+\widehat\rho,
 (-\alpha+\delta)^\vee
 \right\rangle
 =(k+2)-r=\frac pq-r\notin\Z,
\]
because $q\ge2$ and $\gcd(p,q)=1$.  Thus the required real-root pairing is nonintegral for
every $r\ge1$.  Theorems~9.1.3--9.1.4 of \cite{ArakawaW} therefore
give
\[
 H^i_{DS,+}(L_r)=0\quad(i\neq0),
\]
and identify $H(L_r)$ with a nonzero irreducible principal $W$-module.

For $m=r-1$, the rank-one specialization of the transformed conformal
weight is
\[
 \Delta_{DS}(m)
 =\frac{m(m+2)}{4(k+2)}-\frac m2
 =\frac{q(r^2-1)-2p(r-1)}{4p}
 =\frac{(qr-p)^2-(p-q)^2}{4pq}=h_r.
\]
Under the conformal identification
$W^k(\mathfrak{sl}_2)\cong V^{\mathrm{Vir}}_{c_{p,q}}$ established
above, the reduced irreducible is therefore the simple lowest-weight
Virasoro module $S_r$.
\end{proof}


\begin{prop}\label{prop:H-exact}
For every $M\in\KL^k(\mathfrak{sl}_2)$,
\[
 H^i_{DS,+}(M)=0\qquad(i\neq0).
\]
Consequently $H$ is a $\C$-linear exact functor
\[
 H:\KL^k(\mathfrak{sl}_2)\longrightarrow\Oc_{c_{p,q}}.
\]
Moreover, if $\ell(-)$ denotes Jordan--H\"older length, then
\begin{equation}\label{eq:H-length}
 \ell(H(M))=\ell(M).
\end{equation}
\end{prop}

\begin{proof}
The simple objects of $\KL^k(\mathfrak{sl}_2)$ are the modules
$L_r$, $r\ge1$; see \cite[Section~2.2, in particular
Theorem~2.2]{McRaeYang}.  Lemma~\ref{lem:KL-smooth} verifies the underlying smoothness hypothesis,
and Proposition~\ref{prop:H-simple} verifies the cohomological hypothesis of
Lemma~\ref{lem:BRST-devissage} for every simple object.  Hence the asserted
vanishing and exactness follow.  The BRST chain
construction is $\C$-linear on morphisms, and passage to cohomology
preserves this linearity; hence $H$ is $\C$-linear.  If
\[
 0=M_0\subset M_1\subset\cdots\subset M_\ell=M
\]
is a composition series, exactness gives a filtration of $H(M)$ with
successive quotients
\[
 H(M_i)/H(M_{i-1})\cong H(M_i/M_{i-1})\cong S_{r_i},
\]
which are nonzero simple objects by Proposition~\ref{prop:H-simple}.
Hence the resulting filtration of $H(M)$ is itself a composition series,
and \eqref{eq:H-length} follows.  Each successive quotient is
$S_{r_i}=L(c_{p,q},h_{r_i,1})$, hence is one of the simple objects allowed
in the definition of $\Oc_{c_{p,q}}$.  Repeated application of
Lemma~\ref{lem:gr-extension} to the displayed filtration shows that
$H(M)$ is grading restricted.  Thus $H(M)$ has the required composition
factors and belongs to $\Oc_{c_{p,q}}$.  No injectivity of the labels $r\mapsto S_r$ is
needed; every successive quotient is simply a nonzero simple object.
\end{proof}

\begin{cor}\label{cor:H-faithful}
The functor $H$ is faithful.
\end{cor}

\begin{proof}
Let $0\neq f:M\to N$.  Since $H$ is exact,
\[
 \operatorname{Im}H(f)\cong H(\operatorname{Im}f).
\]
The module $\operatorname{Im}f$ is nonzero and therefore has positive
length.  Equation~\eqref{eq:H-length} gives
\[
 \ell\bigl(H(\operatorname{Im}f)\bigr)>0.
\]
Hence $H(f)\neq0$.
\end{proof}

\subsection{Reduction of Weyl modules}

The next calculation is the one place where we need more than reduction of
simple highest-weight modules.  We use only the rank-one category-$\mathcal O$
statement for the particular affine Verma module mapping onto $V_r$; no
generic-level irreducibility or quantum-Langlands duality is involved.

\begin{prop}\label{prop:H-Weyl}
For every $r\geq1$,
\[
 H(V_r)\cong K_r.
\]
\end{prop}

\begin{proof}
Take $\lambda=(r-1)\omega$, $\kappa_{\mathrm{AF}}=t\kappa_0$, and
$\check\mu=0$ in Arakawa--Frenkel.  By Lemma~\ref{lem:BRST-convention-comparison}, the $\check\mu=0$
Arakawa--Frenkel functor is naturally identified with the principal $+$
functor used in this paper.

\emph{Vanishing.}  Their Theorem~2.1 gives
\[
 H^i_{DS,0}(V_r)=0\qquad(i\neq0)
\]
for every value of $\kappa_{\mathrm{AF}}$; no irrationality hypothesis is used here.

\emph{Cyclicity.}  Here
$\gamma=\kappa_{\mathrm{AF}}/\kappa_0=t=p/q\ne-2$.  The rank-one
statement in \cite[Section~4.4]{ArakawaFrenkel} applies for every
complex $\gamma\ne-2$: $H^0_{DS,0}$ is exact on the affine category
$\mathcal O_{\kappa}$ and sends an affine Verma module to the
corresponding Virasoro Verma module.  Write $M^{\kappa}_{r-1}$ for the affine Verma module of horizontal
highest weight $(r-1)\omega$.  The canonical quotient
\[
 M^{\kappa}_{r-1}\twoheadrightarrow V_r
\]
therefore induces a surjection
\[
 \mathcal V(c_{p,q},h_r)
 \cong H^0_{DS,0}(M^{\kappa}_{r-1})
 \twoheadrightarrow H(V_r).
\]
In particular, $H(V_r)$ is generated by the image of the affine
highest-weight vector.  This use of Section~4.4 is independent of the
irrational-level duality theorem; indeed, that section is precisely where
\cite{ArakawaFrenkel} discusses the failure of the duality statement at
rational parameter.

\emph{Character.}  Formula~(4.19) of \cite{ArakawaFrenkel} computes the
relative $\Z_{\ge0}$-graded character.  For $\mathfrak{sl}_2$,
$\check\mu=0$, and $\lambda=(r-1)\omega$, one has
\[
 \langle\lambda+\rho,\check\rho\rangle=\frac r2,
\]
so its two Weyl-group terms specialize to
\[
 z^{r/2}(z^{-r/2}-z^{r/2})=1-z^r.
\]
Hence the relative character is
\[
 \operatorname{ch}^{\mathrm{rel}}_z H(V_r)
 =\frac{1-z^r}{\prod_{m\ge1}(1-z^m)}.
\]
Formula~(4.19) is written in the auxiliary $\Z_{\ge0}$-grading of
Section~4.3 of \cite{ArakawaFrenkel}.  We identify that grading with the
Virasoro grading directly at the present rational level.  By
Lemma~\ref{lem:BRST-Hamiltonian},
\[
 L^{C(M)}_0=L^{\mathrm{aff}}_0-\frac12h_0
 +L^{\mathrm f}_0-J^{\mathrm{gh}}_0.
\]
On the affine generators this gives
\[
 [L^{C(M)}_0,e(n)]=(-n-1)e(n),\qquad
 [L^{C(M)}_0,f(n)]=(-n+1)f(n),\qquad
 [L^{C(M)}_0,h(n)]=-nh(n).
\]
For the ghosts, the Clifford OPE and
$J^{\mathrm{gh}}(z)=:\psi_\alpha(z)\psi_{-\alpha}(z):$ give
\[
 [L^{\mathrm f}_0,\psi_\alpha(n)]=-n\psi_\alpha(n),\qquad
 [L^{\mathrm f}_0,\psi_{-\alpha}(n)]=-n\psi_{-\alpha}(n),
\]
and
\[
 [J^{\mathrm{gh}}_0,\psi_\alpha(n)]=\psi_\alpha(n),\qquad
 [J^{\mathrm{gh}}_0,\psi_{-\alpha}(n)]=-\psi_{-\alpha}(n).
\]
Consequently
\[
 [L^{\mathrm f}_0-J^{\mathrm{gh}}_0,\psi_\alpha(n)]
 =(-n-1)\psi_\alpha(n),
\]
while
\[
 [L^{\mathrm f}_0-J^{\mathrm{gh}}_0,\psi_{-\alpha}(n)]
 =(-n+1)\psi_{-\alpha}(n).
\]
Under $(\psi^{AF}_{\alpha,n})^*\leftrightarrow\psi_{-\alpha}(n)$,
these are exactly the auxiliary degrees introduced in
\cite[Section~4.3]{ArakawaFrenkel} in the grading used to derive the
character formula~(4.19), specialized to $\check\mu=0$.  Explicitly,
that grading is
\[
 \deg e(n)=\deg\psi_\alpha(n)=-n-1,\qquad
 \deg f(n)=\deg\psi^*_\alpha(n)=-n+1,\qquad
 \deg h(n)=-n.
\]
Our normal ordering is fixed so that
\[
 L^{\mathrm f}_0\mathbf1_{\mathrm{gh}}
 =J^{\mathrm{gh}}_0\mathbf1_{\mathrm{gh}}=0.
\]
Let $G_{\mathrm{AF}}$ denote that auxiliary grading operator and put
\[
 A:=L^{C(V_r)}_0-h_r.
\]
By \eqref{eq:DS-weight-dictionary} and the normalization of the ghost
vacuum,
\[
 A(v_r\otimes\mathbf1)=0=G_{\mathrm{AF}}(v_r\otimes\mathbf1).
\]
The commutator formulas above show that $A$ and $G_{\mathrm{AF}}$ have
the same commutator with every affine and ghost mode used in the PBW
spanning set of the BRST complex.  Induction on the length of a PBW
monomial therefore gives
\[
 A\,w=G_{\mathrm{AF}}\,w
\]
for every PBW monomial $w$ applied to
$v_r\otimes\mathbf1$.  Indeed, $V_r$ is generated from its affine
highest-weight vector by the affine PBW operators (with the finite-dimensional
horizontal top generated by the zero mode $f(0)$), and the standard ghost
creation operators generate the ghost Fock space.  These PBW monomials
therefore span the BRST complex.  Hence
$A=G_{\mathrm{AF}}$ as grading operators on $C_+(V_r)$, and consequently on
cohomology.  Thus, at the rational admissible level itself, the auxiliary
degree is exactly $L^{DS}_0-h_r$.  Thus, as a formal graded
character with finite-dimensional homogeneous pieces, the unnormalized
character $\ch_zM=\operatorname{Tr}_M z^{L_0}$ gives
\begin{equation}\label{eq:H-character}
 \ch_z H(V_r)=
 \frac{z^{h_r}(1-z^r)}{\prod_{m\ge1}(1-z^m)}.
\end{equation}
The Verma-module map above is $L_0$-graded.  The Virasoro Verma module is
the algebraic direct sum of its finite-dimensional $L_0$-eigenspaces, so
its graded submodule kernel and its quotient $H(V_r)$ inherit semisimple
$L_0$-actions.  The following character comparison is therefore
coefficientwise in finite-dimensional graded pieces.
On the other hand, the one-row Kac character formula
\cite[Remark~2.2]{McRaeSopin} gives
\begin{equation}\label{eq:Kac-character}
 \ch_z K_r=
 \frac{z^{h_r}(1-z^r)}{\prod_{m\ge1}(1-z^m)}.
\end{equation}

Consequently,
\[
 \ch_z\ker\!\left(\mathcal V(c_{p,q},h_r)
 \twoheadrightarrow H(V_r)\right)
 =\frac{z^{h_r+r}}{\prod_{m\ge1}(1-z^m)}.
\]
In particular, the kernel is zero below level $r$ and one-dimensional at
level $r$.  A nonzero homogeneous vector of minimal relative degree in a proper
graded Virasoro submodule of a Verma module is annihilated by all positive
Virasoro modes, hence is singular.  Thus the level-$r$ kernel is a
singular line.  By the paragraph immediately preceding
\cite[Theorem~4.3]{McRaeSopin}, the Verma module
$\mathcal V(c_{p,q},h_r)$ contains, up to nonzero scalar, a unique singular
vector $v_{r,1}$ at relative degree $r$; the corresponding one-row Verma
quotient is $K_{r,1}$, since the second Kac label is $1<q$.
Thus the nonzero degree-$r$ kernel line is $\C v_{r,1}$, and the submodule
generated by $v_{r,1}$ is contained in the kernel.  By
\cite[Remark~2.2]{McRaeSopin},
\[
 K_r=K_{r,1}\cong
 \mathcal V(c_{p,q},h_r)/\langle v_{r,1}\rangle.
\]
The Verma-module surjection therefore factors through a grading-preserving
surjection
\[
 K_r\twoheadrightarrow H(V_r).
\]
Equations~\eqref{eq:H-character} and \eqref{eq:Kac-character} show that the
kernel of this factor map has zero graded character.  Its graded pieces are
finite dimensional, so every graded piece of the kernel vanishes.  Thus
$H(V_r)\cong K_r$.
\end{proof}

\begin{rmk}\label{rmk:Weyl-reduction-scope}
The proof uses only the all-level Weyl-module vanishing and character formula
of \cite{ArakawaFrenkel}, together with the rank-one Verma-quotient statement
of their Section~4.4.  The grading comparison is proved directly from
Lemma~\ref{lem:BRST-Hamiltonian}.  In particular, neither Theorem~2.2
nor Theorem~5.2 of \cite{ArakawaFrenkel} is used at rational
$\kappa_{\mathrm{AF}}$.
\end{rmk}

\subsection{The exponential comparison}

We finally isolate the convention-sensitive compatibility that will later
rigidify objectwise projective isomorphisms.  At this stage $H$ has not
been equipped with a tensor structure, so the following proposition is stated at the level of the underlying
linear functor, as an equality of exponentials rather than as a
ribbon-functor assertion.

\begin{lem}[Exponential and cohomology]\label{lem:exponential-cohomology}
Let $(C^\bullet,d)$ be a complex and let $T$ be a locally finite chain
endomorphism.  Then the induced endomorphism $H^\bullet(T)$ is locally
finite.  The algebraic exponentials are therefore well-defined and
\[
 H^\bullet(e^T)=e^{H^\bullet(T)}.
\]
\end{lem}

\begin{proof}
For $c\in C^\bullet$, the orbit space
\[
 U_c:=\operatorname{span}_{\C}\{T^nc:n\ge0\}
\]
is finite dimensional and $T$-stable.  Define $e^Tc$ by evaluating the
ordinary matrix exponential on $U_c$; this is independent of the choice
of a finite-dimensional $T$-stable space containing $c$.  Since $T$
commutes with $d$, so does $e^T$, and $e^T$ is a chain automorphism with
inverse $e^{-T}$.

If $c$ is a cocycle, then
\[
 \operatorname{span}_{\C}\{H^\bullet(T)^n[c]:n\ge0\}
\]
is the image in cohomology of the finite-dimensional space $U_c$.
Thus $H^\bullet(T)$ is locally finite.  On this finite-dimensional orbit
space the exponential is a polynomial in the relevant endomorphism, so
\[
 H^\bullet(e^T)[c]=[e^Tc]=e^{H^\bullet(T)}[c].
\]
\end{proof}

\begin{prop}[Exponential compatibility]\label{prop:H-twist}
If $M\in\KL^k_{\bar\epsilon}(\mathfrak{sl}_2)$, then
\begin{equation}\label{eq:H-exponential-compatibility}
 H(\theta^-_M)=e^{2\pi iL^{DS}_0}\big|_{H(M)}.
\end{equation}
Here $\theta^-_M=(-1)^\epsilon e^{2\pi iL^{\mathrm{aff}}_0}$ is the
non-standard affine twist from \eqref{eq:minus-twist}.
\end{prop}

\begin{proof}
The four operators
$L^{\mathrm{aff}}_0$, $h_0$, $L^{\mathrm f}_0$, and
$J^{\mathrm{gh}}_0$ commute pairwise.  In particular
$[L^{\mathrm{aff}}_0,h_0]=0$, the affine and ghost operators act on
separate tensor factors, and on the ghost factor $J^{\mathrm{gh}}_0$ is
the zero mode of a weight-one field for the charged-fermion conformal
structure, so
\[
 [L^{\mathrm f}_0,J^{\mathrm{gh}}_0]=0.
\]  By the explicit Clifford--Fock realization
\cite[(154)--(159)]{ArakawaW}, the ghost space
$\mathcal F_{\mathrm{gh}}$ has a basis of standard ghost monomials.  With
the convention fixed above,
\[
 J^{\mathrm{gh}}_0\one_{\mathrm{gh}}=0,
 \qquad
 [J^{\mathrm{gh}}_0,\psi_\alpha(n)]=\psi_\alpha(n),
 \qquad
 [J^{\mathrm{gh}}_0,\psi_{-\alpha}(n)]=-\psi_{-\alpha}(n).
\]
Thus $J^{\mathrm{gh}}_0$ has integral spectrum on the ghost Fock space.
The ghost conformal grading is integral as well, and the standard ghost
monomials are simultaneous eigenvectors for $L^{\mathrm f}_0$ and
$J^{\mathrm{gh}}_0$.  Hence
\[
 \operatorname{Spec}(L^{\mathrm f}_0-J^{\mathrm{gh}}_0)\subset\Z,
\]
and therefore
\begin{equation}\label{eq:ghost-exponential}
 e^{2\pi i(L^{\mathrm f}_0-J^{\mathrm{gh}}_0)}
 =\id_{\mathcal F_{\mathrm{gh}}}.
\end{equation}
Horizontal $\mathfrak{sl}_2$ acts locally finitely on objects of
$\KL^k(\mathfrak{sl}_2)$ by \cite[Section~2]{McRaeYang}.  Hence every
vector lies in a finite-dimensional horizontal $\mathfrak{sl}_2$-module,
and $h_0$ is semisimple.  For
$M\in\KL^k_{\bar\epsilon}$,
\[
 \operatorname{Spec}(h_0|_M)\subset\epsilon+2\Z,
 \qquad
 e^{-\pi ih_0}=(-1)^\epsilon\id_M.
\]
Because $M$ is grading restricted, every generalized
$L^{\mathrm{aff}}_0$-eigenspace is finite dimensional.  Since
$[L^{\mathrm{aff}}_0,h_0]=0$, each such generalized eigenspace is
$h_0$-stable.  Every vector of $M$ has only finitely many generalized
$L^{\mathrm{aff}}_0$-components, so it lies in a finite-dimensional
subspace stable under both operators.  Every ghost vector is a finite sum
of simultaneous eigenvectors
for $L^{\mathrm f}_0$ and $J^{\mathrm{gh}}_0$.  Consequently every vector
of $C_+(M)$ is contained in a finite-dimensional subspace invariant under
the four mutually commuting operators
\[
 L^{\mathrm{aff}}_0,\qquad h_0,\qquad
 L^{\mathrm f}_0,\qquad J^{\mathrm{gh}}_0.
\]
All exponentials below are taken on such finite-dimensional joint invariant
subspaces.  Moreover,
\[
 [L^{\mathrm{aff}}_0,x(n)]=-n\,x(n)
 \qquad(x\in\mathfrak{sl}_2,\ n\in\Z),
\]
so
\[
 e^{2\pi iL^{\mathrm{aff}}_0}x(n)e^{-2\pi iL^{\mathrm{aff}}_0}=x(n).
\]
Thus $\theta^-_M=(-1)^\epsilon e^{2\pi iL^{\mathrm{aff}}_0}$ is an
affine-module automorphism.  The usual exponential identity for
commuting operators is therefore purely algebraic and gives the
chain-level identity
\begin{align}
 e^{2\pi iL^{C(M)}_0}
 &=e^{2\pi iL^{\mathrm{aff}}_0}
   e^{-\pi ih_0}
   e^{2\pi i(L^{\mathrm f}_0-J^{\mathrm{gh}}_0)}\notag\\
 &=(-1)^\epsilon e^{2\pi iL^{\mathrm{aff}}_0}
   \otimes\id_{\mathcal F_{\mathrm{gh}}}\notag\\
 &=\theta^-_M\otimes\id_{\mathcal F_{\mathrm{gh}}}.
 \label{eq:complex-exponential}
\end{align}
On each of the finite-dimensional joint invariant subspaces just
described, the exponential is the scalar exponential times a finite
polynomial in the nilpotent parts; no analytic convergence is involved.
The left-hand side of \eqref{eq:complex-exponential} commutes with the
BRST differential by Lemma~\ref{lem:BRST-Hamiltonian}.  The right-hand
side is precisely the BRST chain map induced functorially by the
affine-module automorphism $\theta^-_M$, and therefore also commutes
with the differential.  Taking degree-zero cohomology and using
Lemmas~\ref{lem:BRST-Hamiltonian} and
\ref{lem:exponential-cohomology} therefore gives
\eqref{eq:H-exponential-compatibility}.  The chain-level identity is
functorial in $M$: an affine-module morphism commutes with
$L^{\mathrm{aff}}_0$ and $h_0$ and acts trivially on the ghost factor.
Thus \eqref{eq:H-exponential-compatibility} is an equality of natural
automorphisms on each parity subcategory.  Since
$\KL^k=\KL^k_{\bar0}\oplus\KL^k_{\bar1}$ and both sides are additive on
finite direct sums, the same identity holds on every object of $\KL^k$.

As a normalization check, on the lowest class of $V_r$ the reduced
exponential has eigenvalue $e^{2\pi ih_r}$, while the affine expression
has eigenvalue $(-1)^{r-1}e^{2\pi i\Delta^{\mathrm{aff}}_r}$, where
\[
 \Delta^{\mathrm{aff}}_r=\frac{r^2-1}{4(k+2)}.
\]
They agree because the parity of $V_r$ satisfies
$\epsilon\equiv r-1\pmod2$ and
$h_r=\Delta^{\mathrm{aff}}_r-(r-1)/2$.
\end{proof}

\begin{thm}[Drinfeld--Sokolov package]\label{thm:DS-package}
Principal $+$ reduction restricts to a well-defined $\C$-linear functor
\[
 H:\KL^k(\mathfrak{sl}_2)\longrightarrow\Oc_{c_{p,q}}.
\]
On $\KL^k(\mathfrak{sl}_2)$ one has
\[
 H^i_{DS,+}(M)=0\qquad(i\ne0,\,M\in\KL^k),
\]
and principal $+$ reduction is exact and faithful, preserves
Jordan--H\"older length, and satisfies
\[
 H(V_r)\cong K_r,\qquad H(L_r)\cong S_r\qquad(r\ge1).
\]
On the BRST complex the exponential identity
\eqref{eq:complex-exponential} holds, and on cohomology it gives
\eqref{eq:H-exponential-compatibility}.  The last identity is an identity
of natural automorphisms of the underlying linear functor; no monoidal
structure on $H$ is asserted here.
\end{thm}

\begin{proof}
The simple images and higher-cohomology vanishing are
Proposition~\ref{prop:H-simple}; exactness and length preservation are
Proposition~\ref{prop:H-exact}; Weyl images are
Proposition~\ref{prop:H-Weyl}; faithfulness is
Corollary~\ref{cor:H-faithful}; and exponential compatibility is
Proposition~\ref{prop:H-twist}.  These conclusions are obtained entirely
from the BRST analysis of this section.  In particular,
Section~\ref{sec:DS} is logically independent of the projective
reconstruction and of the Virasoro extension theory of
Section~\ref{sec:prelim}.
\end{proof}

\begin{lem}[Normalization of the Drinfeld--Sokolov branch]
\label{lem:FKW-normalization}
With the root-vector normalization fixed in
Lemma~\ref{lem:BRST-convention-comparison}, the functor
$H^0_{DS,+}$ is the rank-one principal quantized
Drinfeld--Sokolov reduction with character
\[
 \Psi(e_\alpha t^n)=\delta_{n,-1}.
\]
For the ordered pair $(p,q)$ it is the branch for which
\[
 H^0_{DS,+}(V_2)\cong K_{2,1}.
\]
Consequently it is the quantum Drinfeld--Sokolov reduction paired with
$F_{p,q}$ in Conjecture~7.16 of \cite{McRaeYang}.  After interchanging
$p$ and $q$, the same normalization gives
\[
 K^{(q,p)}_{2,1}\cong K^{(p,q)}_{1,2},
\]
which is the branch paired there with $F_{q,p}$.
\end{lem}

\begin{proof}
Arakawa's $+$ cohomology is one of the standard quantized principal
Drinfeld--Sokolov reductions originating in
\cite{FKW}; see also \cite{ArakawaVanishing,ArakawaW}.  The explicit
comparison in Lemma~\ref{lem:BRST-convention-comparison} fixes the
principal character to be $\Psi(e_\alpha t^n)=\delta_{n,-1}$, removing
any ambiguity from root-vector or ghost conventions.  Proposition~\ref{prop:H-Weyl}
then gives $H^0_{DS,+}(V_2)\cong K_{2,1}$.  McRae--Yang's
Conjecture~7.16 refers to the two standard quantum Drinfeld--Sokolov
restrictions without imposing a $+$/$-$ notation; their Theorem~7.15,
however, distinguishes the ordered functors by the rigid Kac objects
$K_{2,1}$ and $K_{1,2}$.  Thus the displayed Weyl image fixes which
standard reduction is associated to the ordered pair $(p,q)$.  The
identity $c_{p,q}=c_{q,p}$ and the Kac weight/quotient formulas give
$K^{(q,p)}_{2,1}\cong K^{(p,q)}_{1,2}$ after swapping the parameters,
which gives the second assertion.
\end{proof}

\section{Images of projective modules}\label{sec:images}

From now on write $F=F_{p,q}$ and
\[
 Q_j:=F(P_j).
\]
For $h\in\C$ we abbreviate $M^{[h]}:=M^{[h+\Z]}$.  The following
proof protocol will be used throughout this section.  The functor $F$ is
applied only to morphisms, isomorphisms, and finite biproduct
decompositions in the projective subcategory, together with its strong
monoidal structure.  Since $F$ is $\C$-linear, the images of biproduct
inclusions and projections still satisfy the biproduct identities.
Every short exact sequence involving an object $Q_j=F(P_j)$ is
constructed entirely in $\Oc_{c_{p,q}}$, either from the Virasoro
$K_2$-fusion sequence or from a previously constructed target-side exact
sequence by the exact functors $K_2\boxtimes-$ and $(-)^{[\xi]}$.
In particular, no monomorphism in a nonsplit affine short exact sequence
is asserted to remain monic under $F$.  Twist preservation is used only
after projective faithfulness has been established.  Non-splitting of the
reconstructed target sequences is likewise proved only at that later
stage.

\subsection{Small Weyl modules}

\begin{lem}\label{lem:F-small-Weyl}
For $1\leq r\leq p$,
\[
 F(V_r)\cong K_r.
\]
\end{lem}

\begin{proof}
Recall from Section~\ref{sec:prelim} that $P_r=V_r$ for $1\le r<p$,
while $P_p=V_p=L_p$ is wall simple-projective.  The case $r=1$ follows from the unit constraint
$F(V_1)\cong\mathbf1_{\Oc}=K_1$, while $r=2$ is the defining normalization
$F(V_2)\cong K_2$.  If $p=2$, these two cases exhaust the assertion.
Assume that $F(V_j)\cong K_j$ for every $1\le j\le r$, where
$2\le r\le p-1$.  Since $2\le r\le p-1$, one has $p\nmid r$, and
\cite[Theorem~2.23]{McRaeYang} gives directly the biproduct
decomposition
\[
 V_2\boxt V_r\cong V_{r-1}\oplus V_{r+1}.
\]
Applying the additive strong monoidal functor $F$ gives
\[
\begin{aligned}
 K_{r-1}\oplus F(V_{r+1})
 &\cong F(V_{r-1}\oplus V_{r+1})\\
 &\cong F(V_2\boxt V_r)\\
 &\cong F(V_2)\boxt F(V_r)\\
 &\cong K_2\boxt K_r
 \cong K_{r-1}\oplus K_{r+1}.
\end{aligned}
\]
Here the last isomorphism is the split case of
\eqref{eq:K2-fusion}.  All objects involved have finite length, and
$\Oc_{c_{p,q}}$ is Krull--Schmidt by Section~\ref{sec:prelim}.  Hence
Krull--Schmidt cancellation of the common summand $K_{r-1}$ gives
$F(V_{r+1})\cong K_{r+1}$.
\end{proof}

\subsection{Residue bookkeeping and transported recursions}

\begin{lem}[Residue-sector peeling]\label{lem:residue-sector-peeling}
Let all modules below have finite length.  Suppose $\xi_-\ne\xi_+$ and
\[
 A_-=A_-^{[\xi_-]},\quad B_-=B_-^{[\xi_-]},\qquad
 A_+=A_+^{[\xi_+]},\quad B_+=B_+^{[\xi_+]},
\]
and suppose
\[
 0\longrightarrow A_-\oplus A_+\longrightarrow X
 \longrightarrow B_-\oplus B_+\longrightarrow0
\]
is exact.  Then
\[
 X=X^{[\xi_-]}\oplus X^{[\xi_+]}
\]
and exact residue projection gives
\[
 0\to A_\pm\to X^{[\xi_\pm]}\to B_\pm\to0.
\]
For one fixed $\xi\in\{\xi_-,\xi_+\}$, suppose that
$X=D\oplus D'$ with $D=D^{[\xi]}$ and that $D$ has the same
Jordan--H\"older multiplicities as $X^{[\xi]}$.  Then
$D'^{[\xi]}=0$, and the inclusion $D\hookrightarrow X$ has image exactly
$X^{[\xi]}$.  Finally, if $X$ is supported in a finite subset
$\Sigma\subset\C/\Z$, then every direct summand of $X$ is supported
in $\Sigma$.
\end{lem}

\begin{proof}
Lemma~\ref{lem:block-projection} gives the two displayed exact
sequences.  If $\eta\notin\{\xi_-,\xi_+\}$, exact residue projection gives
$0\to0\to X^{[\eta]}\to0$, hence $X^{[\eta]}=0$.  This proves the asserted
two-sector decomposition.  For the second assertion, the biproduct inclusions and projections are
Virasoro-module maps, hence commute with $L_0$ and preserve generalized
conformal-weight residue sectors.  Exact residue projection of
$X=D\oplus D'$ therefore gives
\[
 X^{[\xi]}=D\oplus D'^{[\xi]}.
\]
Here and below, Jordan--H\"older multiplicities are indexed by
isomorphism classes of simple objects, not by their numerical labels.  In
particular, the possible collision $S_j\cong S_{p-j}$ for $q=2$ does not
affect the argument.  If $D$ and $X^{[\xi]}$ have the same such
multiplicities, Lemma~\ref{lem:length-cancellation} gives
$D'^{[\xi]}=0$.  Hence the
inclusion $D\hookrightarrow X$ has image exactly the residue sector
$X^{[\xi]}$.  More generally, residue projection is additive: if
$X=D_1\oplus D_2$, then
$X^{[\eta]}=D_1^{[\eta]}\oplus D_2^{[\eta]}$ for every $\eta$.  Thus every
direct summand of an object supported in a finite set $\Sigma\subset\C/\Z$
is supported in $\Sigma$.
\end{proof}

\begin{cor}[Residue-sector cancellation]
\label{cor:residue-sector-cancellation}
If $X^{[\xi]}=A\oplus B$ for finite-length modules and
$[X^{[\xi]}]=[A]$ in the Grothendieck group, then $B=0$.
\end{cor}

\begin{proof}
The equality gives $[B]=0$.  Since the Grothendieck group of a length
category is free on simple isomorphism classes, all Jordan--H"older
multiplicities of $B$ vanish and hence $B=0$.
\end{proof}

\begin{lem}[Grothendieck classes in a length category]
\label{lem:length-K0}
For finite-length objects in $\Oc_{c_{p,q}}$, equality in the
Grothendieck group is equality of Jordan--H\"older multiplicities:
\[
 [M]=\sum_{[S]}[M:S][S].
\]
In particular, if $[M]=[S]$ for a simple object $S$, then $M\cong S$.
\end{lem}

\begin{proof}
The Grothendieck group of a length category is free on the isomorphism
classes of simple objects, with coefficients given by Jordan--H\"older
multiplicities.  Thus $[M]=[S]$ forces $\ell(M)=1$ and its unique simple
factor to be $S$.
\end{proof}

\begin{lem}[Residue pairing for a strip]\label{lem:strip-residue-pairing}
Let $a=np-r$ and $b=np+r$, where $n\ge1$ and $1\le r\le p-1$.  Then
\[
 [h_a]=[h_b],
 \qquad h_b-h_a=r(nq-1)\in\Z.
\]
If $1\le r\le p-2$, then $a-1,a+1,b-1,b+1$ are all positive and
\[
 [h_{a+1}]=[h_{b-1}],\qquad
 [h_{a-1}]=[h_{b+1}]
 \quad\text{in }\C/\Z,
\]
and these two residue classes are distinct.  At the far wall put
\[
 \xi_0:=[h_{(n-1)p}]=[h_{(n+1)p}],\qquad
 \xi_1:=[h_{(n-1)p+2}]=[h_{(n+1)p-2}].
\]
Then $\xi_0\ne\xi_1$ for every $n\ge1$.  When $n=1$, the symbol $h_0$
is used only in the arithmetic sense fixed in Section~\ref{sec:prelim};
there is no module $K_0$, and any corresponding module term is omitted.
For $p=2$ the same convention applies to the absent label $0$, and one
has explicitly
\[
 [h_{2n}]\ne[h_{2n-2}]=[h_{2n+2}].
\]
\end{lem}

\begin{proof}
The first congruence is the specialization of
\eqref{eq:weight-difference} giving $h_b-h_a=r(nq-1)$.  The neighbouring and wall congruences are the corresponding formulas in
Corollary~\ref{cor:residue-separation}.  The crossed wall differences there
show $\xi_0\ne\xi_1$ also for $n=1$; the value $h_0$ is only the arithmetic
extension of the Kac weight formula.  For $p=2$, coprimality forces $q$ to
be odd and
\[
 h_{2n}-h_{2n-2}=nq-1-\frac q2\notin\Z,
\]
while $h_{2n+2}-h_{2n-2}\in\Z$.  This gives the final displayed residue
relation.  The convention about the label $0$ is bookkeeping and is never
used as a statement about a Kac module.
\end{proof}

\begin{lem}[Transported projective recursions]\label{lem:transported-recursions}
Applying the additive strong monoidal functor $F$ to the projective
tensor-product decomposition isomorphisms of McRae--Yang gives the
following identities.  If
$p\ge3$,
\begin{align}
 K_2\boxt Q_{np}&\cong Q_{np+1},\label{eq:transport-wall-first}\\
 K_2\boxt Q_{np+1}&\cong2Q_{np}\oplus Q_{np+2},\label{eq:transport-first}\\
 K_2\boxt Q_{np+r}&\cong Q_{np+r-1}\oplus Q_{np+r+1},
   &&2\le r\le p-2,\label{eq:transport-interior}\\
 K_2\boxt Q_{(n+1)p-1}&\cong
 Q_{(n-1)p}\oplus Q_{(n+1)p-2}\oplus Q_{(n+1)p}.
 \label{eq:transport-far-wall}
\end{align}
Here $n\ge1$ and the term $Q_0$ is absent.  If $p=2$, then
\begin{align}
 K_2\boxt Q_{2n}&\cong Q_{2n+1},\label{eq:transport-p2-even}\\
 K_2\boxt Q_{2n+1}&\cong
 Q_{2n-2}\oplus2Q_{2n}\oplus Q_{2n+2},
 \label{eq:transport-p2-odd}
\end{align}
again omitting $Q_0$ when $n=1$.
\end{lem}

\begin{proof}
Each source recursion used here is already an isomorphism in the additive
projective subcategory, not a nonsplit short exact sequence.  For example,
\[
 V_2\boxtimes P_{np}\cong P_{np+1}
 \quad\Longrightarrow\quad
 K_2\boxtimes Q_{np}\cong Q_{np+1};
\]
only preservation of this isomorphism is used.  For \eqref{eq:transport-wall-first}, apply the $r=p$ wall case of
\cite[Theorem~4.8]{McRaeYang} with the theorem's strip index equal to
$n-1$:
\[
 V_2\boxtimes P_{(n-1)p+p}\cong P_{np+1},
\]
that is,
\[
 V_2\boxtimes P_{np}\cong P_{np+1}
 \quad\Longrightarrow\quad
 K_2\boxtimes Q_{np}\cong Q_{np+1}.
\]
Only preservation of this isomorphism is used.  For completeness, the
four cases for $p\ge3$ are obtained from \cite[Theorem~4.8]{McRaeYang}
as follows.  The case $r=p$ with strip index $n-1$ gives
\eqref{eq:transport-wall-first}; the case $r=1$ with strip index $n$
gives \eqref{eq:transport-first}; the cases $2\le r\le p-2$ give
\eqref{eq:transport-interior} without reindexing; and the case $r=p-1$
gives \eqref{eq:transport-far-wall}.  Thus every displayed formula is
the image under $F$ of a biproduct isomorphism in the projective
subcategory.  For $p=2$, both \eqref{eq:transport-p2-even} and
\eqref{eq:transport-p2-odd} are recorded in
\cite[Theorem~1.4(1)]{McRaeYang}; equivalently, the first is the defining
identity $P_{2n+1}=V_2\boxtimes P_{2n}$ and the second is
\cite[Theorem~4.3]{McRaeYang}.  Applying the additive functor $F$ therefore preserves the biproduct decomposition,
and strong monoidality identifies
$F(V_2\boxtimes P)\cong K_2\boxtimes F(P)$.  This gives the displayed
recursions without applying $F$ to any nonsplit exact sequence.
\end{proof}

\subsection{Projective reconstruction}

\begin{prop}[Reconstruction of images of projectives]\label{prop:projective-images}
For $n\geq1$ and $1\leq r\leq p-1$,
\begin{equation}\label{eq:F-projective-filtration}
 0\longrightarrow K_{np-r}\longrightarrow Q_{np+r}
 \longrightarrow K_{np+r}\longrightarrow0.
\end{equation}
Moreover, with
\[
 \xi_{n,r}:=[h_{np-r}]=[h_{np+r}],
\]
one has the single-residue support
\begin{equation}\label{eq:F-projective-residue}
 Q_{np+r}=Q_{np+r}^{[\xi_{n,r}]}.
\end{equation}
Finally,
\begin{equation}\label{eq:F-wall-projective}
 Q_{np}\cong K_{np}\qquad(n\geq1).
\end{equation}
\end{prop}

\begin{proof}
At object level every occurrence of $K_0$ or $Q_0$ below denotes an absent
summand.  The symbol $h_0$ is retained only in congruence calculations and
never denotes the lowest weight of an object.

We first assume $p\ge3$.  We prove simultaneously, by induction over
successive strips, the Kac filtrations, the wall identifications, and the
single-residue support.  More precisely, let $\mathcal J_n$ be the conjunction of the wall
identifications
\[
 Q_{mp}\cong K_{mp}\qquad(1\le m\le n)
\]
and, for every completed strip $1\le m\le n-1$ and $1\le r<p$, the
statements
\[
 0\to K_{mp-r}\to Q_{mp+r}\to K_{mp+r}\to0,
 \qquad
 Q_{mp+r}=Q_{mp+r}^{[h_{mp+r}]}.
\]
Lemma~\ref{lem:F-small-Weyl} gives $Q_j\cong K_j$ for $1\le j\le p$, hence
the base case $\mathcal J_1$.  Suppose $\mathcal J_n$ holds.  The proof of
$\mathcal J_n\Rightarrow\mathcal J_{n+1}$ is a second, inner induction
along the $n$-th strip.  We first construct $Q_{np+1}$ and then successively
construct
\[
 Q_{np+2},\ldots,Q_{(n+1)p-1}.
\]
Thus, at the step from horizontal index $r$ to $r+1$, every vertex strictly
to the left in the same strip has already been reconstructed.  The final
far-wall recursion then determines $Q_{(n+1)p}$.

\smallskip
\noindent\emph{First vertex.}
Using the fixed isomorphism $Q_{np}\cong K_{np}$, form the target-side
isomorphism
\[
 \alpha_n:
 K_2\boxtimes K_{np}
 \xrightarrow{\ \sim\ }
 K_2\boxtimes Q_{np}
 \xrightarrow{\ \sim\ }
 Q_{np+1},
\]
where the second arrow is \eqref{eq:transport-wall-first}.  Transporting
the known Virasoro wall-fusion sequence
\[
 0\longrightarrow K_{np-1}\longrightarrow
 K_2\boxtimes K_{np}\longrightarrow K_{np+1}\longrightarrow0
\]
across $\alpha_n$ gives
\[
 0\longrightarrow K_{np-1}\longrightarrow Q_{np+1}
 \longrightarrow K_{np+1}\longrightarrow0.
\]
No source short exact sequence is used in this step.
Since $h_{np+1}-h_{np-1}=nq-1\in\Z$, both end terms are supported in the
same residue class.  For every $\eta\ne[h_{np+1}]$, exact residue
projection gives
\[
 0\longrightarrow0\longrightarrow Q_{np+1}^{[\eta]}\longrightarrow0,
\]
so $Q_{np+1}^{[\eta]}=0$.  Hence
\[
 Q_{np+1}=Q_{np+1}^{[h_{np+1}]},
\]
which establishes both parts of the induction invariant at the first
vertex.

\smallskip
\noindent\emph{Propagation through the strip.}
Assume the sequence is known for $Q_b$, where
$b=np+r$, $1\le r\le p-2$, and put $a=np-r$.  Thus the inner induction
hypothesis provides the already established target-side exact sequence
\[
 0\longrightarrow K_a\longrightarrow Q_b\longrightarrow K_b
 \longrightarrow0.
\]
Since neither $a$ nor $b$ is divisible by $p$, both relevant $K_2$-fusion
sequences are in the split non-wall case.  More explicitly,
\[
 K_2\boxtimes K_a\cong K_{a-1}\oplus K_{a+1},\qquad
 K_2\boxtimes K_b\cong K_{b-1}\oplus K_{b+1}.
\]
Exactness of the target functor $K_2\boxt-$ therefore gives
\begin{equation}\label{eq:tensored-Q-sequence}
 0\longrightarrow K_{a-1}\oplus K_{a+1}
 \longrightarrow K_2\boxt Q_b
 \longrightarrow K_{b-1}\oplus K_{b+1}\longrightarrow0.
\end{equation}
By Lemmas~\ref{lem:Kac-residue-support} and
\ref{lem:strip-residue-pairing}, put
\[
 \xi_-=[h_{a+1}]=[h_{b-1}],\qquad
 \xi_+=[h_{a-1}]=[h_{b+1}],\qquad \xi_-\ne\xi_+.
\]
The two residue summands are
\begin{align}
 0&\to K_{a+1}\to (K_2\boxt Q_b)^{[\xi_-]}
       \to K_{b-1}\to0,\label{eq:interior-block-minus}\\
 0&\to K_{a-1}\to (K_2\boxt Q_b)^{[\xi_+]}
       \to K_{b+1}\to0,\label{eq:interior-block-plus}
\end{align}
with $\xi_-\ne\xi_+$.  This is a target-side exact sequence obtained by
exact tensoring and residue projection; $F$ is not being applied to the
sequence defining $Q_b$.

For $r=1$, the transported recursion is
\[
 K_2\boxt Q_{np+1}\cong2K_{np}\oplus Q_{np+2}.
\]
Since $K_{np}=K_{np}^{[\xi_-]}$ and $\xi_-\ne\xi_+$, projection of the
transported direct-sum decomposition gives
\[
 (K_2\boxt Q_{np+1})^{[\xi_-]}
 \cong2K_{np}\oplus Q_{np+2}^{[\xi_-]}.
\]
On the other hand, \eqref{eq:interior-block-minus} is
\[
 0\to K_{np}\to (K_2\boxt Q_{np+1})^{[\xi_-]}
   \to K_{np}\to0,
\]
so its Jordan--H\"older class is $2[K_{np}]$.  The summand $2K_{np}$ has exactly the Jordan--H\"older multiplicities
of the $\xi_-$ sector, so Lemma~\ref{lem:residue-sector-peeling}, applied
to the transported biproduct decomposition, gives
$Q_{np+2}^{[\xi_-]}=0$.  Since the whole tensor product is supported in
$\{\xi_-,\xi_+\}$ and $Q_{np+2}$ is a direct summand, the final assertion
of Lemma~\ref{lem:residue-sector-peeling} excludes every other residue.
Consequently
\[
 Q_{np+2}=Q_{np+2}^{[\xi_+]}.
\]
Projecting the transported biproduct decomposition to $\xi_+$ therefore
gives
\[
 (K_2\boxt Q_{np+1})^{[\xi_+]}
 \cong Q_{np+2},
\]
because the summand $2K_{np}$ is supported entirely in $\xi_-$.  Under
this identification, \eqref{eq:interior-block-plus} becomes
\[
 0\to K_{np-2}\to Q_{np+2}\to K_{np+2}\to0.
\]
Thus the new summand $Q_{np+2}$ has no $\xi_-$-component.

Now let $2\le r\le p-2$.  Here
\[
 b-1=np+(r-1),\qquad a+1=np-(r-1),
\]
so the inner induction hypothesis says that the already reconstructed
summand $Q_{b-1}$ is
supported in $\xi_-$ and fits into exactly
\[
 0\to K_{a+1}\to Q_{b-1}\to K_{b-1}\to0.
\]
Thus it has exactly the Jordan--H\"older class of
\eqref{eq:interior-block-minus}.  Projecting
\eqref{eq:transport-interior} gives
\[
 (K_2\boxt Q_b)^{[\xi_-]}
 \cong Q_{b-1}\oplus Q_{b+1}^{[\xi_-]}.
\]
The $\xi_-$-sector and the direct summand $Q_{b-1}$ have the same
Grothendieck class.  Since
\[
 (K_2\boxt Q_b)^{[\xi_-]}
 \cong Q_{b-1}\oplus Q_{b+1}^{[\xi_-]},
\]
Corollary~\ref{cor:residue-sector-cancellation} gives
$Q_{b+1}^{[\xi_-]}=0$.  As $Q_{b+1}$ is a direct summand of a tensor
product supported in $\{\xi_-,\xi_+\}$, it has no other residue support.
Thus $Q_{b+1}=Q_{b+1}^{[\xi_+]}$.  Moreover, the inner induction gives
$Q_{b-1}=Q_{b-1}^{[\xi_-]}$, so $Q_{b-1}^{[\xi_+]}=0$.  Projecting
\eqref{eq:transport-interior} to $\xi_+$ therefore yields
\[
 (K_2\boxt Q_b)^{[\xi_+]}
 \cong Q_{b+1}^{[\xi_+]}
 =Q_{b+1}.
\]
Under this identification, \eqref{eq:interior-block-plus} becomes
\[
 0\to K_{a-1}\to Q_{b+1}\to K_{b+1}\to0.
\]
This propagates both the filtration and the single-residue support through
the strip.  When $p=3$ the range
$2\le r\le p-2$ is empty, so the preceding $r=1$ step already reaches the
far edge: in that case $Q_{np+2}=Q_{(n+1)p-1}$.

\smallskip
\noindent\emph{Far edge and next wall.}
After $Q_{(n+1)p-1}$ has been reconstructed, the completed inner induction
provides the target-side exact sequence
\[
 0\longrightarrow K_{(n-1)p+1}
 \longrightarrow Q_{(n+1)p-1}
 \longrightarrow K_{(n+1)p-1}
 \longrightarrow0.
\]
Tensoring this already established target sequence exactly with $K_2$ and
using the split non-wall $K_2$-fusion decompositions on its end terms, set
\[
 M:=K_2\boxt Q_{(n+1)p-1}.
\]
Put
\[
 \xi_0=[h_{(n-1)p}]=[h_{(n+1)p}],\qquad
 \xi_1=[h_{(n-1)p+2}]=[h_{(n+1)p-2}].
\]
Lemma~\ref{lem:strip-residue-pairing} gives $\xi_0\ne\xi_1$ for every
$n\ge1$.  Exact tensoring and residue projection give, for $n>1$,
\begin{align}
 0&\to K_{(n-1)p}\to M^{[\xi_0]}
      \to K_{(n+1)p}\to0,\label{eq:far-wall-sector-zero}\\
 0&\to K_{(n-1)p+2}\to M^{[\xi_1]}
      \to K_{(n+1)p-2}\to0.\label{eq:far-wall-sector-one}
\end{align}
For $n=1$ the $\xi_0$ sector is instead the isomorphism
\begin{equation}\label{eq:far-wall-sector-zero-base}
 M^{[\xi_0]}\cong K_{2p},
\end{equation}
while \eqref{eq:far-wall-sector-one} remains valid.  In every case
$M^{[\eta]}=0$ for $\eta\notin\{\xi_0,\xi_1\}$.  On the other hand,
\eqref{eq:transport-far-wall} gives
\[
 M\cong Q_{(n-1)p}\oplus Q_{(n+1)p-2}\oplus Q_{(n+1)p},
\]
with $Q_0$ omitted.  By the outer induction hypothesis,
$Q_{(n-1)p}=K_{(n-1)p}$ is supported in $\xi_0$ when $n>1$, whereas the
inner induction has already reconstructed $Q_{(n+1)p-2}$ and placed it in
$\xi_1$.  Projection of the direct-sum
recursion to $\xi_1$ therefore gives
\[
 M^{[\xi_1]}
 \cong Q_{(n+1)p-2}\oplus Q_{(n+1)p}^{[\xi_1]}.
\]
Indeed, the inner induction gives
\[
 0\to K_{(n-1)p+2}\to Q_{(n+1)p-2}
 \to K_{(n+1)p-2}\to0,
\]
so this summand has exactly the Jordan--H\"older multiplicities of
\eqref{eq:far-wall-sector-one}.  Lemma~\ref{lem:residue-sector-peeling}, with
$D=Q_{(n+1)p-2}$ in the $\xi_1$ sector, therefore yields
\[
 Q_{(n+1)p}^{[\xi_1]}=0.
\]
Since no other residue occurs in $M$, the new wall summand is supported in
$\xi_0$; explicitly,
\[
 Q_{(n+1)p}=Q_{(n+1)p}^{[\xi_0]}.
\]

If $n>1$, projection to $\xi_0$ gives
\[
 M^{[\xi_0]}\cong K_{(n-1)p}\oplus Q_{(n+1)p}.
\]
In the Grothendieck group, comparison with
\eqref{eq:far-wall-sector-zero} gives
\[
 [M^{[\xi_0]}]=[K_{(n-1)p}]+[K_{(n+1)p}]
 =[K_{(n-1)p}]+[Q_{(n+1)p}],
\]
and therefore $[Q_{(n+1)p}]=[K_{(n+1)p}]$.  Applying the length
homomorphism $\ell:K_0(\Oc_{c_{p,q}})\to\Z$ gives
$\ell(Q_{(n+1)p})=1$, because $K_{(n+1)p}=S_{(n+1)p}$ is simple.  Hence
$Q_{(n+1)p}$ is simple.  Equality of Grothendieck classes identifies its
unique composition factor with $S_{(n+1)p}$, and therefore
\[
 Q_{(n+1)p}\cong K_{(n+1)p}.
\]
When $n=1$, the terms $K_0,Q_0$ do not occur.  Equation
\eqref{eq:far-wall-sector-zero-base} gives $M^{[\xi_0]}\cong K_{2p}$.  On the
other hand, the transported direct-sum decomposition and the fact that
$Q_{2p-2}$ is supported in $\xi_1$ give
$M^{[\xi_0]}\cong Q_{2p}$.  Hence $Q_{2p}\cong K_{2p}$.  Thus
$\mathcal J_{n+1}$ holds.

For $p=2$ we use the following induction invariant $\mathcal I_n$.
For every $1\le m\le n$ one has $Q_{2m}\cong K_{2m}$, and for every
$1\le m\le n-1$ the odd vertex satisfies
\[
 0\to K_{2m-1}\to Q_{2m+1}\to K_{2m+1}\to0,
 \qquad Q_{2m+1}=Q_{2m+1}^{[h_{2m+1}]}.
\]
The base case $\mathcal I_1$ is $Q_2\cong K_2$, supplied by
Lemma~\ref{lem:F-small-Weyl}.  Assume $\mathcal I_n$.  Equation
\eqref{eq:transport-p2-even} identifies
\[
 Q_{2n+1}\cong K_2\boxt K_{2n}.
\]
The wall fusion sequence therefore gives
\[
 0\to K_{2n-1}\to Q_{2n+1}\to K_{2n+1}\to0.
\]
The two end terms have congruent lowest weights.  Exact projection to every
other residue sector is zero, so
\[
 Q_{2n+1}=Q_{2n+1}^{[h_{2n+1}]}.
\]
This establishes the required filtration and single-residue support for the
new odd vertex.

Set $M:=K_2\boxt Q_{2n+1}$.  Exact tensoring of the displayed filtration,
using the split non-wall fusion sequences on both end terms, gives
\[
 0\to K_{2n-2}\oplus K_{2n}\to M
 \to K_{2n}\oplus K_{2n+2}\to0,
\]
with $K_0$ omitted for $n=1$.  Put
\[
 \xi_{\mathrm w}:=[h_{2n}],\qquad
 \xi_{\mathrm o}:=[h_{2n-2}]=[h_{2n+2}].
\]
By Lemma~\ref{lem:strip-residue-pairing},
\[
 h_{2n}-h_{2n-2}=nq-1-\frac q2\notin\Z,
\]
so the residues $\xi_{\mathrm w}$ and $\xi_{\mathrm o}$ are distinct.  Exact residue projection
gives
\begin{align*}
 0&\to K_{2n}\to M^{[\xi_{\mathrm w}]}
     \to K_{2n}\to0,\\
 0&\to K_{2n-2}\to M^{[\xi_{\mathrm o}]}
     \to K_{2n+2}\to0,
\end{align*}
with the first term of the second sequence omitted when $n=1$.
Meanwhile \eqref{eq:transport-p2-odd} gives
\[
 M\cong Q_{2n-2}\oplus2Q_{2n}\oplus Q_{2n+2},
\]
where $Q_0$ is absent.  If $n>1$, the induction hypothesis gives
\[
 Q_{2n-2}=K_{2n-2}=K_{2n-2}^{[\xi_{\mathrm o}]},
\]
whereas for $n=1$ the summand $Q_0$ is absent.  Hence in either case
$Q_{2n-2}^{[\xi_{\mathrm w}]}=0$.  Since
$2Q_{2n}=2K_{2n}$ is supported entirely in $\xi_{\mathrm w}$,
projection of \eqref{eq:transport-p2-odd} gives
\[
 M^{[\xi_{\mathrm w}]}
 \cong2K_{2n}\oplus Q_{2n+2}^{[\xi_{\mathrm w}]}.
\]
The wall-sector exact sequence has Jordan--H\"older class
$2[K_{2n}]$.  Lemma~\ref{lem:residue-sector-peeling}, applied to the
summand $2K_{2n}$, gives $Q_{2n+2}^{[\xi_{\mathrm w}]}=0$.  Hence the new even summand is entirely
supported in $\xi_{\mathrm o}$.  If $n>1$, the already known wall
$Q_{2n-2}=K_{2n-2}$ is also supported in $\xi_{\mathrm o}$, and projection
of \eqref{eq:transport-p2-odd} gives
\[
 M^{[\xi_{\mathrm o}]}
 \cong K_{2n-2}\oplus Q_{2n+2}.
\]
Comparison with the $\xi_{\mathrm o}$ exact sequence therefore gives
$[Q_{2n+2}]=[K_{2n+2}]$.  Since Grothendieck classes record
Jordan--H\"older multiplicities and $K_{2n+2}$ is simple, this equality
forces $Q_{2n+2}$ to have length one; hence
$Q_{2n+2}\cong K_{2n+2}$.  For $n=1$, the $Q_0$ summand is absent, so
\[
 M\cong2K_2\oplus Q_4.
\]
Since $2K_2$ is supported in $\xi_{\mathrm w}$, projection to
$\xi_{\mathrm o}$ gives
\[
 M^{[\xi_{\mathrm o}]}\cong Q_4.
\]
On the other hand, the second residue-sector sequence, with its $K_0$
term omitted, gives
\[
 M^{[\xi_{\mathrm o}]}\cong K_4.
\]
Hence $Q_4\cong K_4$.  Thus the new odd filtration and the wall identification together establish
$\mathcal I_{n+1}$, completing the induction.

\end{proof}

\begin{cor}[Composition factors]\label{cor:Q-composition-factors}
For $n\ge1$ and $1\le r\le p-1$,
\begin{equation}\label{eq:Q-global-JH}
 [Q_{np+r}]
 =[S_{np-r}]+2[S_{np+r}]+[S_{(n+2)p-r}].
\end{equation}
\end{cor}

\begin{proof}
Add the Jordan--H\"older classes of the two Kac modules in
\eqref{eq:F-projective-filtration}.  The Kac sequence for
$K_{np-r}=K_{(n-1)p+(p-r)}$ has simple factors
$S_{np-r}$ and $S_{np+r}$, while that for $K_{np+r}$ has factors
$S_{np+r}$ and $S_{(n+2)p-r}$.
\end{proof}

\begin{rmk}\label{rmk:no-left-exactness}
Until projective faithfulness is established, the reconstruction uses only
the additive strong monoidal structure of $F$: $F$ is applied to
isomorphisms, finite biproduct decompositions, and morphisms.  The target
short exact sequences are either the known $K_2$-fusion sequences or are
obtained from previously established target sequences by exact tensoring
with $K_2$ and exact conformal-residue projection.  No nonsplit short exact
sequence in the source is transported through $F$.
\end{rmk}

\subsection{Faithfulness on projectives}

\begin{lem}[Temperley--Lieb ideal detection]\label{lem:TL-ideal-detection}
Let $u\in\C^\times$ satisfy $\operatorname{ord}(u^2)=p\ge3$ and set
$\delta=u+u^{-1}$.  With
\[
 [m]_u:=\frac{u^m-u^{-m}}{u-u^{-1}},\qquad
 \ell:=\min\{m\ge1:[m]_u=0\},
\]
one has $\ell=p$.  Every nonzero proper tensor ideal of $\TL(\delta)$
contains the first-critical one-row path idempotent
$p_{[p-1]}(u)$ of Goodman--Wenzl; this idempotent is evaluable and is
the $(p-1)$-strand Jones--Wenzl projector $JW_{p-1}(u)$.
\end{lem}

\begin{proof}
Since $p\ge3$, one has $u\ne\pm1$, so the denominator in $[m]_u$ is
nonzero.  Thus
\[
 [m]_u=0\quad\Longleftrightarrow\quad u^{2m}=1,
\]
and $\operatorname{ord}(u^2)=p$ gives
\[
 \ell=p,\qquad [p]_u=0,\qquad [j]_u\ne0\quad(1\le j<p).
\]
Since $u^2$ is a primitive $p$-th root of unity, this is exactly the
root-of-unity specialization of \cite[Section~2.4]{GoodmanWenzl}, with
their critical parameter $\ell=p$.  Proposition~2.1 of
\cite{GoodmanWenzl} identifies the negligible ideal with the tensor
ideal generated by the first critical path idempotent
$p_{[p-1]}(u)$, and Theorem~3.3 identifies it as the unique nonzero
proper tensor ideal in this specialization.

The sign convention used for quantum $\mathfrak{sl}_2$ is obtained by
putting
\[
 \mathbb v:=-u.
\]
Then
\[
 u+u^{-1}=-\mathbb v-\mathbb v^{-1}=-[2]_{\mathbb v},
 \qquad
 [m]_{\mathbb v}=(-1)^{m-1}[m]_u.
\]
Hence the Goodman--Wenzl category and the Temperley--Lieb category in
the tilting convention have exactly the same loop parameter and the
same first critical index.  Because $[j]_u\ne0$ for $1\le j<p$, the
one-row path reaches the first critical line only at its endpoint, and
the corresponding path idempotent is evaluable.  The usual recursive
Jones--Wenzl formula has no vanishing denominator before that endpoint,
so
\[
 p_{[p-1]}(u)=JW_{p-1}(u).
\]
Equivalently, both are the unique evaluable idempotent with coefficient
one on the identity diagram and satisfying the Jones--Wenzl absorption
relations.  Therefore every nonzero proper tensor ideal contains
$JW_{p-1}(u)$.
\end{proof}



\begin{lem}[First-critical projector under the tilting realization]
\label{lem:first-critical-tilting-projector}
Let
\[
 \zeta=e^{\pi iq/p},\qquad u=-\zeta,\qquad
 \delta=u+u^{-1}=-\zeta-\zeta^{-1},
\]
with $p\ge3$, and let
\[
 \iota:\TL(\delta)\xrightarrow{\sim}\TL_{\C,\zeta}
\]
be the identity-on-diagrams monoidal identification.  Then
\[
 \iota\bigl(JW_{p-1}(u)\bigr)=E_{p-1}^{(\infty,p)},
\]
where $E_{p-1}^{(\infty,p)}$ is the mixed Jones--Wenzl idempotent of
\cite{SuttonTubbenhauerWedrichZhu}.  Under the
Temperley--Lieb/tilting equivalence its retract is the indecomposable
tilting module $T(p-1)$.
\end{lem}

\begin{proof}
Since $u^2=\zeta^2=e^{2\pi iq/p}$ is a primitive $p$-th root of unity,
\[
 [p]_\zeta=0,\qquad [j]_\zeta\ne0\quad(1\le j<p),
\]
so the complex specialization has mixed characteristic
$(\infty,p)$.  The integer $p=[1,0]_{\infty,p}$ is an eve in the sense
of \cite[Definition~2.5]{SuttonTubbenhauerWedrichZhu}.  For an eve the
mixed projector at the first critical position specializes to the
ordinary one-row Jones--Wenzl projector.  The comparison of diagrammatic
conventions in \cite[Section~1.3(b)]{SuttonTubbenhauerWedrichZhu}
states that, over $\C$ at a root of unity, these mixed projectors agree
with the Goodman--Wenzl projectors after matching circle values.  Our
choice $u=-\zeta$ gives
\[
 u+u^{-1}=-\zeta-\zeta^{-1},
\]
so the identification $\iota$ matches the two conventions and hence
\[
 \iota\bigl(JW_{p-1}(u)\bigr)=E_{p-1}^{(\infty,p)}.
\]
Finally, \cite[Theorem~3.18]{SuttonTubbenhauerWedrichZhu} identifies the
retract defined by the mixed projector $E_{v-1}$ with $T(v-1)$.
Taking $v=p$ gives the claimed retract $T(p-1)$.
\end{proof}

\begin{prop}\label{prop:F-faithful-projectives-large}
If $p\geq3$, the restriction of $F$ to the projective subcategory is
faithful.
\end{prop}

\begin{proof}
Put $\zeta=e^{\pi iq/p}$.  By \cite[Proposition~6.4 and Theorem~6.6]{McRaeYang} there is a
monoidal equivalence
\begin{equation}\label{eq:tilting-projective-equivalence}
 \Phi:\mathcal T_\zeta\xrightarrow{\sim}\mathcal P^k,
 \qquad \Phi(X)=V_2,\qquad \Phi(T_\lambda)=P_{\lambda+1}.
\end{equation}
Let $G:=F\circ\Phi$.  Suppose that $G$ is not faithful.  Choose
$0\ne f_0:A\to B$ with $G(f_0)=0$.  After decomposing $A$ and $B$ into indecomposable tiltings, choose
summands $T$ of $A$ and $U$ of $B$, and write
\[
 \iota_T^A:T\longrightarrow A,
 \qquad
 \pi_U^B:B\longrightarrow U
\]
for the corresponding inclusion and projection.  Choose $T,U$ so that
\[
 f:=\pi_U^B\circ f_0\circ\iota_T^A:T\longrightarrow U
\]
is nonzero.  Functoriality gives $G(f)=0$.  By \cite[Remark~2.11]{SuttonTubbenhauerWedrichZhu}, every
indecomposable tilting is a direct summand of a tensor power of the
fundamental tilting module $X=T(1)$.
Choose retract data
\[
 T\xrightarrow{j_T}X^{\otimes m}\xrightarrow{r_T}T,
 \qquad
 U\xrightarrow{j_U}X^{\otimes n}\xrightarrow{r_U}U,
 \qquad r_Tj_T=\id_T,\quad r_Uj_U=\id_U.
\]
Then
\[
 r_U\circ(j_U\circ f\circ r_T)\circ j_T=f\ne0,
\]
so $j_U\circ f\circ r_T\ne0$.  Functoriality also gives
$G(j_U\circ f\circ r_T)=0$.  Set
\[
 u:=-\zeta=-e^{\pi iq/p},\qquad
 \delta:=u+u^{-1}=-\zeta-\zeta^{-1},
 \qquad \TL_\delta:=\TL(\delta).
\]
Since $\gcd(p,q)=1$,
\[
 \operatorname{ord}(u^2)=p,\qquad
 [p]_u=0,\qquad [j]_u\ne0\quad(1\le j<p).
\]
Thus the first critical index is $p$, and $JW_{p-1}(u)$ is evaluable
because its recursion involves only the nonzero denominators
$[1]_u,\ldots,[p-1]_u$.  Goodman--Wenzl use the root parameter $t=u$.  In the convention of
\cite{SuttonTubbenhauerWedrichZhu}, the specialized Temperley--Lieb
category $\TL_{\C,\zeta}$ has circle value
\[
 -[2]_\zeta=-\zeta-\zeta^{-1}=\delta.
\]
Thus the identity on Temperley--Lieb diagrams gives a strict
$\C$-linear monoidal identification
\[
 \iota:\TL_\delta\xrightarrow{\sim}\TL_{\C,\zeta}.
\]
Moreover
\[
 [m]_{-\zeta}=(-1)^{m-1}[m]_\zeta,
 \qquad
 [m]_u=0\iff u^{2m}=1.
\]
Since $\operatorname{ord}(u^2)=\operatorname{ord}(\zeta^2)=p$, the
Goodman--Wenzl critical index and the quantum characteristic of the
complex specialization in \cite{SuttonTubbenhauerWedrichZhu} are both
$\ell=p$, with characteristic parameter $\mathsf p=\infty$.  Composing
$\iota$ with the functor of
\cite[Proposition~2.13]{SuttonTubbenhauerWedrichZhu} gives
\begin{equation}\label{eq:TL-fullness}
 \Psi:\TL_\delta\longrightarrow\mathcal T_\zeta,
 \qquad1\longmapsto X,
\end{equation}
which is full on Hom spaces between tensor powers and induces an
equivalence after additive idempotent completion; see also
\cite[Remark~2.15]{SuttonTubbenhauerWedrichZhu}.  For every $m,n\ge0$, set
\[
 J(m,n):=\ker\!\left(
 \Hom_{\TL_\delta}(m,n)\xrightarrow{\,G\circ\Psi\,}
 \Hom(G\Psi(m),G\Psi(n))\right).
\]
Because $G\circ\Psi$ is $\C$-linear and strong monoidal, the spaces
$J(m,n)$ form a $\C$-linear family stable under arbitrary pre- and
postcomposition and under tensoring on either side; hence $J$ is a tensor
ideal.  By fullness choose $\widetilde f$
with $\Psi(\widetilde f)=j_U\circ f\circ r_T$.  Since the latter morphism
is nonzero, necessarily $\widetilde f\ne0$; thus $\widetilde f\in J(m,n)$
and $J\ne0$.
The ideal is proper because
\[
 (G\circ\Psi)(\id_{\mathbf1})
 =\id_{(G\circ\Psi)(\mathbf1)}\ne0;
\]
in fact $(G\circ\Psi)(\mathbf1)\cong K_1$.

Since $\operatorname{ord}(u^2)=p$ by $\gcd(p,q)=1$,
Lemma~\ref{lem:TL-ideal-detection} applies.  Equivalently, by
\cite[Theorem~3.3]{GoodmanWenzl}, the nonzero proper tensor ideal $J$
is the negligible ideal, and \cite[Proposition~2.1]{GoodmanWenzl}
therefore gives
\[
 JW_{p-1}(u)=p_{[p-1]}(u)\in J.
\]
By Lemma~\ref{lem:first-critical-tilting-projector}, the retract of
$\Psi(JW_{p-1}(u))$ is precisely the indecomposable tilting module
\[
 T(p-1)=T_{p-1}.
\]
Thus, setting
\[
 e:=\Psi(JW_{p-1}(u))\in
 \End_{\mathcal T_\zeta}(X^{\otimes(p-1)}),
\]
the retract defined by $e$ is $T_{p-1}$.  Choose maps
\[
 T_{p-1}\xrightarrow{i}X^{\otimes(p-1)}
 \xrightarrow{\pi}T_{p-1},
 \qquad \pi i=\id_{T_{p-1}},\qquad i\pi=e.
\]
Since $JW_{p-1}(u)\in J$, one has
$(G\circ\Psi)(JW_{p-1}(u))=G(e)=0$.
Because $e=i\pi$, one has $e\circ i=i$, and hence
\[
 G(i)=G(e)\circ G(i)=0.
\]
Therefore
\[
 \id_{G(T_{p-1})}=G(\pi)\circ G(i)=0.
\]
Thus $G(T_{p-1})$ is a zero object.  Under \eqref{eq:tilting-projective-equivalence},
$T_{p-1}$ corresponds to $P_p$, so this says $F(P_p)=0$, contradicting
Lemma~\ref{lem:F-small-Weyl}, which gives $F(P_p)=K_p\ne0$.
\end{proof}

\begin{lem}\label{lem:p-two-faithful}
If $p=2$, then $F|_{\mathcal P^k}$ is faithful.
\end{lem}

\begin{proof}
The wall projectives $P_{2n}$ are isolated and their identities have
nonzero images by Proposition~\ref{prop:projective-images}.  For $p=2$ the non-wall reflection chain is $s_i=2i+1$, so the
corresponding projectives are
\[
 P_1,P_3,P_5,\ldots.
\]
For $m\ge1$, \eqref{eq:transport-p2-even} and
Proposition~\ref{prop:projective-images} give
\[
 Q_{2m+1}\cong K_2\boxt Q_{2m}\cong K_2\boxt K_{2m}.
\]
Since $2\mid2m$, the wall-fusion result
\cite[Theorem~6.5 and Remark~6.11]{McRaeSopin} says that this Virasoro
module is logarithmic.  This is an external target-side fusion input and
does not use Proposition~\ref{prop:Q-logarithmic}; hence there is no
circularity.  Therefore $e^{2\pi iL_0}|_{Q_{2m+1}}$ is not diagonalizable
by Lemma~\ref{lem:exponential-logarithmic}.  By twist preservation,
\[
 F(\theta^-_{P_{2m+1}})=\theta_{Q_{2m+1}}
 =e^{2\pi iL_0}|_{Q_{2m+1}}.
\]
Together with \eqref{eq:source-twist-nilpotent}, this gives
\[
 e^{2\pi iL_0}|_{Q_{2m+1}}
 =\lambda_{2m+1}
  \bigl(\id+\gamma_{2m+1}F(N_{2m+1})\bigr),
 \qquad \gamma_{2m+1}\ne0.
\]
Thus $F(N_{2m+1})\ne0$.

Choose adjacent arrows on the source reflection chain as in
Proposition~\ref{prop:source-Hom-package}, with
$\nu_{i+1}=x_i\circ y_i$ a nonzero generator of the radical of the positive
vertex.  Since $\nu_{i+1}$ is a nonzero scalar multiple of the chosen
$N_{s_{i+1}}$,
\[
 0\ne F(\nu_{i+1})=F(x_i)F(y_i).
\]
Consequently both adjacent arrows have nonzero image.  The source Hom
package now shows injectivity on every one-dimensional adjacent Hom space.
On a positive vertex, $\id$ and $F(N)$ are linearly independent: if
$F(N)=c\id$, then $0=F(N)^2=c^2\id$, so $c=0$, contradicting
$F(N)\ne0$.  At the initial vertex the endomorphism space is already
one-dimensional.  These are all nonzero Hom spaces between indecomposable projectives when
$p=2$, by Proposition~\ref{prop:source-Hom-package}.  Thus $F$ is injective on every indecomposable projective Hom space.  A
morphism between finite direct sums is a matrix of component morphisms, so
componentwise injectivity implies faithfulness on the whole projective
subcategory.
\end{proof}

\begin{cor}[Projective faithfulness]\label{cor:F-faithful-projectives}
For every $p\geq2$, the restriction $F|_{\mathcal P^k}$ is faithful.
\end{cor}

\begin{proof}
Combine Proposition~\ref{prop:F-faithful-projectives-large} and
Lemma~\ref{lem:p-two-faithful}.
\end{proof}

\begin{prop}[Logarithmicity of non-wall images]\label{prop:Q-logarithmic}
For $n\ge1$ and $1\le r\le p-1$, the module $Q_{np+r}$ is logarithmic.
Consequently the sequence \eqref{eq:F-projective-filtration} is non-split.
\end{prop}

\begin{proof}
Projective faithfulness gives $F(N_{np+r})\ne0$.  Twist preservation and
\eqref{eq:source-twist-nilpotent} give
\[
 e^{2\pi iL_0}|_{Q_{np+r}}
 =\lambda_{np+r}
 \bigl(\id+\gamma_{np+r}F(N_{np+r})\bigr),
 \qquad \gamma_{np+r}\ne0.
\]
Moreover
\[
 F(N_{np+r})^2=0,\qquad F(N_{np+r})\ne0.
\]
Since $\gamma_{np+r}\ne0$, the minimal polynomial of
$\id+\gamma_{np+r}F(N_{np+r})$ contains $(X-1)^2$.  Hence this operator,
and therefore the exponential above, is not diagonalizable.  Lemma~\ref{lem:exponential-logarithmic} implies that
$L_0$ is not semisimple.  Both end terms of
\eqref{eq:F-projective-filtration} are ordinary; hence a split extension
would be ordinary, a contradiction.
\end{proof}

\begin{thm}[Projective reconstruction and faithfulness]\label{thm:projective-control}
For every $n\ge1$,
\[
 F(P_{np})\cong K_{np}.
\]
For every $n\ge1$ and $1\le r\le p-1$,
\[
 0\to K_{np-r}\to F(P_{np+r})\to K_{np+r}\to0
\]
is non-split, its middle term is logarithmic, and it satisfies
\[
 F(P_{np+r})=F(P_{np+r})^{[h_{np+r}]},
\qquad
 [F(P_{np+r})]
 =[S_{np-r}]+2[S_{np+r}]+[S_{(n+2)p-r}].
\]
Moreover $F|_{\mathcal P^k}$ is faithful.  No preservation by $F$ of a
nonsplit short exact sequence is used in these conclusions, nor do they
invoke Nakano's extension classification, the Hom calculations of
Section~\ref{sec:fullness}, or Drinfeld--Sokolov reduction.
\end{thm}

\begin{proof}
Combine Proposition~\ref{prop:projective-images},
Corollary~\ref{cor:Q-composition-factors},
Corollary~\ref{cor:F-faithful-projectives}, and
Proposition~\ref{prop:Q-logarithmic}.
\end{proof}


\section{Logarithmic images of affine projectives and full faithfulness}\label{sec:fullness}

The preceding section gives the Kac filtrations and projective faithfulness
without assuming exactness of $F$.  Here $P(\tau)$ denotes the Virasoro module introduced above as the middle
term of a nonzero class in
\[
 \Ext^1_{\mathscr C}\bigl(K(\tau),L(h_{\alpha_2})\bigr);
\]
Theorem~\ref{thm:Nakano-input}, applied through the one-row interface
of Corollary~\ref{cor:Nakano-one-row}, identifies this isomorphism class
with the module denoted $P(\tau)$ in
\cite[Definition~5.13]{Nakano}.  The notation does not assert
projectivity in $\Oc_{c_{p,q}}$.  No Ext group below is obtained by a
blanket identification of $\mathscr C$ with the ambient category of
\cite[Definition~3.14]{Nakano}; the only staggered input used here is the
one-row result Proposition~\ref{prop:one-row-staggered-quotient} and its
Yoneda consequence Corollary~\ref{cor:Nakano-one-row}.  This section proves the
following structural statement.

\begin{thm}[Projective images and full faithfulness]\label{thm:projective-classification}
For every indecomposable projective $P_j$,
\[
 F(P_j)\cong
 \begin{cases}
  K_j,&1\le j\le p-1,\\
  K_j=S_j,&p\mid j,\\
  P(\tau_{n,r}),&j=np+r,\ n\ge1,\ 1\le r\le p-1.
 \end{cases}
\]
Moreover, as $\C$-algebras,
\[
 \End_{\Oc_{c_{p,q}}} F(P_j)\cong
 \begin{cases}
  \C,&1\le j\le p-1\text{ or }p\mid j,\\
  \C[\varepsilon]/(\varepsilon^2),
      &j=np+r,\ n\ge1,\ 1\le r\le p-1.
 \end{cases}
\]
The modules in the third case are logarithmic and indecomposable, and
$F|_{\mathcal P^k}$ is fully faithful.
\end{thm}

The proof occupies the remainder of the section.  The dependence is
asymmetric and will be kept explicit.  Section~\ref{sec:images} proves
projective faithfulness without using any target Hom-space dimension.  That
faithfulness is then used there to establish logarithmicity of the non-wall
objects $Q_j$, which is an input to the recognition argument below.  Once
the logarithmic extensions and their distinguished nonsplit quotients have
been identified, the target Hom upper bounds follow from these established
Virasoro short exact sequences.  Projective faithfulness is used once more
only after those upper bounds are known, to provide the matching lower
bounds.  Thus no target Hom-space dimension, and no target-side fullness
statement, enters the proof of projective faithfulness.

\subsection{Recognition of logarithmic extensions}

Throughout this section, Hom and unqualified Yoneda Ext spaces are taken
in the abelian category $\Oc_{c_{p,q}}$.  When all endpoints lie in
$\mathscr C$, Lemma~\ref{lem:Nakano-exact-compatibility} identifies the
same extension space with $\Ext^1_{\mathscr C}$.  All pushouts and
pullbacks below are therefore ordinary Virasoro-module constructions in
$\Oc_{c_{p,q}}$.

We shall repeatedly use the six-term Hom--Yoneda exact sequences of
Lemma~\ref{lem:prelim-Hom-Yoneda}; their connecting maps are the usual
pushout and pullback maps in the ambient abelian category.

\begin{lem}[Unique simple submodule of a nonsplit length-two extension]
\label{lem:length-two-socle}
Let
\[
 0\longrightarrow T'\longrightarrow A\longrightarrow T\longrightarrow0
\]
be a nonsplit extension of nonisomorphic simple objects in a finite-length
abelian category.  Then $T'$ is the unique simple subobject of $A$; in
particular $\Soc(A)=T'$.
\end{lem}

\begin{proof}
If a simple subobject $U\subset A$ is not contained in $T'$, its composite
with $A\twoheadrightarrow T$ is nonzero and hence an isomorphism.  This
would give a section of $A\twoheadrightarrow T$, contradicting
nonsplitting.  Thus every simple subobject lies in $T'$, and simplicity of
$T'$ gives the claim.  In particular, every isomorphism out of $A$ carries
this distinguished subobject onto the socle of the target.
\end{proof}

The following lemma will also be used in Section~\ref{sec:comparison}.

\begin{lem}[Recognition lemma]\label{lem:Nakano-recognition}
Let $n\ge1$, $1\le r\le p-1$, and put
\[
 a=np-r,\qquad b=np+r.
\]
Suppose that $E\in\Oc_{c_{p,q}}$ is logarithmic and fits into a short
exact sequence
\begin{equation}\label{eq:recognition-extension}
 0\longrightarrow K_a\longrightarrow E\longrightarrow K_b
 \longrightarrow0.
\end{equation}
Then, for the distinguished submodule $S_b\hookrightarrow K_a$ in
$0\to S_b\to K_a\to S_a\to0$, viewed inside $E$, the induced sequence
\[
 0\longrightarrow S_a\longrightarrow E/S_b\longrightarrow K_b
 \longrightarrow0
\]
is nonsplit, and
\[
 E/S_b\cong K(\tau_{n,r}),\qquad
 E\cong P(\tau_{n,r}).
\]
\end{lem}

\begin{proof}
Proposition~\ref{prop:Felder-labels} gives
\[
 \tau_{n,r}\in\mathcal T_{p_+,p_-}\setminus\mathcal T^0_{p_+,p_-},
 \qquad h_{\alpha_1}=h_a<h_b=h_{\alpha_2}<h_{(n+2)p-r}=h_{\alpha_3},
\]
together with
\[
 K_a\cong A(\tau_{n,r})=L(h_{\alpha_1},h_{\alpha_2}),\qquad
 K_b\cong B(\tau_{n,r})=L(h_{\alpha_2},h_{\alpha_3}).
\]
Fix these endpoint identifications.  The nonsplit Kac sequence
\[
 0\to S_b\to K_a\to S_a\to0
\]
and Lemma~\ref{lem:length-two-socle} show that $S_b=\Soc(K_a)$.
The same lemma applied to \eqref{eq:Nakano-tower-A} shows that the chosen
isomorphism $K_a\cong A(\tau_{n,r})$ carries $S_b$ onto the distinguished
$L(h_{\alpha_2})$.

By Lemma~\ref{lem:Nakano-membership} and the Serre property,
$K_a,K_b,E\in\mathscr C$.  Under the endpoint identifications,
\eqref{eq:recognition-extension} is literally the Virasoro-module
configuration
\[
 0\longrightarrow L(h_{\alpha_1},h_{\alpha_2})
 \longrightarrow E\longrightarrow
 L(h_{\alpha_2},h_{\alpha_3})\longrightarrow0
\]
used in Corollary~\ref{cor:Nakano-one-row}; there is no reversal of the
Yoneda variables.  Since $K_a$ and $K_b$ are ordinary, a split sequence
would make $E\cong K_a\oplus K_b$ ordinary.  Thus logarithmicity implies
that the extension is nonsplit.  Proposition~\ref{prop:one-row-staggered-quotient}
gives directly that
\begin{equation}\label{eq:recognition-first-quotient}
 0\longrightarrow S_a\longrightarrow E/S_b
 \longrightarrow K_b\longrightarrow0
\end{equation}
is nonsplit.  Its class is therefore
a nonzero element of
\[
 \Ext^1_{\mathscr C}(K_b,S_a)
 \cong
 \Ext^1_{\mathscr C}\bigl(B(\tau_{n,r}),L(h_{\alpha_1})\bigr)
 \cong\C
\]
by \eqref{eq:Nakano-first-Ext}.  Lemma~\ref{lem:one-dimensional-ext}
gives an isomorphism
\[
 \phi:E/S_b\xrightarrow{\sim}K(\tau_{n,r}).
\]
Fix such a $\phi$ and transport the quotient map to obtain
\begin{equation}\label{eq:recognition-second-extension}
 0\longrightarrow S_b\longrightarrow E
 \xrightarrow{\ \phi\circ\pi\ }K(\tau_{n,r})\longrightarrow0.
\end{equation}

This second sequence is also nonsplit.  Indeed, a retraction
$E\to S_b$ would restrict along $K_a\hookrightarrow E$ to a retraction
of $S_b\hookrightarrow K_a$, contradicting the nonsplit one-row Kac
sequence.  Hence \eqref{eq:recognition-second-extension} represents a
nonzero class in
\[
 \Ext^1_{\mathscr C}(K(\tau_{n,r}),S_b)
 =
 \Ext^1_{\mathscr C}
 \bigl(K(\tau_{n,r}),L(h_{\alpha_2})\bigr)
 \cong\C
\]
by Corollary~\ref{cor:Nakano-one-row}.  By the definition of $P(\tau)$
after Theorem~\ref{thm:Nakano-input} and
Lemma~\ref{lem:one-dimensional-ext}, every nonzero class
has an isomorphic middle term.  Therefore
$E\cong P(\tau_{n,r})$.
\end{proof}

\begin{cor}\label{cor:Q-Ptau}
For $n\geq1$ and $1\leq r\leq p-1$,
\[
 Q_{np+r}\cong P(\tau_{n,r}),
\]
and, with $b=np+r$, the quotient by the distinguished $S_b$ satisfies
\[
 Q_b/S_b\cong K(\tau_{n,r}),
\]
while the induced sequence
\[
 0\longrightarrow S_{np-r}\longrightarrow Q_b/S_b
 \longrightarrow K_b\longrightarrow0
\]
is nonsplit.
\end{cor}

\begin{proof}
Apply Lemma~\ref{lem:Nakano-recognition} to
Proposition~\ref{prop:Q-logarithmic} and
\eqref{eq:F-projective-filtration}.
\end{proof}

\subsection{Hom spaces along a reflection chain}

Fix a non-wall reflection chain and write
\[
 \mathsf S_i=S_{s_i},\qquad
 \mathsf K_i=K_{s_i},\qquad
 \mathsf Q_i=Q_{s_i}.
\]

\begin{lem}\label{lem:chain-simple-distinct}
The simple modules $\mathsf S_i$ on a fixed reflection chain are pairwise
nonisomorphic.
\end{lem}

\begin{proof}
If two distinct labels $s_i,s_j$ on the chain gave isomorphic simples,
then
Corollary~\ref{cor:positive-label-collisions} would give $q=2$ and their
sum equal to $p$.  But every $s_i$ with $i\ge1$ is greater than $p$, so
this cannot occur; the labels themselves are distinct by
Lemma~\ref{lem:reflection-partition}.
\end{proof}

\begin{lem}\label{lem:disjoint-constituents}
If two finite-length modules have no common simple Jordan--H\"older
constituent, then there is no nonzero morphism between them in either
direction.
\end{lem}

\begin{proof}
The image of a nonzero morphism is a nonzero finite-length module and
therefore has a simple constituent common to the source and the target.
\end{proof}

\begin{lem}\label{lem:chain-filtrations}
For every $i\geq0$,
\begin{equation}\label{eq:A-chain}
 0\longrightarrow\mathsf S_{i+1}\longrightarrow\mathsf K_i
 \longrightarrow\mathsf S_i\longrightarrow0
\end{equation}
is non-split.  For $i\geq1$,
\begin{equation}\label{eq:Q-chain}
 0\longrightarrow\mathsf K_{i-1}\longrightarrow\mathsf Q_i
 \longrightarrow\mathsf K_i\longrightarrow0,
\end{equation}
and $\mathsf Q_0=\mathsf K_0$.
\end{lem}

\begin{proof}
If $s_i=n_ip+r_i$ with $1\leq r_i\leq p-1$, then
$s_{i-1}=n_ip-r_i$ for $i\ge1$ and
$s_{i+1}=(n_i+2)p-r_i$.  Thus \eqref{eq:A-chain} is the Kac sequence and
\eqref{eq:Q-chain} is Proposition~\ref{prop:projective-images}.
\end{proof}

\begin{lem}[Nakano triple at a positive chain vertex]
\label{lem:chain-Nakano-triple}
For $i\ge1$, one has $s_i>p$.  Hence in the Euclidean decomposition
$s_i=n_ip+r_i$ with $1\le r_i\le p-1$ one has $n_i\ge1$.  Put
$\tau_i:=\tau_{n_i,r_i}$.  Then
\[
 (h_{\alpha_1},h_{\alpha_2},h_{\alpha_3})
 =(h_{s_{i-1}},h_{s_i},h_{s_{i+1}}),
\]
\[
 A(\tau_i)\cong K_{s_{i-1}}=\mathsf K_{i-1},\qquad
 B(\tau_i)\cong K_{s_i}=\mathsf K_i.
\]
If
\[
 \sigma_i:\mathsf S_i\hookrightarrow\mathsf K_{i-1}
 \hookrightarrow\mathsf Q_i
\]
is the distinguished simple submodule, then
\begin{equation}\label{eq:distinguished-quotient}
 \mathsf Q_i/\sigma_i(\mathsf S_i)\cong K(\tau_i),
\end{equation}
and the induced sequence
\begin{equation}\label{eq:distinguished-quotient-extension}
 0\longrightarrow\mathsf S_{i-1}\longrightarrow
 \mathsf Q_i/\sigma_i(\mathsf S_i)\longrightarrow
 \mathsf K_i\longrightarrow0
\end{equation}
is non-split.
\end{lem}

\begin{proof}
The first two assertions are Proposition~\ref{prop:Felder-labels} together
with the reflection identities in Lemma~\ref{lem:chain-filtrations}.
Both the quotient assertion and the nonsplitting of
\eqref{eq:distinguished-quotient-extension} are exactly the corresponding
statements of Corollary~\ref{cor:Q-Ptau}.
\end{proof}

We keep the notation
\begin{equation}\label{eq:distinguished-sigma}
 \sigma_i:\mathsf S_i\hookrightarrow\mathsf K_{i-1}
 \hookrightarrow\mathsf Q_i
\end{equation}
for this distinguished submodule.

In the Grothendieck group these filtrations give
\begin{equation}\label{eq:Q-JH-vector}
 [\mathsf Q_0]=[\mathsf S_0]+[\mathsf S_1],\qquad
 [\mathsf Q_i]=[\mathsf S_{i-1}]+2[\mathsf S_i]+[\mathsf S_{i+1}]
 \quad(i\ge1).
\end{equation}
The possible collision of simple labels when $q=2$ is treated separately
below.

\begin{lem}[Hom spaces for a nonsplit length-two module]
\label{lem:length-two-Hom}
Let
\[
 0\longrightarrow T'\longrightarrow A\longrightarrow T\longrightarrow0
\]
be nonsplit, where $T$ and $T'$ are nonisomorphic simple lowest-weight
Virasoro modules.  Then
\begin{equation}\label{eq:length-two-Hom}
 \End(A)=\C,\qquad \Hom(A,T')=0,\qquad \Hom(A,T)=\C.
\end{equation}
\end{lem}

\begin{proof}
Since $T$ and $T'$ are nonisomorphic simple objects,
\[
 \End(T)=\End(T')=\C,\qquad
 \Hom(T,T')=\Hom(T',T)=0.
\]
Applying $\Hom(-,T')$ gives
\[
 0\longrightarrow\Hom(A,T')\longrightarrow\End(T')
 \xrightarrow{\partial}\Ext^1(T,T').
\]
The connecting morphism sends $\id_{T'}$ to the class of the given
extension.  Since the extension is nonsplit, this class is nonzero.  As
$\End(T')=\C$, the map $\partial$ is injective, so $\Hom(A,T')=0$.

Applying $\Hom(-,T)$ gives
\[
 0\longrightarrow\End(T)\longrightarrow\Hom(A,T)
 \longrightarrow\Hom(T',T)=0,
\]
hence $\Hom(A,T)\cong\C$.  Finally, applying $\Hom(A,-)$ to the original
sequence yields
\[
 0\longrightarrow\Hom(A,T')\longrightarrow\End(A)
 \longrightarrow\Hom(A,T).
\]
The first term is zero and the last is one-dimensional.  Since
$\id_A\ne0$, it follows that $\End(A)\cong\C$.
\end{proof}

\begin{lem}[Kac--Kac Hom spaces]\label{lem:Kac-Kac-Hom}
For $i,j\ge0$,
\[
 \Hom(\mathsf K_i,\mathsf K_j)=
 \begin{cases}
  \C,&j=i,\\
  \C,&j=i-1\text{ and }i\ge1,\\
  0,&\text{otherwise}.
 \end{cases}
\]
In particular, $\Hom(\mathsf K_0,\mathsf K_j)=0$ for $j\ge1$.
\end{lem}

\begin{proof}
The simples on a fixed chain are pairwise nonisomorphic; this is
Lemma~\ref{lem:chain-simple-distinct}.  The length-two calculation of
Lemma~\ref{lem:length-two-Hom} gives
\[
 \Hom(\mathsf K_i,\mathsf S_i)=\C,
 \qquad
 \Hom(\mathsf K_i,\mathsf S_{i+1})=0.
\]
If $k\notin\{i,i+1\}$, then $\mathsf S_k$ is not a Jordan--H\"older
constituent of $\mathsf K_i$.  Lemma~\ref{lem:disjoint-constituents}
then gives $\Hom(\mathsf K_i,\mathsf S_k)=0$.  Hence
\[
 \Hom(\mathsf K_i,\mathsf S_k)=
 \begin{cases}\C,&k=i,\\0,&k\ne i.\end{cases}
\]
Apply $\Hom(\mathsf K_i,-)$ to \eqref{eq:A-chain} for $\mathsf K_j$.
If $j=i-1$, the nonzero map is the explicit factorization
\[
 \mathsf K_i\twoheadrightarrow\mathsf S_i
 \hookrightarrow\mathsf K_{i-1}.
\]
Equivalently, the relevant left-exact segment is
\[
 0\longrightarrow\Hom(\mathsf K_i,\mathsf S_i)
 \longrightarrow\Hom(\mathsf K_i,\mathsf K_{i-1})
 \longrightarrow\Hom(\mathsf K_i,\mathsf S_{i-1})=0,
\]
so the space is one-dimensional.  If $j=i$, \eqref{eq:length-two-Hom} gives
$\End(\mathsf K_i)=\C$.  For all other $j$, both neighboring simple Hom
terms vanish, and exactness gives zero.
\end{proof}

\begin{lem}[Pushout by a quotient]\label{lem:pushout-quotient}
Let
\[
 0\longrightarrow A\longrightarrow E\longrightarrow B\longrightarrow0
\]
be an extension in an abelian category, and let $q:A\twoheadrightarrow C$
have kernel $K$.  Identify $K$ with its image under the composite
$K\hookrightarrow A\hookrightarrow E$.  The pushout of the extension
along $q$ is then canonically represented by
\[
 0\longrightarrow C\longrightarrow E/K\longrightarrow B\longrightarrow0.
\]
\end{lem}

\begin{proof}
The composite $K\hookrightarrow A\hookrightarrow E$ is monic.  Passing to
$E/K$ identifies the image of $A$ with $A/K\cong C$ and leaves quotient
$B$.  The resulting square with $A\to C$ is a pushout by the universal
property of the cokernel of $K\to A$.
\end{proof}

\begin{lem}\label{lem:Q-simple-Hom}
For $i\geq1$,
\[
 \Hom(\mathsf Q_i,\mathsf S_j)=
 \begin{cases}
  \C,&j=i,\\
  0,&j\neq i.
 \end{cases}
\]
\end{lem}

\begin{proof}
By Lemmas~\ref{lem:chain-simple-distinct},
\ref{lem:disjoint-constituents}, and \ref{lem:length-two-Hom},
\[
 \Hom(\mathsf K_i,\mathsf S_j)=\delta_{ij}\C,\qquad
 \Hom(\mathsf K_{i-1},\mathsf S_j)=\delta_{i-1,j}\C.
\]
Applying $\Hom(-,\mathsf S_j)$ to \eqref{eq:Q-chain} therefore gives the
explicit exact sequence
\begin{equation}\label{eq:Q-simple-long-exact}
 0\to \delta_{ij}\C
 \to\Hom(\mathsf Q_i,\mathsf S_j)
 \to\delta_{i-1,j}\C
 \xrightarrow{\partial_j}\Ext^1_{\Oc}(\mathsf K_i,\mathsf S_j),
\end{equation}
where $\delta_{uv}$ is the Kronecker symbol.  Thus the Hom space is zero
unless $j=i-1$ or $i$.  For $j=i$, the canonical composite
$\mathsf Q_i\twoheadrightarrow\mathsf K_i\twoheadrightarrow\mathsf S_i$
is an epimorphism, hence nonzero, and
\eqref{eq:Q-simple-long-exact} makes the space one-dimensional.

For $j=i-1$, let
$\pi_{i-1}:\mathsf K_{i-1}\twoheadrightarrow\mathsf S_{i-1}$ be the
head projection, and let
\[
 \xi_i\in\Ext^1_{\Oc}(\mathsf K_i,\mathsf K_{i-1})
\]
be the Yoneda class of \eqref{eq:Q-chain}.  By
Lemma~\ref{lem:length-two-Hom},
\[
 \Hom(\mathsf K_{i-1},\mathsf S_{i-1})=\C\pi_{i-1}.
\]
The connecting morphism is covariant in the second Yoneda variable and is
therefore given by pushout in $\Oc_{c_{p,q}}$:
\[
 \partial_{i-1}(\pi_{i-1})=(\pi_{i-1})_*\xi_i
 \in\Ext^1_{\Oc}(\mathsf K_i,\mathsf S_{i-1}).
\]
Since $\ker\pi_{i-1}=\mathsf S_i$, embedded in $\mathsf Q_i$ as
$\sigma_i(\mathsf S_i)$, Lemma~\ref{lem:pushout-quotient} identifies this
pushout with
\[
 0\longrightarrow\mathsf S_{i-1}\longrightarrow
 \mathsf Q_i/\sigma_i(\mathsf S_i)\longrightarrow\mathsf K_i
 \longrightarrow0.
\]
By Lemma~\ref{lem:chain-Nakano-triple}, this sequence is nonsplit in
$\Oc_{c_{p,q}}$.  Its Yoneda class in
$\Ext^1_{\Oc}(\mathsf K_i,\mathsf S_{i-1})$ is therefore nonzero, so
$\partial_{i-1}$ is nonzero.  Its domain is
one-dimensional, hence $\partial_{i-1}$ is injective.  Exactness of \eqref{eq:Q-simple-long-exact} gives
$\Hom(\mathsf Q_i,\mathsf S_{i-1})=0$.
\end{proof}

\begin{lem}\label{lem:Q-Kac-Hom}
For $i\geq1$,
\[
 \Hom(\mathsf Q_i,\mathsf K_j)=
 \begin{cases}
  \C,&j=i-1\text{ or }j=i,\\
  0,&\text{otherwise}.
 \end{cases}
\]
\end{lem}

\begin{proof}
Apply $\Hom(\mathsf Q_i,-)$ to \eqref{eq:A-chain}.  Its beginning is
\[
 0\to\Hom(\mathsf Q_i,\mathsf S_{j+1})
 \to\Hom(\mathsf Q_i,\mathsf K_j)
 \to\Hom(\mathsf Q_i,\mathsf S_j)
 \to\Ext^1_{\Oc}(\mathsf Q_i,\mathsf S_{j+1}).
\]
Lemma~\ref{lem:Q-simple-Hom} gives the vanishing outside $j=i-1,i$.
If $j=i-1$, the relevant segment is
\[
 0\to\Hom(\mathsf Q_i,\mathsf S_i)
 \to\Hom(\mathsf Q_i,\mathsf K_{i-1})
 \to\Hom(\mathsf Q_i,\mathsf S_{i-1})=0,
\]
so $\Hom(\mathsf Q_i,\mathsf K_{i-1})\cong\C$.  If $j=i$, the relevant
segment begins
\[
 0=\Hom(\mathsf Q_i,\mathsf S_{i+1})
 \to\Hom(\mathsf Q_i,\mathsf K_i)
 \to\Hom(\mathsf Q_i,\mathsf S_i),
\]
so $\Hom(\mathsf Q_i,\mathsf K_i)$ injects into a one-dimensional space.  The canonical composite
\[
 \mathsf Q_i\twoheadrightarrow\mathsf K_i
 \twoheadrightarrow\mathsf S_i
\]
is an epimorphism, hence nonzero.  Therefore
$\mathsf Q_i\twoheadrightarrow\mathsf K_i$ is nonzero, and this Hom
space is also one-dimensional.
\end{proof}

\begin{prop}[Target upper bounds]\label{prop:target-Hom-upper}
Within a fixed reflection chain,
\[
 \dim\Hom(\mathsf Q_i,\mathsf Q_j)\le
 \begin{cases}
  2,&i=j\geq1,\\
  1,&i=j=0,\\
  1,&|i-j|=1,\\
  0,&|i-j|\geq2.
 \end{cases}
\]
\end{prop}

\begin{proof}
For $i,j\ge1$, apply $\Hom(\mathsf Q_i,-)$ to
$0\to\mathsf K_{j-1}\to\mathsf Q_j\to\mathsf K_j\to0$.  The
relevant left-exact segment is
\[
 0\longrightarrow\Hom(\mathsf Q_i,\mathsf K_{j-1})
 \longrightarrow\Hom(\mathsf Q_i,\mathsf Q_j)
 \longrightarrow\Hom(\mathsf Q_i,\mathsf K_j).
\]
Consequently
\[
 \dim\Hom(\mathsf Q_i,\mathsf Q_j)
 \le
 \dim\Hom(\mathsf Q_i,\mathsf K_{j-1})
 +\dim\Hom(\mathsf Q_i,\mathsf K_j).
\]
Lemma~\ref{lem:Q-Kac-Hom} says
\[
 \dim\Hom(\mathsf Q_i,\mathsf K_t)
 =\mathbf1_{t=i-1}+\mathbf1_{t=i},
\]
so the right-hand side is $2$ for $j=i$, $1$ for $|i-j|=1$, and $0$
for $|i-j|\ge2$.

At the left boundary, $\mathsf Q_0=\mathsf K_0$ and
Lemma~\ref{lem:Kac-Kac-Hom} gives $\End(\mathsf Q_0)=\C$.  For $j\ge2$,
applying $\Hom(\mathsf K_0,-)$ to \eqref{eq:Q-chain} gives
$\Hom(\mathsf Q_0,\mathsf Q_j)=0$, while the reverse space vanishes by
Lemma~\ref{lem:Q-Kac-Hom}.  Moreover,
\[
 \Hom(\mathsf Q_0,\mathsf Q_1)\cong\C,
 \qquad
 \Hom(\mathsf Q_1,\mathsf Q_0)\cong\C.
\]
Indeed, $\mathsf Q_0=\mathsf K_0$, and applying
$\Hom(\mathsf K_0,-)$ to
$0\to\mathsf K_0\to\mathsf Q_1\to\mathsf K_1\to0$, together with
Lemma~\ref{lem:Kac-Kac-Hom}, gives the first equality; the second is
Lemma~\ref{lem:Q-Kac-Hom}.
\end{proof}

\begin{rmk}[Logical order of the Hom calculation]\label{rmk:Hom-no-circularity}
The dependence is
\[
\begin{aligned}
 \text{projective faithfulness}
 &\Longrightarrow \text{logarithmicity}
 \Longrightarrow \text{recognition and distinguished nonsplitting},\\
 &\Longrightarrow \text{target Hom upper bounds}.
\end{aligned}
\]
Projective faithfulness itself was proved in Section~\ref{sec:images}
without using any target Hom-space dimension.  After the upper bounds are
known, the same faithfulness supplies lower bounds from the source Hom
spaces, yielding the exact target dimensions.  Full faithfulness is then a
consequence of those dimensions, not an input to them.
\end{rmk}

\begin{prop}\label{prop:target-Hom-table}
Within a fixed reflection chain,
\begin{equation}\label{eq:target-Hom-table}
 \dim\Hom(\mathsf Q_i,\mathsf Q_j)=
 \begin{cases}
  2,&i=j\geq1,\\
  1,&i=j=0,\\
  1,&|i-j|=1,\\
  0,&|i-j|\geq2.
 \end{cases}
\end{equation}
\end{prop}

\begin{proof}
Let $d_{ij}$ denote the right-hand side of
\eqref{eq:target-Hom-table}.  Proposition~\ref{prop:source-Hom-package},
projective faithfulness, and Proposition~\ref{prop:target-Hom-upper} give
\[
 d_{ij}
 =\dim\Hom_{\KL^k}(P_{s_i},P_{s_j})
 \le \dim\Hom_{\Oc}(\mathsf Q_i,\mathsf Q_j)
 \le d_{ij}.
\]
Hence equality holds throughout.
\end{proof}

\begin{cor}[Local endomorphism algebras]\label{cor:target-local-endomorphisms}
For $n\ge1$ and $1\le r\le p-1$,
\[
 \End(Q_{np+r})
 =\C\id\oplus\C F(N_{np+r})
 \cong\C[\varepsilon]/(\varepsilon^2).
\]
In particular, $Q_{np+r}$ is indecomposable.
\end{cor}

\begin{proof}
The label $np+r$ occurs at a positive index on its reflection chain by
Lemma~\ref{lem:reflection-partition}.  Equation~\eqref{eq:target-Hom-table}
gives $\dim\End(Q_{np+r})=2$.  Projective faithfulness gives
$F(N_{np+r})\neq0$, while
$F(N_{np+r})^2=F(N_{np+r}^2)=0$.  A nonzero square-zero
endomorphism cannot be a scalar multiple of the identity, so $\id$ and
$F(N_{np+r})$ are linearly independent.  Since the endomorphism space has
dimension two, they form a basis.  The map
\[
 \C[\varepsilon]/(\varepsilon^2)\longrightarrow\End(Q_{np+r}),
 \qquad \varepsilon\longmapsto F(N_{np+r}),
\]
is therefore an algebra isomorphism.  In particular the endomorphism
algebra is local, so $Q_{np+r}$ is indecomposable.
\end{proof}

The following consequence is not used in the proof of full faithfulness;
we record it here for the comparison argument in Section~\ref{sec:comparison}.

The full-faithfulness argument below is already complete at the level of
Hom-space dimensions.  We record next a canonical radical element only
for the comparison with Drinfeld--Sokolov reduction in
Section~\ref{sec:comparison}; it is not used to prove fullness.

\begin{cor}[Canonical conformal radical]\label{cor:target-conformal-radical}
Let $P=P_{np+r}$ and $Q=F(P)$, with $n\ge1$ and $1\le r\le p-1$.
Let $D_P$ be the canonical affine conformal nilpotent of
Lemma~\ref{lem:source-L0-radical}.  For every actual generalized
$L_0$-eigenvalue $\lambda$ of $Q$, let $Q^{\mathrm{gen}}_\lambda$ be the
corresponding generalized eigenspace and define
\[
 D_Q|_{Q^{\mathrm{gen}}_\lambda}
 :=(L_0-\lambda\id)|_{Q^{\mathrm{gen}}_\lambda}.
\]
Then these operators assemble to a Virasoro-module endomorphism and
\begin{equation}\label{eq:canonical-radical-matching}
 F(D_P)=D_Q,
 \qquad
 \rad\End(Q)=\C D_Q=\C F(N_P).
\end{equation}
\end{cor}

\begin{proof}
By Proposition~\ref{prop:projective-images}, $Q$ is supported in the
single residue sector $\xi_{n,r}$.  The Virasoro commutator
$[L_0,L_m]=-mL_m$ gives
\[
 D_{\lambda-m}L_m=L_mD_\lambda,
\]
so the nilpotent parts on the actual generalized $L_0$-eigenspaces
assemble to a Virasoro-module endomorphism $D_Q$.  It is nonzero because $Q$ is logarithmic.  It is globally locally
nilpotent in the algebraic sense: every vector is a finite sum of
generalized $L_0$-eigenvectors, and on each summand the corresponding
nilpotent part is nilpotent, so one common power annihilates that vector.
Hence $D_Q$ cannot be invertible: for any nonzero $v$ there is an $N$
with $D_Q^Nv=0$, whereas invertibility of $D_Q$ would make $D_Q^N$
invertible.
Corollary~\ref{cor:target-local-endomorphisms} identifies $\End(Q)$ with a
local dual-number algebra, so every nonunit belongs to its Jacobson
radical.  Hence
\[
 0\ne D_Q\in\rad\End(Q).
\]
Thus $D_Q$ is intrinsic to the Virasoro module $Q$ and does not depend on
a chosen generator of $\rad\End(Q)$.  The radical is one-dimensional,
so $D_Q$ spans it and $D_Q^2=0$.

On the source,
\[
 \theta^-_P=\lambda_P(\id+2\pi iD_P).
\]
Applying $F$ and using preservation of the minus twist gives
\[
 e^{2\pi iL_0}|_Q
 =\lambda_P\bigl(\id+2\pi iF(D_P)\bigr).
\]
Because $D_P$ is a nonzero scalar multiple of $N_P$ and $F$ is faithful on
projectives, $F(D_P)\ne0$; moreover $F(D_P)^2=0$.  Thus $F(D_P)$ is a
nonunit of the same local algebra and lies in its radical.

Every actual generalized eigenvalue $\lambda$ of $Q$ lies in the single
residue class $\xi_{n,r}$, so the scalar $e^{2\pi i\lambda}$ is independent
of $\lambda$; denote it by $\mu_Q$.  On each finite-dimensional generalized
eigenspace one has $L_0=\lambda\id+D_Q$ and $D_Q^2=0$, hence
\[
 e^{2\pi iL_0}=\mu_Q(\id+2\pi iD_Q).
\]
These identities are compatible with the algebraic direct-sum
decomposition into generalized eigenspaces and therefore assemble to the
global identity on $Q$.  Moreover, for every Virasoro mode $L_m$,
\[
 e^{2\pi iL_0}L_me^{-2\pi iL_0}=e^{-2\pi im}L_m=L_m,
\]
so $e^{2\pi iL_0}|_Q$ is a Virasoro-module automorphism and hence an
element of $\End(Q)$.  Since
\[
 \End(Q)/\rad\End(Q)\cong\C,
\]
the images of the two exponential identities in this quotient give
$\mu_Q=\lambda_P$.  Subtracting their scalar
parts then yields
\[
 2\pi i\lambda_PF(D_P)=2\pi i\lambda_PD_Q,
\]
and hence $F(D_P)=D_Q$.  Finally,
Lemma~\ref{lem:source-L0-radical} gives $\C D_P=\C N_P$, so projective
faithfulness yields the remaining equality in
\eqref{eq:canonical-radical-matching}.
\end{proof}

\subsection{Separation of reflection chains and walls}

For a finite-length module $M$, let $\operatorname{JH}(M)$ denote the set
of isomorphism classes of its simple Jordan--H\"older constituents, with
multiplicities forgotten.

\begin{lem}[Composition-factor collisions]\label{lem:composition-collisions}
For a chain beginning at $u$, write $Q_i^u:=Q_{s_i(u)}$.  If $q>2$,
distinct reflection chains have disjoint simple
composition-factor sets.  If $q=2$, then $p$ is odd and the only chains
which can share a simple are those beginning at $r$ and $p-r$.  Their only
common simple is $S_r\cong S_{p-r}$, and
\[
 \operatorname{JH}(Q_i^r)\cap
 \operatorname{JH}(Q_j^{p-r})=\varnothing
 \qquad\text{unless }i,j\in\{0,1\}.
\]
Thus the common isomorphism class occurs only in objects of chain index
$0$ or $1$.
\end{lem}

\begin{proof}
Corollary~\ref{cor:positive-label-collisions} shows that a collision of
distinct positive labels forces $q=2$ and complementary labels summing to
$p$.  Since every chain label $s_i$ with $i\ge1$ is larger than $p$, only
the initial labels $r,p-r$ can collide.  Equations
\eqref{eq:A-chain}--\eqref{eq:Q-chain} show where the corresponding simple
occurs.
\end{proof}

\begin{lem}\label{lem:cross-chain}
The Hom spaces between images of projectives belonging to distinct reflection
chains are zero.
\end{lem}

\begin{proof}
For $q>2$, use Lemmas~\ref{lem:composition-collisions} and
\ref{lem:disjoint-constituents}.  Suppose $q=2$ and compare the chains beginning at $r$ and $p-r$, with
$1\le r<p-r$.  Fix once and for all an isomorphism
\[
 \iota:S_{p-r}\xrightarrow{\sim}S_r,
\]
and write $S:=S_r$, using $\iota$ to identify the two simple heads below.
Put
\[
 A=K_r,\qquad B=K_{p-r}.
\]
Their Kac sequences, after using $\iota$ on the head of $B$, are
\[
 0\to S_{2p-r}\to A\to S\to0,
 \qquad
 0\to S_{p+r}\to B\to S\to0,
\]
with nonisomorphic submodules.  Thus
\[
 \operatorname{JH}(Q_0^r)=\{S,S_{2p-r}\},\qquad
 \operatorname{JH}(Q_0^{p-r})=\{S,S_{p+r}\},
\]
while \eqref{eq:Q-chain} gives
\[
 \operatorname{JH}(Q_1^r)=\{S,S_{2p-r},S_{2p+r}\},\qquad
 \operatorname{JH}(Q_1^{p-r})=\{S,S_{p+r},S_{3p-r}\},
\]
with multiplicities suppressed.  Let $\pi_A:A\twoheadrightarrow S$ be the
natural head map and let $\pi_B:B\twoheadrightarrow S$ be the natural
head map to $S_{p-r}$ followed by $\iota$.  Lemma~\ref{lem:length-two-Hom} gives
\[
 \Hom(A,S)=\C\pi_A,\qquad \Hom(B,S)=\C\pi_B.
\]
Thus, if $f:A\to B$, then $\pi_Bf=c\pi_A$ for some $c\in\C$.
If $c=0$, the image lies in the simple submodule $S_{p+r}\subset B$; since
$S_{p+r}$ is not a Jordan--H\"older constituent of $A$,
Lemma~\ref{lem:disjoint-constituents} gives $f=0$.  If $c\neq0$, the
restriction of $f$ to $S_{2p-r}\subset A$ is zero because
$\Hom(S_{2p-r},B)=0$ by disjoint constituents.  Hence $f$ factors through
$\pi_A$ as $f=\bar f\pi_A$.  The equality
$\pi_B\bar f=c\id_S$ makes $c^{-1}\bar f$ a section of $\pi_B$, contrary
to nonsplitting.  Thus $\Hom(A,B)=0$, and the same argument with $A$ and
$B$ interchanged gives $\Hom(B,A)=0$.

By Lemma~\ref{lem:composition-collisions}, all cross-chain pairs are
already disjoint except the boundary indices $i,j\in\{0,1\}$.  Thus it
remains, in both morphism directions, to check exactly the four index
pairs $(0,0),(1,0),(0,1),(1,1)$.  The index-zero pair has just been
settled:
\[
 \Hom(Q_0^r,Q_0^{p-r})=
 \Hom(Q_0^{p-r},Q_0^r)=0.
\]
For $f:Q^r_1\to Q^{p-r}_0=K_{p-r}$, the composite with
$K_{p-r}\twoheadrightarrow S$ is zero by
Lemma~\ref{lem:Q-simple-Hom}; hence the image lies in $S_{p+r}$, which is
not a constituent of $Q^r_1$, so $f=0$.  Conversely, if
$g:K_{p-r}\to Q^r_1$, its composite with
$Q^r_1\twoheadrightarrow K_{2p-r}$ vanishes by disjoint constituents;
therefore $g$ factors through $K_r\subset Q^r_1$, and the index-zero
calculation gives $g=0$.  Thus
\[
 \Hom(Q_1^r,Q_0^{p-r})=
 \Hom(Q_0^{p-r},Q_1^r)=0.
\]
Interchanging $r$ and $p-r$ gives, by the same two explicit
factorizations,
\[
 \Hom(Q_0^r,Q_1^{p-r})=
 \Hom(Q_1^{p-r},Q_0^r)=0.
\]
Finally apply $\Hom(Q_1^r,-)$ to
\[
 0\to Q_0^{p-r}\to Q_1^{p-r}\to K_{p+r}\to0.
\]
The Hom spaces to both end terms vanish: the first by the preceding
boundary calculation and the second by disjoint Jordan--H\"older
constituents.  Hence $\Hom(Q_1^r,Q_1^{p-r})=0$.  Reversing $r$ and
$p-r$ gives $\Hom(Q_1^{p-r},Q_1^r)=0$.  Thus the four boundary cases can
be summarized as
\[
\begin{array}{c|cc}
 &Q_0^{p-r}&Q_1^{p-r}\\ \hline
 Q_0^r&0&0\\
 Q_1^r&0&0
\end{array}
\]
for Hom spaces from the $r$-chain to the $(p-r)$-chain, and the reverse
matrix is likewise zero.
\end{proof}

\begin{lem}[Wall objects]\label{lem:wall-Hom}
For $m,n\geq1$,
\[
 \Hom(Q_{np},Q_{mp})=
 \begin{cases}\C,&m=n,\\0,&m\neq n,\end{cases}
\]
and every Hom space between a wall object and the image of a non-wall
projective is zero in either direction.
\end{lem}

\begin{proof}
Proposition~\ref{prop:projective-images} gives
$Q_{np}=K_{np}=S_{np}$.  Corollary~\ref{cor:positive-label-collisions}
shows that a wall label cannot collide with a positive non-wall label, nor
can two distinct wall labels collide.  The assertion follows from
simplicity and Lemma~\ref{lem:disjoint-constituents}.
\end{proof}

\subsection{Full faithfulness and morphisms between projectives}

\begin{thm}\label{thm:F-fully-faithful}
The restriction
\[
 F|_{\mathcal P^k}:\mathcal P^k\longrightarrow\Oc_{c_{p,q}}
\]
is fully faithful.
\end{thm}

\begin{proof}
For indecomposable projectives $P_i,P_j$, projective faithfulness gives
an injection
\[
 \Hom_{\KL^k}(P_i,P_j)
 \hookrightarrow
 \Hom_{\Oc_{c_{p,q}}}(F(P_i),F(P_j)).
\]
The source dimensions are those of
Proposition~\ref{prop:source-Hom-package}.  Within one reflection chain,
Proposition~\ref{prop:target-Hom-table} gives the matching target upper
bounds; between distinct chains, including the $q=2$ boundary collision,
Lemma~\ref{lem:cross-chain} gives the required target vanishing; and wall
objects are covered by Lemma~\ref{lem:wall-Hom}.  Hence the displayed
injection is an isomorphism for every pair of indecomposable projectives.
Every projective object is a finite direct sum of indecomposable
projectives.  Since Hom between finite biproducts is the corresponding
matrix space of component Hom spaces, full faithfulness on all pairs of
indecomposable projectives implies full faithfulness on $\mathcal P^k$.
\end{proof}

\begin{cor}[Projective endomorphism algebras along finite reflection segments]\label{cor:finite-zigzag-algebra}
For a reflection chain and every $N\ge0$, $F$ induces an algebra
isomorphism
\[
 \End_{\KL^k}\!\left(\bigoplus_{i=0}^N P_{s_i}\right)
 \xrightarrow{\ \sim\ }
 \End_{\Oc}\!\left(\bigoplus_{i=0}^N Q_{s_i}\right).
\]
In particular, after choosing nonzero adjacent generators, the target
category has
\[
 \dim\Hom(Q_{s_i},Q_{s_j})=
 \begin{cases}
  1,&i=j=0,\\
  2,&i=j\ge1,\\
  1,&|i-j|=1,\\
  0,&|i-j|\ge2,
 \end{cases}
\]
and, for $i\ge1$, if
\[
 \nu_i=x_{i-1}\circ y_{i-1}
 \in\rad\End_{\KL^k}(P_{s_i})
\]
is the source radical loop fixed in
Proposition~\ref{prop:source-Hom-package}, then
\[
 \rho_i:=F(\nu_i)
\]
is nonzero and
\[
 \End(Q_{s_i})=\C\id\oplus\C\rho_i,
 \qquad \rho_i^2=0,
 \qquad \rad\End(Q_{s_i})=\C\rho_i.
\]
For $i\ge1$, every nonzero length-two backtracking endomorphism of
$Q_{s_i}$ is a nonzero scalar multiple of $\rho_i$; at the initial vertex
$Q_{s_0}$ every length-two backtracking endomorphism is zero.  Every
length-two path whose endpoints have distance two is zero.
\end{cor}

\begin{proof}
Full faithfulness identifies every matrix entry and respects composition.
The Hom dimensions are Proposition~\ref{prop:target-Hom-table}, and the
endomorphism assertion is Corollary~\ref{cor:target-local-endomorphisms}.
For $i\ge1$, a nonzero backtracking composite factors through an adjacent
indecomposable $Q_{s_{i\pm1}}$, which is not isomorphic to $Q_{s_i}$.  If
the composite were invertible, one of the adjacent maps would split and
$Q_{s_i}$ would be a direct summand of that adjacent indecomposable, a
contradiction.  Hence the composite is a nonunit and therefore lies in
\[
 \rad\End(Q_{s_i})=\C\rho_i.
\]
At $Q_{s_0}$ the endomorphism algebra is $\C$; a nonzero backtracking
composite would therefore be invertible and is excluded by the same
splitting argument.  Finally, a length-two path joining vertices at
distance two vanishes because the corresponding Hom space is zero.
\end{proof}

\begin{proof}[Proof of Theorem~\ref{thm:projective-classification}]
For $1\le j\le p-1$, $P_j=V_j$ and
Lemma~\ref{lem:F-small-Weyl} gives $F(P_j)=K_j$.  Here
\[
 0\longrightarrow S_{2p-j}\longrightarrow K_j
 \longrightarrow S_j\longrightarrow0
\]
is nonsplit, and its two simple factors are nonisomorphic by
Corollary~\ref{cor:positive-label-collisions}; and Lemma~\ref{lem:length-two-Hom} gives
$\End(K_j)=\C$.  If
$p\mid j$, Proposition~\ref{prop:projective-images} gives
$F(P_j)=K_j=S_j$, again with scalar endomorphisms.  For
$j=np+r$, $n\ge1$, logarithmicity is Proposition~\ref{prop:Q-logarithmic},
the identification with $P(\tau_{n,r})$ is Corollary~\ref{cor:Q-Ptau}, and
the local endomorphism algebra and indecomposability are
Corollary~\ref{cor:target-local-endomorphisms}.  Full faithfulness is
Theorem~\ref{thm:F-fully-faithful}.
\end{proof}

\begin{cor}[Reflection-chain realization]\label{cor:blockwise-projective-realization}
For each non-wall reflection chain $\{s_i\}_{i\ge0}$,
\[
 F:\operatorname{add}\{P_{s_i}:i\ge0\}
 \xrightarrow{\ \sim\ }
 \operatorname{add}\{Q_{s_i}:i\ge0\}
\]
is an equivalence onto a full additive Virasoro subcategory.  The initial
object is the ordinary Kac module $Q_{s_0}=K_{s_0}$, while for every
$i\ge1$ the object $Q_{s_i}$ is the logarithmic module $P(\tau_i)$ from
Nakano's extension theory.  Wall projectives are sent to the isolated
simple modules $K_{np}$.
\end{cor}

\begin{proof}
Full faithfulness, Proposition~\ref{prop:projective-images}, and
Corollary~\ref{cor:Q-Ptau} give the asserted equivalence on finite direct
sums of the listed indecomposables.  If $\operatorname{add}$ is understood
to include direct summands, let $e\in\End_{\Oc}(F(P))$ be an idempotent,
where $P$ is such a finite direct sum in the fixed chain.  Full
faithfulness gives a unique idempotent
$\widetilde e\in\End_{\KL^k}(P)$ with $F(\widetilde e)=e$.  Since
$\KL^k(\mathfrak{sl}_2)$ is abelian, $\widetilde e$ splits; its image is a
direct summand of the projective object $P$ and is therefore projective.
By Krull--Schmidt, every indecomposable summand of this image is
isomorphic to one of the $P_{s_i}$ already occurring in the chosen
finite direct sum $P$.  Hence the lifted summand remains in
$\operatorname{add}\{P_{s_i}:i\ge0\}$.  Applying $F$ identifies it with
the target retract defined by $e$.  Thus the equivalence is essentially
surjective onto the full additive, idempotent-complete closure appearing
in the statement.
\end{proof}

\begin{cor}[Virasoro realization of the affine projective tensor subcategory]
\label{cor:projective-essential-image}
Let $\mathcal Q_{p,q}$ denote the smallest replete full subcategory of
$\Oc_{c_{p,q}}$ that is closed under finite direct sums and direct summands
and contains
\[
 \{K_r:1\le r<p\}
 \cup\{K_{np}:n\ge1\}
 \cup\{P(\tau_{n,r}):n\ge1,\ 1\le r<p\}.
\]
Then
\[
 F|_{\mathcal P^k}:\mathcal P^k\xrightarrow{\ \sim\ }\mathcal Q_{p,q}
\]
is a $\C$-linear equivalence.  The subcategory $\mathcal Q_{p,q}$ is
closed under the fusion product of its objects, and with the inherited
constraints this equivalence is braided strong monoidal.  Thus the
root-of-unity projective/tilting tensor category is realized as a full
replete idempotent-complete additive monoidal subcategory of Virasoro
representation theory; the subcategory contains both ordinary and
logarithmic objects.  No claim is made that $\mathcal Q_{p,q}$ is closed
under kernels, cokernels, or extensions in $\Oc_{c_{p,q}}$.
\end{cor}

\begin{proof}
Let $\operatorname{Im}^{\mathrm{repl}}(F|_{\mathcal P^k})$ denote the
replete full essential image of the projective restriction.  By
Theorem~\ref{thm:F-fully-faithful},
\[
 F|_{\mathcal P^k}:\mathcal P^k
 \longrightarrow \operatorname{Im}^{\mathrm{repl}}(F|_{\mathcal P^k})
\]
is a $\C$-linear equivalence, and Theorem~\ref{thm:projective-classification}
identifies the indecomposable isomorphism classes occurring in this image
with the displayed list.  It remains to verify closure under direct
summands.  Indeed, full faithfulness gives an algebra isomorphism
\[
 \End_{\KL^k}(P)\xrightarrow{\sim}\End_{\Oc}(F(P)),
\]
so every idempotent $e$ of $F(P)$ has a unique idempotent lift
$\widetilde e$ in $\End_{\KL^k}(P)$.  Since $\KL^k(\mathfrak{sl}_2)$ is
abelian and direct summands of projective objects are projective,
$\widetilde e$ splits inside $\mathcal P^k$; applying $F$ realizes the
retract defined by $e$ as the image of a projective object.  Hence the
replete full essential image is precisely $\mathcal Q_{p,q}$.

By \cite[Theorem~6.6]{McRaeYang}, $\mathcal P^k$ is monoidal, and $F$ is
a braided tensor functor by \cite[Theorem~7.15]{McRaeYang}.  Its tensor
unit belongs to the image:
$\mathbf1=K_1\cong F(P_1)$.  Thus, for any $X,Y\in\mathcal Q_{p,q}$, there exist projectives $P,Q$
with $X\cong F(P)$ and $Y\cong F(Q)$, and
\[
 X\boxtimes Y\cong F(P)\boxtimes F(Q)
 \cong F(P\boxtimes Q)\in\mathcal Q_{p,q}.
\]
The associativity, unit, and braiding constraints are those inherited
from $\Oc_{c_{p,q}}$ and restrict to $\mathcal Q_{p,q}$ because the
subcategory is full and fusion closed.  Since $F|_{\mathcal P^k}$ is fully faithful and essentially surjective
onto $\mathcal Q_{p,q}$, while its tensor and braiding constraints are the
restrictions of those of $F$, it is a braided strong monoidal equivalence
$\mathcal P^k\xrightarrow{\sim}\mathcal Q_{p,q}$.
\end{proof}

\section{Comparison with Drinfeld--Sokolov reduction}\label{sec:comparison}

We now compare $F$ with the homological functor $H$.  The comparison on
projectives uses projective faithfulness from Section~\ref{sec:images} and
the projective Hom-space calculation of Section~\ref{sec:fullness},
together with the exactness and faithfulness of $H$.  For each non-wall
indecomposable projective
\[
 P=P_{np+r},\qquad n\ge1,\quad1\le r\le p-1,
\]
Corollary~\ref{cor:target-conformal-radical} identifies the image of the
canonical affine conformal nilpotent with the intrinsic nilpotent part of
the Virasoro Hamiltonian:
\[
 F(D_P)=D_{F(P)}.
\]
The corresponding equality $H(D_P)=D_{H(P)}$ will be proved below from
Proposition~\ref{prop:H-twist}.  The exponential identities identify the
canonical conformal nilpotent exactly under both functors: they fix a
distinguished nonzero element of the one-dimensional radical, rather than
merely its line.  No normalization of the auxiliary generator $N_P$ is
involved.  This exact identification removes the scalar ambiguity in the
reverse adjacent arrows.  The right exactness of $F$ from
\cite[Theorem~7.15]{McRaeYang} enters only in
Proposition~\ref{prop:right-exact-extension}, after naturality on
projectives has been established.

\subsection{Projective objects under reduction}

Since $H$ is exact, apply it to \eqref{eq:source-projective-filtration}.
Here exactness of $H$, proved independently in Section~\ref{sec:DS}, is
essential; no exactness of $F$ is being used.
After choosing the endpoint identifications
$H(V_{np-r})\cong K_{np-r}$ and $H(V_{np+r})\cong K_{np+r}$ from
Proposition~\ref{prop:H-Weyl}, we obtain, for $n\ge1$ and
$1\le r\le p-1$,
\begin{equation}\label{eq:H-projective-filtration}
 0\longrightarrow K_{np-r}\longrightarrow H(P_{np+r})
 \longrightarrow K_{np+r}\longrightarrow0.
\end{equation}
Put
\[
 a=np-r,\qquad b=np+r,\qquad c=(n+2)p-r.
\]
By Proposition~\ref{prop:Felder-labels},
\[
 h_a<h_b<h_c.
\]
The one-row Kac sequences are
\[
 0\longrightarrow S_b\longrightarrow K_a\longrightarrow S_a
 \longrightarrow0,
 \qquad
 0\longrightarrow S_c\longrightarrow K_b\longrightarrow S_b
 \longrightarrow0.
\]
Thus the two simple constituents of each endpoint Kac module are
nonisomorphic, and Lemma~\ref{lem:length-two-Hom} gives
\[
 \End(K_a)=\End(K_b)=\C.
\]
Changing either endpoint identification in
\eqref{eq:H-projective-filtration} therefore rescales the associated
Yoneda class by a nonzero scalar.  In particular, its vanishing or
nonvanishing is independent of these choices.

\begin{prop}\label{prop:projective-object-comparison}
For every indecomposable projective $P_j$, there is an isomorphism
\[
 F(P_j)\cong H(P_j).
\]
More precisely,
\[
 F(P_j)\cong H(P_j)\cong
 \begin{cases}
  K_j,&1\le j\le p-1,\\
  K_j=S_j,&p\mid j,\\
  P(\tau_{n,r}),&j=np+r,\ n\ge1,\ 1\le r\le p-1.
 \end{cases}
\]
In the third case both $F(P_j)$ and $H(P_j)$ are logarithmic.
No compatibility of these objectwise isomorphisms with source morphisms is
asserted at this stage.
\end{prop}

\begin{proof}
If $1\le j\le p-1$, then $P_j=V_j$ and both functors give $K_j$ by
Lemma~\ref{lem:F-small-Weyl} and Proposition~\ref{prop:H-Weyl}.  If
$p\mid j$, then $P_j=V_j=L_j$; hence
\[
 F(P_j)\cong K_j,
 \qquad
 H(P_j)\cong H(V_j)\cong K_j=S_j,
\]
by Proposition~\ref{prop:projective-images} and
Proposition~\ref{prop:H-Weyl}.

Let $j=b=np+r$ with $n\ge1$ and $1\le r\le p-1$, and set
$E=H(P_b)$.  Both end terms in \eqref{eq:H-projective-filtration} belong to
$\mathscr C$ by Lemma~\ref{lem:Nakano-membership}.  Since $\mathscr C$
is Serre, the middle term also satisfies $E\in\mathscr C$.  By Lemma~\ref{lem:source-L0-radical}, write
$D_{P_b}=\delta_{P_b}N_{P_b}$ with $\delta_{P_b}\ne0$.  Proposition~\ref{prop:H-twist}
then gives the precise identity
\begin{equation}\label{eq:H-projective-exponential}
\begin{aligned}
 e^{2\pi iL^{DS}_0}\big|_E
 &=H(\theta^-_{P_b})\\
 &=\lambda_{P_b}\bigl(\id_E+2\pi iH(D_{P_b})\bigr)\\
 &=\lambda_{P_b}\bigl(
    \id_E+2\pi i\delta_{P_b}H(N_{P_b})
   \bigr).
\end{aligned}
\end{equation}
Since $H$ is faithful,
\[
 U:=2\pi i\delta_{P_b}H(N_{P_b})\ne0,\qquad U^2=0.
\]
Thus the operator $T:=\lambda_{P_b}^{-1}e^{2\pi iL^{DS}_0}|_E=
\id+U$ satisfies
\[
 (T-\id)^2=0,\qquad T\ne\id.
\]
Its minimal polynomial therefore contains $(X-1)^2$, so
$e^{2\pi iL^{DS}_0}|_E$ is not diagonalizable.
Lemma~\ref{lem:exponential-logarithmic} implies that $E$ is logarithmic.
Since both $K_a$ and $K_b$ are ordinary, the exact sequence
\eqref{eq:H-projective-filtration} cannot split: otherwise
$E\cong K_a\oplus K_b$ would have semisimple $L_0$.  Thus $E$
represents a nonzero logarithmic extension class.  We now verify
explicitly that Lemma~\ref{lem:Nakano-recognition} applies.  The crucial input, Proposition~\ref{prop:one-row-staggered-quotient},
is a statement about an \emph{arbitrary} logarithmic middle term with these
one-row endpoints.  For $a\ge2$, its proof first verifies every condition
of \cite[Definitions~3.14 and~5.1]{Nakano} for that very middle term and
then invokes the universal quantifier in \cite[Theorem~5.11]{Nakano},
which is stated for an arbitrary logarithmic representative of the
relevant Ext class.  The exceptional case $a=1$ is proved separately in
Lemma~\ref{lem:vacuum-edge-staggered}, with the two prime branches of the
right braid-type Verma module treated according to their actual grades.
Thus no distinguished deformation is being substituted for the module
$E$ constructed here, and no blanket identification of ambient categories
is used.  By exactness of $H$,
\[
 0\longrightarrow K_a\longrightarrow E\longrightarrow K_b
 \longrightarrow0,
 \qquad a=np-r,\quad b=np+r,
\]
is exact with precisely the submodule--quotient orientation required
there.  By Proposition~\ref{prop:Felder-labels},
\[
 \tau_{n,r}\in
 \mathcal T_{p_+,p_-}\setminus\mathcal T^0_{p_+,p_-},
\]
and Lemma~\ref{lem:Nakano-membership}, together with the Serre property,
places all three terms in $\mathscr C$.  Since $E$ is logarithmic, every
hypothesis of Lemma~\ref{lem:Nakano-recognition} is satisfied.  Hence
\[
 H(P_b)=E\cong P(\tau_{n,r}).
\]
Corollary~\ref{cor:Q-Ptau} gives the same identification for $F(P_b)$.
\end{proof}

\begin{cor}[Canonical conformal radical under reduction]
\label{cor:H-projective-radical}
Let $P=P_{np+r}$ with $n\ge1$ and $1\le r\le p-1$, and set $E=H(P)$.
Let $D_P$ be the canonical affine conformal nilpotent of
Lemma~\ref{lem:source-L0-radical}.  For each actual generalized
$L^{DS}_0$-eigenvalue $\lambda$ of $E$, define
\[
 D_E|_{E^{\mathrm{gen}}_\lambda}
 :=(L^{DS}_0-\lambda\id)|_{E^{\mathrm{gen}}_\lambda}.
\]
Then the operators $D_E$ assemble to a Virasoro-module endomorphism and
\begin{equation}\label{eq:H-canonical-radical}
 H(D_P)=D_E,
 \qquad
 \rad\End H(P)=\C D_E=\C H(N_P).
\end{equation}
In particular,
\[
 \End H(P)=\C\id\oplus\C D_E
 \cong\C[\varepsilon]/(\varepsilon^2).
\]
\end{cor}

\begin{proof}
Put $\xi_E=[h_{np-r}]=[h_{np+r}]$.  Both end terms in
\eqref{eq:H-projective-filtration} are supported in the residue sector
$\xi_E$.  For every $\eta\ne\xi_E$, exact residue projection gives
\[
 0\longrightarrow0\longrightarrow E^{[\eta]}\longrightarrow0,
\]
so $E=E^{[\xi_E]}$.  On actual generalized $L^{DS}_0$-eigenspaces, the
Virasoro commutator $[L^{DS}_0,L_m]=-mL_m$ gives
\[
 D_{\lambda-m}L_m=L_mD_\lambda,
\]
so the operators above assemble to a Virasoro-module endomorphism $D_E$;
indeed, $L_m$ maps the generalized eigenspace of eigenvalue $\lambda$
into that of eigenvalue $\lambda-m$.  The module $E$ is logarithmic by
Proposition~\ref{prop:projective-object-comparison}, so $D_E\ne0$.

At this point Proposition~\ref{prop:projective-object-comparison} is used
only as an objectwise identification.  The following argument imports
only the already established target endomorphism algebra from
Corollary~\ref{cor:target-local-endomorphisms}; it neither uses nor
presupposes a natural comparison between $F$ and $H$.  No choice made here
enters the naturality construction below.  Choose an arbitrary
Virasoro-module isomorphism $E\xrightarrow{\sim}F(P)$.  At this point it
is used only to transport the dimension and locality of the endomorphism
algebra; no radical generator and no naturality datum is transported
through it.  Conjugation and
Corollary~\ref{cor:target-local-endomorphisms} give
\[
 \dim\End(E)=2,\qquad \End(E)\text{ is local}.
\]
Hence $\dim\rad\End(E)=1$.  Every vector of $E$ is a finite sum of generalized
$L^{DS}_0$-eigenvectors, and $D_E$ is nilpotent on each such generalized
eigenspace.  Thus $D_E$ is locally nilpotent and nonzero, hence is not
invertible.  Since $\End(E)$ is a finite-dimensional local algebra, its
nonunits are precisely its Jacobson radical.  Therefore
\[
 0\ne D_E\in\rad\End(E),
\]
and the one-dimensionality of the radical gives
$\rad\End(E)=\C D_E$.  Write $D_E^2=aD_E$.  The Jacobson radical of the
finite-dimensional algebra $\End(E)$ is nilpotent; if $a\ne0$, then
$D_E^m=a^{m-1}D_E\ne0$ for every $m\ge1$, a contradiction.  Hence
$a=0$ and
\[
 D_E^2=0.
\]

On the source,
\[
 \theta^-_P=\lambda_P(\id+2\pi iD_P).
\]
Applying $H$ and using Proposition~\ref{prop:H-twist} gives
\begin{equation}\label{eq:H-DP-exponential}
 e^{2\pi iL^{DS}_0}|_E
 =\lambda_P\bigl(\id+2\pi iH(D_P)\bigr).
\end{equation}
Faithfulness of $H$ gives $H(D_P)\ne0$, while $H(D_P)^2=0$.  Hence
$H(D_P)$ is a nonunit of the same local algebra and therefore belongs to
its radical.  On the other hand, every actual generalized eigenvalue $\lambda$ of
$E$ lies in the single residue class $\xi_E$.  Choose any representative
$\lambda_0$ of this class and put $\mu_E:=e^{2\pi i\lambda_0}$; this is
independent of the representative.  Since $D_E^2=0$, the Jordan decomposition on
each generalized eigenspace gives the global identity
\[
 e^{2\pi iL^{DS}_0}|_E=\mu_E(\id+2\pi iD_E).
\]
Because Virasoro modes shift conformal weights by integers, this
exponential commutes with every Virasoro mode and hence belongs to
$\End(E)$.  Thus the same endomorphism $e^{2\pi iL^{DS}_0}|_E$ admits the
two expressions
\[
 \lambda_P\bigl(\id+2\pi iH(D_P)\bigr)
 \qquad\text{and}\qquad
 \mu_E\bigl(\id+2\pi iD_E\bigr).
\]
Since both $H(D_P)$ and $D_E$ belong to $\rad\End(E)$, subtracting
the two exponential expressions shows
\[
 (\lambda_P-\mu_E)\id_E\in\rad\End(E).
\]
But $\id_E\notin\rad\End(E)$, so $\lambda_P=\mu_E$.  This scalar
is nonzero; substituting back gives
\[
 H(D_P)=D_E.
\]
Finally Lemma~\ref{lem:source-L0-radical} gives
$D_P=\delta_PN_P$ for some $\delta_P\in\C^\times$.  Since $H$ is
faithful,
\[
 \C H(N_P)=\C H(D_P)=\C D_E=\rad\End H(P),
\]
which is the remaining equality in \eqref{eq:H-canonical-radical}.
\end{proof}

\begin{rmk}\label{rmk:objectwise-not-natural}
Proposition~\ref{prop:projective-object-comparison} is only an objectwise
statement.  No canonical choice of the middle isomorphisms, and no
compatibility with projective morphisms, is asserted at this point.  The
next subsection supplies precisely this compatibility.
\end{rmk}

\subsection{Canonical conformal radicals and naturality}

The exponential will not make the individual objectwise isomorphisms
canonical.  Its role is narrower and decisive: it removes the scalar
obstruction to satisfying naturality simultaneously for the two
orientations of every adjacent edge.

\begin{lem}[Intrinsic conformal-radical compatibility]
\label{lem:twist-locking}
Let $P=P_{np+r}$ with $n\ge1$ and $1\le r\le p-1$.  Every
Virasoro-module isomorphism
\[
 \eta_P:F(P)\xrightarrow{\sim}H(P)
\]
satisfies
\begin{equation}\label{eq:twist-locking}
 \eta_PF(u)=H(u)\eta_P
 \qquad\text{for every }u\in\rad\End(P).
\end{equation}
\end{lem}

\begin{proof}
The exact normalization is essential here: knowing only that the two
radicals are one-dimensional would determine the desired identity only up
to a nonzero scalar, which would not suffice to fix the reverse-edge
obstruction below.  Let $D_P$ be the canonical affine conformal nilpotent.
By Corollary~\ref{cor:target-conformal-radical}, $F(D_P)$ is exactly the
nilpotent part of the target Virasoro Hamiltonian on $F(P)$, while
Corollary~\ref{cor:H-projective-radical} identifies $H(D_P)$ with the
corresponding nilpotent part on $H(P)$.  A Virasoro-module isomorphism commutes with $L_0$, so for every actual
generalized eigenvalue $\lambda$ one has
\[
 \eta_P\bigl(F(P)^{\mathrm{gen}}_\lambda\bigr)
 =H(P)^{\mathrm{gen}}_\lambda,
 \qquad
 \eta_P(L_0-\lambda)=(L_0-\lambda)\eta_P
\]
on that generalized eigenspace.  Hence it intertwines the nilpotent parts
and
\[
 \eta_PF(D_P)=H(D_P)\eta_P.
\]
Since $\rad\End(P)=\C D_P$, every $u\in\rad\End(P)$ has the form
$u=aD_P$ for some $a\in\C$, and \eqref{eq:twist-locking} follows by
linearity.  No monoidal structure on $H$ is used.
\end{proof}

Thus radical-loop naturality is automatic for every objectwise
Virasoro-module isomorphism.  In constructing a natural comparison on a
reflection chain, only one orientation of each adjacent edge needs to be
normalized; the reverse orientation will then be forced.

\begin{lem}[Mixed projective Hom dimensions]
\label{lem:mixed-projective-Hom}
Fix a non-wall reflection chain $\{s_i\}_{i\ge0}$ and put
$\Pi_i=P_{s_i}$.  For any $\mathscr G,\mathscr G'\in\{F,H\}$,
\[
 \dim\Hom(\mathscr G(\Pi_i),\mathscr G'(\Pi_j))
 =
 \begin{cases}
  2,&i=j\ge1,\\
  1,&i=j=0,\\
  1,&|i-j|=1,\\
  0,&|i-j|\ge2.
 \end{cases}
\]
\end{lem}

\begin{proof}
Choose objectwise isomorphisms
$\phi_j:H(\Pi_j)\xrightarrow{\sim}F(\Pi_j)$ from
Proposition~\ref{prop:projective-object-comparison}.  For each choice $\mathscr G,\mathscr G'\in\{F,H\}$,
pre- and postcomposition with the appropriate $\phi_i^{\pm1}$ and
$\phi_j^{\pm1}$ gives a vector-space isomorphism
\[
 \Hom(\mathscr G(\Pi_i),\mathscr G'(\Pi_j))
 \cong \Hom(F(\Pi_i),F(\Pi_j)).
\]
The latter dimension is given by Proposition~\ref{prop:target-Hom-table}.
No naturality of the auxiliary isomorphisms $\phi_j$ is used here; they
serve only to compare the underlying vector spaces.
\end{proof}

\begin{lem}[Projective exhaustion]\label{lem:projective-exhaustion}
Every indecomposable projective object of $\KL^k(\mathfrak{sl}_2)$ is
isomorphic to exactly one $P_j$.  Consequently
\[
 \mathcal P^k=\operatorname{add}\{P_j:j\ge1\}.
\]
\end{lem}

\begin{proof}
Since $R$ has finite length, a nonzero indecomposable projective $R$
has a maximal proper subobject; choose a simple quotient
$q:R\twoheadrightarrow L_j$.  Let
$p_j:P_j\twoheadrightarrow L_j$ be the projective cover.  Projectivity
gives morphisms $a:P_j\to R$ and $b:R\to P_j$ satisfying
$q a=p_j$ and $p_j b=q$.  Hence $p_jba=p_j$.  Since $p_j$ is essential,
$ba$ is an epimorphism; as $P_j$ has finite length, $ba$ is an
automorphism.  Thus $a$ is a split monomorphism, so $P_j$ is a nonzero
direct summand of $R$.  Indecomposability gives $R\cong P_j$.  The label
is unique because the simple heads are pairwise nonisomorphic: the lowest
affine conformal-degree subspace of $L_j$ is the $j$-dimensional
irreducible horizontal $\mathfrak{sl}_2$-module.  Hence
$L_j\cong L_\ell$ forces $j=\ell$.  Finally, induction on length decomposes
every projective into a finite direct sum of indecomposable projectives,
proving the last assertion.
\end{proof}

\begin{prop}\label{prop:natural-on-projectives}
There is a natural isomorphism of $\C$-linear functors
\[
 \eta:F|_{\mathcal P^k}\xRightarrow{\ \sim\ }H|_{\mathcal P^k}.
\]
On each non-wall reflection chain the components can be chosen
inductively.  Exponential compatibility forces the scalar obstruction for
the reverse adjacent morphisms to be trivial; no uniqueness of the
individual components is asserted.
\end{prop}

\begin{proof}
Fix a non-wall reflection chain $\{s_i\}_{i\ge0}$, put
$\Pi_i=P_{s_i}$, and choose adjacent generators
\[
 \Pi_i\underset{y_i}{\stackrel{x_i}{\rightleftarrows}}\Pi_{i+1}
\]
as nonzero generators of the two one-dimensional adjacent Hom spaces in
Proposition~\ref{prop:source-Hom-package}, normalized only so that
\begin{equation}\label{eq:zigzag-loop-section6}
 \nu_{i+1}:=x_i\circ y_i\ne0
\end{equation}
spans $\rad\End(\Pi_{i+1})$.

Choose an isomorphism $\eta_0:F(\Pi_0)\to H(\Pi_0)$.  Suppose $\eta_i$
has been chosen and take any isomorphism
$\eta'_{i+1}:F(\Pi_{i+1})\to H(\Pi_{i+1})$.  By
Lemma~\ref{lem:mixed-projective-Hom}, both
$\eta'_{i+1}\circ F(x_i)$ and $H(x_i)\circ\eta_i$ lie in the same one-dimensional
mixed Hom space.  They are nonzero: $F(x_i)\ne0$ by projective faithfulness of $F$,
$H(x_i)\ne0$ by faithfulness of $H$, and both $\eta_i$ and
$\eta'_{i+1}$ are isomorphisms.  Hence there is a unique
$a_i\in\C^\times$ such that
\[
 \eta'_{i+1}\circ F(x_i)=a_i H(x_i)\circ\eta_i.
\]
Set $\eta_{i+1}:=a_i^{-1}\eta'_{i+1}$.  Then
\begin{equation}\label{eq:forward-naturality}
 \eta_{i+1}\circ F(x_i)=H(x_i)\circ\eta_i.
\end{equation}
The reverse mixed Hom space is also one-dimensional.  Both
$\eta_i\circ F(y_i)$ and $H(y_i)\circ\eta_{i+1}$ are nonzero, by
faithfulness of $F$ on projectives, faithfulness of $H$, and invertibility
of the chosen components.  Hence
\begin{equation}\label{eq:reverse-scalar}
 \eta_i\circ F(y_i)=c_iH(y_i)\circ\eta_{i+1}
\end{equation}
for a unique $c_i\in\C^\times$.  Composing with the forward relation gives
\begin{align*}
 \eta_{i+1}\circ F(\nu_{i+1})
 &=\eta_{i+1}\circ F(x_i)\circ F(y_i)\\
 &=H(x_i)\circ\eta_i\circ F(y_i)\\
 &=c_iH(x_i)\circ H(y_i)\circ\eta_{i+1}\\
 &=c_iH(\nu_{i+1})\circ\eta_{i+1}.
\end{align*}
Lemma~\ref{lem:twist-locking} gives
$\eta_{i+1}\circ F(\nu_{i+1})=
H(\nu_{i+1})\circ\eta_{i+1}$.  Subtracting from the
preceding equality gives
\[
 (c_i-1)H(\nu_{i+1})\circ\eta_{i+1}=0.
\]
Since $\nu_{i+1}\ne0$, faithfulness of $H$ gives
$H(\nu_{i+1})\ne0$, and hence
$H(\nu_{i+1})\circ\eta_{i+1}\ne0$.  Therefore $c_i=1$.  Thus the normalization along $x_i$ automatically
enforces naturality along the reverse arrow $y_i$ as well.  This argument
uses only the fixed nonzero loop $\nu_{i+1}=x_i y_i$; no assertion about
the opposite composite $y_i x_i$ is required.

By Proposition~\ref{prop:source-Hom-package}, the displayed elements form
bases of every nonzero Hom space between indecomposable projectives on the
fixed reflection chain:
\[
 \C\id_{\Pi_0},\qquad \C x_i,\qquad \C y_i,
 \qquad \C\id_{\Pi_i}\oplus\C\nu_i\ (i\ge1),
\]
and all remaining Hom spaces vanish.  Naturality has been checked on these
bases: for $x_i,y_i$ above, for $\nu_i$ by
Lemma~\ref{lem:twist-locking}, and for identities automatically.  Linearity
therefore gives naturality for every morphism on the chain.  Hom spaces between distinct reflection chains vanish by
Proposition~\ref{prop:source-Hom-package}.

For every wall projective $P_{np}$, choose any isomorphism
\[
 \eta_{np}:F(P_{np})\xrightarrow{\sim}H(P_{np}),
\]
using the identifications
$F(P_{np})\cong K_{np}\cong H(P_{np})$.  Wall projectives are isolated and
have scalar endomorphism algebras, so naturality there is automatic.
By Lemma~\ref{lem:projective-exhaustion}, these components cover all
indecomposable projective objects.

For $R=0$ take the unique isomorphism $0\to0$.  For each
indecomposable $P_j$ set $\alpha_{P_j}=\id_{P_j}$.  For every other
nonzero projective $R$, fix once and for all an isomorphism
\[
 \alpha_R:R\xrightarrow{\sim}
 \bigoplus_j P_j^{\oplus m_j(R)}
\]
with only finitely many nonzero multiplicities.  No independence of this
choice is asserted or needed.  Since $F$ and $H$ are additive, we use their canonical finite-biproduct identifications
\[
 F\!\left(\bigoplus_jP_j^{\oplus m_j(R)}\right)
 \cong\bigoplus_jF(P_j)^{\oplus m_j(R)},
 \qquad
 H\!\left(\bigoplus_jP_j^{\oplus m_j(R)}\right)
 \cong\bigoplus_jH(P_j)^{\oplus m_j(R)}.
\]
With these identifications understood, define
\[
 \eta_R:=
 H(\alpha_R)^{-1}
 \left(\bigoplus_j\eta_{P_j}^{\oplus m_j(R)}\right)
 F(\alpha_R).
\]
No canonicity or independence of the chosen Krull--Schmidt isomorphisms is
asserted.  Fix these choices.  We now verify that the resulting family is
natural.  If $f:R\to R'$ is any projective morphism, the matrix of
$\alpha_{R'}f\alpha_R^{-1}$ has entries between indecomposable
projectives, on which naturality has already been proved.  Matrix
multiplication therefore gives
$\eta_{R'}F(f)=H(f)\eta_R$.
\end{proof}

\begin{lem}[Radical-unipotent ambiguity]\label{lem:radical-annihilates-arrows}
With the notation of Lemma~\ref{lem:mixed-projective-Hom}, let
$\nu_i$ span the radical of $\End(\Pi_i)$ for $i\ge1$.  For every
adjacent generator $x_i:\Pi_i\to\Pi_{i+1}$ and
$y_i:\Pi_{i+1}\to\Pi_i$ one has, whenever the indicated radical is
defined,
\[
 H(\nu_{i+1})H(x_i)=0,\qquad H(x_i)H(\nu_i)=0,
\]
\[
 H(\nu_i)H(y_i)=0,\qquad H(y_i)H(\nu_{i+1})=0.
\]
Consequently, suppose isomorphisms
$\eta_j:F(\Pi_j)\xrightarrow{\sim}H(\Pi_j)$ have been chosen and satisfy
some collection of adjacent naturality relations.  At any positive vertex
$i$, replacing $\eta_i$ by
$(\id+aH(\nu_i))\circ\eta_i$ preserves every adjacent naturality equation
already satisfied.  It also preserves naturality for the radical loop
$\nu_i$, because $H(\nu_i)^2=0$.
\end{lem}

\begin{proof}
Postcomposition defines a linear endomorphism
\[
 T:\Hom(H(\Pi_i),H(\Pi_{i+1}))\longrightarrow
   \Hom(H(\Pi_i),H(\Pi_{i+1})),
 \qquad T(f)=H(\nu_{i+1})f.
\]
By Proposition~\ref{prop:source-Hom-package} and
Lemma~\ref{lem:source-L0-radical},
\[
 \rad\End(\Pi_{i+1})=\C N_{s_{i+1}},
 \qquad N_{s_{i+1}}^2=0.
\]
Since $\nu_{i+1}$ is a nonzero element of this radical,
$\nu_{i+1}^2=0$.  Hence $T^2=0$.  The Hom space is one-dimensional by Lemma~\ref{lem:mixed-projective-Hom};
therefore $T$ is scalar, and the only square-zero scalar is zero.  Hence
$T=0$.  The other three
statements follow identically from pre- or postcomposition on the
corresponding one-dimensional adjacent Hom space.  Since
$(\id+aH(\nu_i))^{-1}=\id-aH(\nu_i)$, these are radical-unipotent
automorphisms and the final assertion follows.  More generally, for any
family of scalars $(a_i)_{i\ge1}$, the simultaneous replacements
\[
 \eta_i\longmapsto(\id+a_iH(\nu_i))\circ\eta_i
 \qquad(i\ge1)
\]
preserve all adjacent naturality equations: on each edge the two possible
correction terms vanish by the four identities above.  They also preserve
naturality for every radical loop, since the additional correction is a
multiple of $H(\nu_i)^2=0$.  Consequently, if the original family is natural on the full projective
reflection chain, then simultaneous independent choices of the $a_i$
again give a natural isomorphism on that chain.
\end{proof}

\subsection{Extension from projectives}

The following elementary projective-density statement is the
natural-transformation analogue of the right-exact extension principle
used in \cite[Theorem~6.7]{McRaeYang}.  We include the proof because the
uniqueness of the extended natural transformation will be used explicitly.

\begin{prop}[Projective determination of natural transformations between right exact functors]
\label{prop:right-exact-extension}
Let $\mathcal A$ be a $\C$-linear abelian category with enough projectives,
let $\mathcal P\subset\mathcal A$ be its full projective subcategory, and
let $\mathcal B$ be a $\C$-linear abelian category.  Here ``right exact''
means that an exact sequence $A\to B\to C\to0$ is sent to an exact sequence
$G(A)\to G(B)\to G(C)\to0$.  For right exact additive $\C$-linear functors $G_1,G_2:\mathcal A\to\mathcal B$, restriction
induces a bijection
\[
 \operatorname{Nat}(G_1,G_2)
 \xrightarrow{\ \sim\ }
 \operatorname{Nat}(G_1|_{\mathcal P},G_2|_{\mathcal P}).
\]
Thus a natural transformation between right exact functors is uniquely
determined by, and can be reconstructed from, its projective restriction.
This is the fully faithful form of the right-exact extension principle;
compare \cite[Theorem~6.7]{McRaeYang}.  In particular, every natural
isomorphism on projectives extends uniquely
to a natural isomorphism on $\mathcal A$.
\end{prop}

\begin{proof}
Enough projectives is sufficient here; no finite global-dimension
hypothesis is needed.  Restriction is injective: if
$\theta:G_1\Rightarrow G_2$ vanishes on projectives and
$p:P\twoheadrightarrow M$ is a projective epimorphism, then
\[
 \theta_MG_1(p)=G_2(p)\theta_P=0.
\]
Right exactness makes $G_1(p)$ epic, so $\theta_M=0$.

For surjectivity, choose a projective epimorphism
$p:P_0\twoheadrightarrow M$ and then a projective epimorphism
$u:P_1\twoheadrightarrow\ker p$.  Let
$d:P_1\to P_0$ be the composite of $u$ with the inclusion
$\ker p\hookrightarrow P_0$.  Thus
\[
 P_1\xrightarrow d P_0\xrightarrow p M\longrightarrow0
\]
is a projective presentation in the required sense.  The map $d$ need
not be injective; no projective resolution or finite global-dimension
hypothesis is being assumed.

Let $\eta^{\mathcal P}$ be a natural transformation on projectives.
Right exactness gives
\[
 G_1(P_1)\xrightarrow{G_1(d)}G_1(P_0)
 \xrightarrow{G_1(p)}G_1(M)\to0,
\]
and similarly for $G_2$.  In particular, $G_1(p)$ is epic and exhibits
$G_1(M)$ as a cokernel of $G_1(d)$.  Naturality for $d$ gives the
explicit identity
\[
 G_2(p)\eta^{\mathcal P}_{P_0}G_1(d)
 =G_2(p)G_2(d)\eta^{\mathcal P}_{P_1}=0.
\]
The cokernel universal property therefore gives a unique map
\[
 \eta_M:G_1(M)\longrightarrow G_2(M)
\]
such that
\begin{equation}\label{eq:eta-projective-presentation}
 \eta_MG_1(p)=G_2(p)\eta^{\mathcal P}_{P_0}.
\end{equation}
The morphism $\eta_M$ is therefore determined by the projective
epimorphism $p:P_0\twoheadrightarrow M$ alone; the object $P_1$ is used
only to verify that the right-hand side of
\eqref{eq:eta-projective-presentation} annihilates
$\operatorname{im}G_1(d)$.  In particular the construction is independent
of the choice of $P_1\twoheadrightarrow\ker p$.

The map is also independent of the chosen projective epimorphism $p$.  Indeed, let
$p':P'_0\twoheadrightarrow M$ be another projective epimorphism and let
$\eta'_M$ be the map constructed from it.  Projectivity of $P_0$ gives
$a:P_0\to P'_0$ with $p'a=p$.  Then
\begin{align*}
 \eta'_M G_1(p)
 &=\eta'_M G_1(p')G_1(a)\\
 &=G_2(p')\eta^{\mathcal P}_{P'_0}G_1(a)\\
 &=G_2(p')G_2(a)\eta^{\mathcal P}_{P_0}\\
 &=G_2(p)\eta^{\mathcal P}_{P_0}\\
 &=\eta_MG_1(p).
\end{align*}
Since $G_1(p)$ is epic, $\eta'_M=\eta_M$.  A lift in the opposite
direction is unnecessary: equality after precomposition with one
projective epimorphism already determines a morphism out of $G_1(M)$.

For a morphism $f:M\to N$, choose projective epimorphisms
$p:P\twoheadrightarrow M$ and $q:Q\twoheadrightarrow N$.  Projectivity of
$P$ gives $a:P\to Q$ with $qa=fp$.  Using
\eqref{eq:eta-projective-presentation} twice and naturality of
$\eta^{\mathcal P}$ for $a$, we compute
\begin{align*}
 G_2(f)\eta_MG_1(p)
 &=G_2(f)G_2(p)\eta^{\mathcal P}_P\\
 &=G_2(q)G_2(a)\eta^{\mathcal P}_P\\
 &=G_2(q)\eta^{\mathcal P}_QG_1(a)\\
 &=\eta_NG_1(q)G_1(a)\\
 &=\eta_NG_1(f)G_1(p).
\end{align*}
Since $G_1(p)$ is epic, $G_2(f)\eta_M=\eta_NG_1(f)$.  Thus $\eta$ is
natural.  Taking $p=\id_P$ and $P_1=0$ shows that its restriction is
$\eta^{\mathcal P}$, and \eqref{eq:eta-projective-presentation} shows
uniqueness.  If $\eta^{\mathcal P}$ is invertible, apply the same construction to
$(\eta^{\mathcal P})^{-1}$, obtaining $\zeta:G_2\Rightarrow G_1$.
The composites $\zeta\eta$ and $\eta\zeta$ restrict to the identity
natural transformations on projectives.  By the uniqueness just proved,
they are the identity transformations on all of $\mathcal A$.  Hence the
extension of an invertible projective transformation is a natural
isomorphism.
\end{proof}

\subsection{Main theorem and consequences}

\begin{proof}[Proof of Theorem~\ref{thm:main}]
At this point exactness of $F$ has nowhere been used.  The comparison has
three steps:
\[
 F|_{\mathcal P^k}\cong H|_{\mathcal P^k}
 \quad\Longrightarrow\quad
 F\cong H
 \quad\Longrightarrow\quad
 F\text{ is exact and faithful}.
\]
By this stage, the projective comparison has been obtained using the full
faithfulness of $F|_{\mathcal P^k}$, the exactness and faithfulness of
$H$, the objectwise projective comparison, and the two exponential
compatibilities.  No preservation of kernels or monomorphisms by $F$ has
been used.  The first arrow is Proposition~\ref{prop:natural-on-projectives}.
Only now do we invoke right exactness of $F$, supplied by
\cite[Theorem~7.15]{McRaeYang}.  The category
$\KL^k(\mathfrak{sl}_2)$ is a $\C$-linear abelian category with enough
projectives, and $\Oc_{c_{p,q}}$ is a $\C$-linear abelian category.  The
functor $F$ is additive, $\C$-linear, and right exact, while $H$ is
additive and $\C$-linear and is right exact because it is exact.
Therefore every hypothesis of Proposition~\ref{prop:right-exact-extension}
is satisfied.  That proposition gives the unique extension of the chosen
projective comparison,
\begin{equation}\label{eq:main-natural-iso}
 F_{p,q}\xRightarrow{\ \sim\ }
 H^0_{DS,+}\big|_{\KL^k(\mathfrak{sl}_2)}.
\end{equation}
Exactness is invariant under natural isomorphism of additive functors, so
exactness of $H$ implies exactness of $F$.  If $f:M\to N$ and $F(f)=0$,
naturality of \eqref{eq:main-natural-iso} gives
\[
 H(f)=\eta_N\circ F(f)\circ\eta_M^{-1}=0.
\]
Faithfulness of $H$ therefore implies $f=0$, so $F$ is faithful.
Projective full faithfulness is Theorem~\ref{thm:F-fully-faithful}.
No monoidal structure on $H$ is used in this implication.  By
Lemma~\ref{lem:FKW-normalization}, the functor on the right-hand side of
\eqref{eq:main-natural-iso} is the quantized Drinfeld--Sokolov reduction
referred to in \cite[Conjecture~7.16]{McRaeYang}.
\end{proof}



\begin{lem}[Rigid normalization under parameter interchange]
\label{lem:parameter-swap-rigid-normalization}
Put $c=c_{p,q}=c_{q,p}$.  Both notations
$\Oc_{c_{q,p}}$ and $\Oc_{c_{p,q}}$ refer to the vertex tensor category
of finite-length modules for the same universal Virasoro vertex algebra
$V_c^{\mathrm{Vir}}$.  We therefore identify them as the same HLZ
braided tensor category, with the same tensor product, tensor unit,
braiding, and twist.  In this category set
\[
 A:=\mathcal K^{(q,p)}_{2,1},\qquad
 B:=\mathcal K^{(p,q)}_{1,2}.
\]
There is an isomorphism $\phi:A\xrightarrow{\sim}B$ which identifies the
chosen self-duality data of \cite[Theorem~7.15]{McRaeYang}:
\[
 e_B\circ(\phi\boxtimes\phi)=e_A,
 \qquad
 (\phi^{-1}\boxtimes\phi^{-1})\circ i_B=i_A.
\]
\end{lem}

\begin{proof}
By the parameter-interchange identification used in the proof of
\cite[Theorem~7.15]{McRaeYang}, in the braided tensor category fixed in
the statement the Kac module $\mathcal K^{(q,p)}_{2,1}$ corresponds to
$\mathcal K^{(p,q)}_{1,2}$.  Fix a Virasoro-module isomorphism
\[
 \phi_0:A\xrightarrow{\sim}B.
\]
We first record that
\[
 \End(A)=\End(B)=\C.
\]
Indeed, by \cite[Remark~2.2]{McRaeSopin}, both modules are
highest-weight quotients of Virasoro Verma modules and are generated by
their one-dimensional lowest-weight spaces.  A Virasoro endomorphism
preserves the lowest conformal-weight line, hence acts there by a scalar;
since that line generates the module, the endomorphism is the same scalar
multiple of the identity.  This argument includes the boundary case
$q=2$ and does not use a length-two description.

Both $A$ and $B$ are the rigid self-dual Kac objects used in
\cite[Theorem~7.15]{McRaeYang}.  Rigidity and the preceding endomorphism
calculation give
\[
 \dim\Hom(A\boxtimes A,\mathbf1)
 =\dim\Hom(\mathbf1,A\boxtimes A)=1,
\]
and similarly for $B$.  Hence there exist $a,b\in\C^\times$ such that
\[
 e_B\circ(\phi_0\boxtimes\phi_0)=a e_A,
 \qquad
 (\phi_0^{-1}\boxtimes\phi_0^{-1})\circ i_B=b i_A.
\]
Transporting the duality of $B$ across $\phi_0$ gives evaluation
$a e_A$ and coevaluation $b i_A$ on $A$.  Comparison of the two
zig-zag identities gives $ab=1$.  Choose $t\in\C^\times$ with
$t^2=a^{-1}$ and put $\phi=t\phi_0$.  Then
\[
 e_B\circ(\phi\boxtimes\phi)=e_A,
\]
and
\[
 (\phi^{-1}\boxtimes\phi^{-1})\circ i_B
 =t^{-2}b i_A=ab\,i_A=i_A.
\]
Thus $\phi$ identifies the chosen self-duality data.
\end{proof}

\begin{cor}[Transposed reduction]\label{cor:transposed-reduction}
Let $k'=-2+q/p$.  Since $c_{q,p}=c_{p,q}=c$, the two target
notations denote the same HLZ braided tensor category of finite-length
$V_c^{\mathrm{Vir}}$-modules, as fixed in
Lemma~\ref{lem:parameter-swap-rigid-normalization}.  The second
McRae--Yang functor satisfies
\[
 F_{q,p}:\KL^{k'}(\mathfrak{sl}_2)\longrightarrow\Oc_{c_{p,q}},
 \qquad
 F_{q,p}\cong H^0_{DS,+}\big|_{\KL^{k'}(\mathfrak{sl}_2)}
\]
as $\C$-linear functors.  In particular, $F_{q,p}$ is exact and faithful.
\end{cor}

\begin{proof}
Apply Theorem~\ref{thm:main} to the ordered pair $(q,p)$.  This gives
\[
 F^{(q,p)}\cong
 H^0_{DS,+}\big|_{\KL^{k'}(\mathfrak{sl}_2)}
\]
as $\C$-linear functors.  Lemma~\ref{lem:parameter-swap-rigid-normalization} identifies the
distinguished object for the ordered pair $(q,p)$ with
$\mathcal K^{(p,q)}_{1,2}$ together with the evaluation and coevaluation
maps fixed in \cite[Theorem~7.15]{McRaeYang}.  Via the isomorphism $\phi$
of that lemma, regard the distinguished generator identification of
$F^{(q,p)}$ as an identification with the rigid datum
$(\mathcal K^{(p,q)}_{1,2},e_{\mathcal K_{1,2}},i_{\mathcal K_{1,2}})$.
The two identities in Lemma~\ref{lem:parameter-swap-rigid-normalization}
say precisely that this identification preserves the evaluation and
coevaluation data in the universal property of
\cite[Theorem~6.8]{McRaeYang}.  That theorem therefore yields a tensor
natural isomorphism
\[
 F^{(q,p)}\xRightarrow{\ \sim\ }F_{q,p}.
\]
Both functors carry the braided tensor structures specified in
\cite[Theorem~7.15]{McRaeYang}; hence this is an isomorphism of braided
tensor functors.  Therefore
\[
 F_{q,p}\cong
 H^0_{DS,+}\big|_{\KL^{k'}(\mathfrak{sl}_2)}
\]
as $\C$-linear functors.  Exactness and faithfulness follow from
Theorem~\ref{thm:main} applied to $(q,p)$.
\end{proof}

\begin{cor}[Concrete consequences]\label{cor:main-consequences}
For every $r\ge1$,
\[
 F_{p,q}(V_r)\cong K_r,\qquad
 F_{p,q}(L_r)\cong S_r.
\]
The functor $F_{p,q}$ is faithful, and its restriction to projectives is
fully faithful.
\end{cor}

\begin{proof}
The object identifications follow from
\eqref{eq:main-natural-iso} and Propositions~\ref{prop:H-Weyl} and
\ref{prop:H-simple}.  Faithfulness and projective full faithfulness are
already part of Theorem~\ref{thm:main}.
\end{proof}

\begin{cor}[Drinfeld--Sokolov reduction on projectives]
\label{cor:DS-projective-fully-faithful}
For projective objects $P,Q$, reduction induces an isomorphism
\[
 \Hom_{\KL^k}(P,Q)\xrightarrow{\ \sim\ }
 \Hom_{\Oc_{c_{p,q}}}(H(P),H(Q)).
\]
For each reflection chain, $H$ identifies the full additive subcategory
generated by its affine projectives with the full additive subcategory
generated by their Virasoro reductions.  On every finite segment it therefore identifies the endomorphism algebra
of the corresponding finite direct sum, and hence all matrix Hom spaces
between the projectives in that segment.
\end{cor}

\begin{proof}
For the natural isomorphism $\eta:F\Rightarrow H$, naturality gives the
commutative square
\[
\begin{array}{ccc}
 \Hom_{\KL^k}(P,Q) & \xrightarrow{\ F_*\ } & \Hom_{\Oc_{c_{p,q}}}(F(P),F(Q))\\[3pt]
 \big\downarrow {H_*} &&
 \big\downarrow {\,a\mapsto\eta_Q\circ a\circ\eta_P^{-1}}\\[3pt]
 \Hom_{\Oc_{c_{p,q}}}(H(P),H(Q)) & = & \Hom_{\Oc_{c_{p,q}}}(H(P),H(Q)).
\end{array}
\]
The top arrow is an isomorphism by Theorem~\ref{thm:F-fully-faithful},
and the right vertical map is the vector-space isomorphism
$\operatorname{Ad}_\eta(a)=\eta_Qa\eta_P^{-1}$.  Hence the left arrow is
an isomorphism; equivalently,
$H|_{\mathcal P^k}$ is fully faithful.  Since $H$ is fully faithful on projectives and is, by construction,
essentially surjective onto the additive subcategory generated by the
objects $H(P_{s_i})$, it induces an equivalence of these full additive
subcategories.  On each finite segment, the resulting equivalence
identifies the endomorphism algebra of the corresponding finite direct
sum, as in Corollary~\ref{cor:finite-zigzag-algebra}.
\end{proof}

\begin{rmk}[Non-canonicity of the comparison]
\label{rmk:comparison-noncanonical}
The comparison on projectives is not canonical, and its residual freedom
can be described explicitly.  At the initial vertex of a non-wall chain
every automorphism is scalar, while for $i\ge1$ every automorphism of
$H(\Pi_i)$ has the form
\[
 c_i\bigl(\id+a_iH(\nu_i)\bigr),
 \qquad c_i\in\C^\times,\quad a_i\in\C.
\]
Naturality along a nonzero adjacent arrow forces $c_{i+1}=c_i$.
Lemma~\ref{lem:radical-annihilates-arrows} imposes no condition on the
parameters $a_i$, because the radical correction terms annihilate the
adjacent arrows.  Thus the scalar is constant on each connected non-wall
reflection chain, while the radical-unipotent parameters at positive
vertices are independent.  Wall components may be scaled independently
because they are isolated.  What is rigid is the scalar obstruction
between the two orientations of an adjacent edge, not the individual
component isomorphisms.  Once a projective natural isomorphism has been
fixed, Proposition~\ref{prop:right-exact-extension} determines its
extension to the whole category uniquely.
\end{rmk}

\begin{cor}[Transport of the McRae--Yang braided tensor structure]\label{cor:transported-tensor}
Fix a natural isomorphism of underlying $\C$-linear functors
$\eta:F_{p,q}\xRightarrow{\sim}H$ as in
\eqref{eq:main-natural-iso}.  Transporting the given braided tensor
structure of $F_{p,q}$ along this fixed $\eta$ equips the underlying
$\C$-linear functor $H$ with the unique strong braided monoidal structure,
relative to this fixed $\eta$, for which $\eta$ is a braided monoidal
natural isomorphism.  For this structure $H$ preserves the source minus twist
$\theta^-_M=(-1)^\epsilon e^{2\pi iL^{\mathrm{aff}}_0}$ and the target
Virasoro twist $\theta_N=e^{2\pi iL_0}$.  The construction is made relative
to the chosen natural isomorphism $\eta$;
no independence of this choice is asserted.
\end{cor}

\begin{proof}
Write the tensorator and unit constraint of $F$ as
\[
 J^F_{X,Y}:F(X)\boxtimes F(Y)\longrightarrow F(X\boxtimes Y),
 \qquad
 J^F_0:\mathbf1\longrightarrow F(\mathbf1).
\]
This is a transport-of-structure construction.  Define
\[
 J^H_{X,Y}
 :=\eta_{X\boxtimes Y}\,J^F_{X,Y}
   (\eta_X^{-1}\boxtimes\eta_Y^{-1}),
 \qquad
 J^H_0:=\eta_{\mathbf1}\circ J^F_0.
\]
Naturality of $J^H$ follows from naturality of $J^F$ and $\eta$.
The pentagon and unit identities are obtained by conjugating the
corresponding diagrams for $F$, and the braided-functor compatibility is
obtained in the same way.  The formulas also prove uniqueness for the fixed $\eta$.  Independently of
any tensorator on $H$, Proposition~\ref{prop:H-twist} already gives the
linear identity
\[
 H(\theta^-_X)=e^{2\pi iL_0}|_{H(X)}.
\]
Theorem~7.15 of \cite{McRaeYang} says that the original braided tensor
structure on $F$ preserves the same source minus twist and the standard
Virasoro twist.  Consequently, after transporting the tensor structure
along $\eta$, the displayed linear identity is precisely the
twist-preservation axiom for $H$.
\end{proof}

\begin{rmk}\label{rmk:linear-vs-monoidal}
Theorem~\ref{thm:main} concerns the underlying $\C$-linear functors.
Corollary~\ref{cor:transported-tensor} gives a braided tensor structure on
the restricted Drinfeld--Sokolov functor by transport.  It is a
transport-of-structure consequence of the natural comparison, not an independent construction of a BRST tensorator.  In particular,
we make no claim that this transported tensorator agrees with a tensorator
obtained directly from BRST complexes, should such a construction be
available.
\end{rmk}


\begin{thebibliography}{99}

\bibitem{ArakawaVanishing}
T.~Arakawa,
\emph{Vanishing of cohomology associated to quantized Drinfeld--Sokolov reduction},
Int. Math. Res. Not. 2004 (2004), no.~15, 729--767.

\bibitem{ArakawaW}
T.~Arakawa,
\emph{Representation theory of $W$-algebras},
Invent. Math. 169 (2007), no.~2, 219--320.


\bibitem{ArakawaFrenkel}
T.~Arakawa and E.~Frenkel,
\emph{Quantum Langlands duality of representations of $\mathcal W$-algebras},
Compos. Math. 155 (2019), no.~12, 2235--2262.

\bibitem{BoeNakanoWiesner}
B.~D.~Boe, D.~K.~Nakano and E.~Wiesner,
\emph{Category $\mathcal O$ for the Virasoro algebra: cohomology and Koszulity},
Pacific J. Math. 234 (2008), no.~1, 1--21.

\bibitem{FKW}
E.~Frenkel, V.~Kac and M.~Wakimoto,
\emph{Characters and fusion rules for $W$-algebras via quantized Drinfeld--Sokolov reduction},
Comm. Math. Phys. 147 (1992), no.~2, 295--328.

\bibitem{CreutzigEtAl}
T.~Creutzig, C.~Jiang, F.~Orosz Hunziker, D.~Ridout and J.~Yang,
\emph{Tensor categories arising from the Virasoro algebra},
Adv. Math. 380 (2021), Paper No.~107601, 34 pp.

\bibitem{GoodmanWenzl}
F.~M.~Goodman and H.~Wenzl,
\emph{Ideals in the Temperley--Lieb category},
Appendix to M.~H.~Freedman, \emph{A magnetic model with a possible Chern--Simons phase},
Comm. Math. Phys. 234 (2003), no.~1, 129--183.

\bibitem{KytolaRidout}
K.~Kyt\"ol\"a and D.~Ridout,
\emph{On staggered indecomposable Virasoro modules},
J. Math. Phys. 50 (2009), 123503.

\bibitem{McRaeSopin}
R.~McRae and V.~Sopin,
\emph{Fusion and (non)-rigidity of Virasoro Kac modules in logarithmic minimal models at $(p,q)$-central charge},
Phys. Scr. 99 (2024), 035233, 42 pp.

\bibitem{McRaeYang}
R.~McRae and J.~Yang,
\emph{The non-semisimple Kazhdan--Lusztig category for affine $\mathfrak{sl}_2$ at admissible levels},
Proc. Lond. Math. Soc. (3) 130 (2025), e70043.

\bibitem{Nakano}
H.~Nakano,
\emph{Projective covers of the simple modules for the triplet $W$-algebra $\mathcal W_{p_+,p_-}$},
arXiv:2305.12448v2 [math.RT], 9 August 2026.

\bibitem{SuttonTubbenhauerWedrichZhu}
L.~Sutton, D.~Tubbenhauer, P.~Wedrich and J.~Zhu,
\emph{$\mathrm{SL}_2$ tilting modules in the mixed case},
Selecta Math. (N.S.) 29 (2023), no.~3, Paper No.~39, 40 pp.

\end{thebibliography}

\end{document}
